%!TEX TS-program = xelatex
%!TEX encoding = UTF-8 Unicode
%
% IST submission — generated by scripts/build_ist_submission.sh (Stage 5)
% Source markdown: 论文初稿P2_EN.md
% Document class: elsarticle (Elsevier; IST default)
%
\documentclass[review,12pt,authoryear]{elsarticle}

\usepackage[utf8]{inputenc}
\usepackage[T1]{fontenc}
\usepackage{lmodern}
\usepackage{amsmath,amssymb,amsthm}
\usepackage{mathtools}
\usepackage{textcomp}
\usepackage{graphicx}
\usepackage{booktabs}
\usepackage{longtable}
\usepackage{tabularx}
\usepackage{array}
\usepackage{hyperref}
\usepackage{url}
\usepackage[normalem]{ulem}
\usepackage{listings}
\usepackage{xcolor}
\usepackage{fontspec}
% Use system Times New Roman; fontspec will fall back automatically
% when the font lacks a glyph (xelatex-native unicode handling).
\setmainfont{Times New Roman}[Ligatures=TeX]
\setmonofont{Menlo}
% Math via unicode-math is unavailable in this minimal install;
% pandoc emits Greek / math symbols in body text as raw Unicode.
% xelatex handles them via the default font; unsupported glyphs are
% rendered as black boxes (acceptable for first-pass review PDF).

% Pandoc passthrough macros
\providecommand{\tightlist}{\setlength{\itemsep}{0pt}\setlength{\parskip}{0pt}}

\title{When Same-Prompt LLM Source Diversity Doesn't Help: An Ablation of Semantic Mutation Operators in Metamorphic Testing for Single-Output Scientific Computing Kernels}

\author[a]{Meng Li}
\affiliation[a]{organization={Author Affiliation 1 (placeholder)},
                country={China}}

\author[b]{(co-first author placeholder)}
\affiliation[b]{organization={Author Affiliation 2 (placeholder)},
                country={China}}

\begin{abstract}
\noindent\textbf{Context.} Metamorphic Testing (MT) addresses the test oracle problem in scientific computing software, but the fault-detection capability of metamorphic relations (MRs) has lacked a domain-aware adequacy metric: classical Mutation Score (MS) operates on syntactic AST mutations and does not capture domain semantics such as conservation laws, monotonicity, or convergence order.
\textbf{Objective.} We propose Semantic Mutation Score (SMS), built on five domain-semantic mutation operators (Conservation, Monotonicity, Convergence, Trajectory, Fidelity-order breaks) that degenerate to classical MS in the syntactic limit (modulo $D_S$-measure-zero subsets, see Section~9 for the formal statement).
\textbf{Method.} We instantiate a 12-PUT $\times$ 5-MP matrix (60 cells, average 24.3 LLM-generated mutants per cell, $N=20$ AVP repetitions) across four classes of single-output scientific computing kernels. A three-layer Layered Root-Cause Analysis (LRCA) classifier separates legitimate semantic faults from artifacts/tolerance/OOD/statistical-noise. We design a three-stage ablation --- v3 (same-source, P1-aligned), v3b (same-source, data-driven primary MP), v4 (cross-source over Claude/GPT/DeepSeek) --- to isolate contributions of MR alignment design and LLM source diversity.
\textbf{Results.} The pre-registered H2 large-effect threshold (Cliff's $\delta \geq 0.474$, Romano 2006) is not met under the pre-registered point-estimate criterion in the primary v3 analysis ($\delta = 0.323$, 95\% CI $[0.017, 0.622]$). Exploratory follow-ups (v3b: $\delta = 0.446$; v4: $\delta = 0.439$) remain below 0.474; the v3 $\to$ v3b and v3b $\to$ v4 contrasts are reported separately rather than as a single ratio. Cross-source pooling raises mutant quality (mean C1\_share $0.164 \to 0.209$) and class-c mean SMS by $+91.4\%$. Friedman $\chi^2 = 15.30$, $p = 0.0041$. SMS shows near-zero rank correlation with simple pattern coverage ($\rho = 0.16$, $n = 12$, $p = 0.61$, v4 primary); statistical power is insufficient to support an ``orthogonal'' claim.
\textbf{Conclusion.} P2 contributes a three-layer methodology for domain-semantic mutation: (Layer 1) formal necessary conditions for ``semantic mutation'', instantiated as five meta-mutation operator classes; (Layer 2) $E_1 \wedge E_2$ equivalence judgment as the conservative complete instantiation; (Layer 3) AST-normalized empirical traceability proving P2 mutants are not a subset of syntactic-mutant pools (12-PUT empirics: 5.14\% overall AST overlap, with HP/SI/TF, 159/292 mutants, categorically unreachable at 0\%). The 60-cell empirical audit demonstrates, within scope, that LLM-mutant $+$ current-MR-design produces medium- not large-effect; this is an auxiliary finding under the methodology backbone, not the paper's main contribution.
\end{abstract}

\begin{keyword}
metamorphic testing \sep mutation testing \sep semantic mutation operators \sep single-output scientific computing kernels \sep LLM-generated mutants \sep ablation study \sep metamorphic relation adequacy \sep same-prompt cross-source mutant pool \sep Cliff's delta \sep Friedman test
\end{keyword}

\begin{document}
\maketitle

\subsection{Section 1 $\cdot$ Paper Identity and
Claims}\label{section-1-paper-identity-and-claims}

\subsubsection{1.1 Terminology Unification (Throughout the
Paper)}\label{terminology-unification-throughout-the-paper}

{\def\LTcaptype{none} % do not increment counter
\begin{longtable}[]{@{}
  >{\raggedright\arraybackslash}p{(\linewidth - 4\tabcolsep) * \real{0.3333}}
  >{\raggedright\arraybackslash}p{(\linewidth - 4\tabcolsep) * \real{0.3333}}
  >{\raggedright\arraybackslash}p{(\linewidth - 4\tabcolsep) * \real{0.3333}}@{}}
\toprule\noalign{}
\begin{minipage}[b]{\linewidth}\raggedright
Use
\end{minipage} & \begin{minipage}[b]{\linewidth}\raggedright
Do Not Use
\end{minipage} & \begin{minipage}[b]{\linewidth}\raggedright
Rationale
\end{minipage} \\
\midrule\noalign{}
\endhead
\bottomrule\noalign{}
\endlastfoot
\textbf{4 representative classes of scientific computing programs} &
\st{4 paradigms} & ``Paradigm'' carries strong semantics in the SE/PL
community as ``programming paradigm = OO/FP/Logic,'' which differs from
this paper's meaning of ``scientific computing program categories'' \\
\textbf{cross-class consistency} & \st{cross-paradigm consistency} &
Aligns with terminology unification \\
\end{longtable}
}

The English text uniformly uses \emph{four representative classes of
single-output scientific computing kernels} and \emph{cross-class
consistency}, avoiding \emph{paradigm}. \textbf{The scope of the
contribution claim is strictly bounded to single-output
\texttt{float\ $\to$\ float} kernels (each PUT source code \textless{} 2 KB,
total 12 PUTs across 4 classes: numeric / probabilistic / surrogate /
ML)}; the title, §1.6, §6.5, and §7.5 use this scope consistently. Where
``scientific computing programs'' appears in the body text, it functions
only as a general domain term (as in §1.3 LLM-mutant literature review,
§1.6 P-series roadmap) and \textbf{does not constitute a methodological
claim over industrial-scale, multi-module, or multi-output scientific
computing software}; §6.5 stakeholder pain-points and §6.5.3 V\&V
documentation discussion are also constrained to this single-output
kernels scope.

\subsubsection{1.2 Core Claims}\label{core-claims}

This paper proposes a \textbf{three-layer methodological framework}
around domain-semantic mutation operators (explicitly revised in P2 R2):

\begin{itemize}
\tightlist
\item
  \textbf{Layer 1 --- Definitional} (§3.2.0): Establishes necessary
  conditions for semantic mutation: (a) cross-function-boundary
  replacement, (b) carrying domain knowledge, (c) altering algorithmic
  class. The five meta mutation operator classes (CE/OS/HP/TF/SI =
  mut\_C/M/G/T/F) are specializations of these necessary conditions
  across the four PUT classes.
\item
  \textbf{Layer 2 --- Operational} (§2.3 / §4.4): Provides equivalence
  criteria E1 $\wedge$ E2 as a conservative complete instantiation of the
  necessary conditions. Justifies the selection among three candidate
  definitions (semantic equivalence / output equivalence / both).
\item
  \textbf{Layer 3 --- Applied} (§3.2.6.3): Uses equivalence detection
  tools to trace P2 mutants back to syntactic mutants, providing
  positive empirical evidence that the P2 mutant pool $\not\subset$ syntactic mutant
  pool (full 12-PUT cosmic-ray empirics: \textbf{5.14\%} overall AST
  overlap rate, with three classes --- HP/SI/TF, totaling 159 mutants
  --- at \textbf{0/0/0} overlap).
\end{itemize}

Additionally, this paper introduces three supporting tools:
\textbf{Semantic Mutation Score (SMS)}, following the classic
\texttt{killed/(mut−equiv)} structure, plus an \textbf{engineering
attribution layer} using likely root cause analysis (LRCA) to conduct an
empirical audit demonstration of MR sets across 60 cells spanning four
representative classes of scientific computing programs. The audit
reports: (a) tool implementation feasibility, (b) SMS behavior on meta
pattern slices, (c) cross-class consistency, and (d) empirical
differences from existing MR metrics (§5). \textbf{The 60-cell audit is
an empirical demonstration following the establishment of the
three-layer methodological framework, not the paper's main contribution}
--- the primary contribution lies in the methodological framework of
Layers 1-3.

\textbf{Theoretical commitment}: This paper's SMS strictly reduces to
the classic Jia \& Harman syntactic Mutation Score (MS) in the
degenerate limit --- classic MS is a special case of SMS when all
extension dimensions are disabled. This establishes this paper's
identity as a true generalization of classic mutation testing.

This paper does not claim ``how adequacy should be defined'' --- that
proposition is reserved for P4 (TOSEM) unified theoretical work.

\subsubsection{1.3 Three-Year Roadmap Position and Related
Work}\label{three-year-roadmap-position-and-related-work}

\paragraph{1.3.1 Three-Year Roadmap}\label{three-year-roadmap}

P2 = $\gamma$ dual-track parallel development (P2 leverages P1's 12-PUT
infrastructure). Concurrent papers:

\begin{itemize}
\tightlist
\item
  \textbf{P1} (SANER 2027): MR meta pattern empirical audit
\item
  \textbf{P2} (this paper, IST 2027 Q3)
\item
  \textbf{P2-CN} (Nuclear Power Engineering 2027 Q4): Nuclear
  engineering practice details, cites P2
\item
  \textbf{P4} (TOSEM): Unified theory + three-pillar coupling theorems,
  cites P2 empirical foundation
\item
  \textbf{P5} (\emph{Nuclear Engineering and Design}, NED, 2028):
  Regulatory extension, cites P2-CN
\end{itemize}

\paragraph{1.3.2 Related Work}\label{related-work}

Overview of semantic-aware mutation work in the past 5 years:

{\def\LTcaptype{none} % do not increment counter
\begin{longtable}[]{@{}
  >{\raggedright\arraybackslash}p{(\linewidth - 4\tabcolsep) * \real{0.3333}}
  >{\raggedright\arraybackslash}p{(\linewidth - 4\tabcolsep) * \real{0.3333}}
  >{\raggedright\arraybackslash}p{(\linewidth - 4\tabcolsep) * \real{0.3333}}@{}}
\toprule\noalign{}
\begin{minipage}[b]{\linewidth}\raggedright
Work
\end{minipage} & \begin{minipage}[b]{\linewidth}\raggedright
Domain
\end{minipage} & \begin{minipage}[b]{\linewidth}\raggedright
Difference from P2
\end{minipage} \\
\midrule\noalign{}
\endhead
\bottomrule\noalign{}
\endlastfoot
Humbatova, Jahangirova \& Tonella (DeepCrime, ISSTA 2021) & Deep
learning systems (real-fault-based mutation) & Single class (deep
learning); P2 includes ML as one of 4 classes \\
Tip, Bell \& Schäfer (LLMorpheus, arXiv 2024) & JavaScript LLM mutants &
Single language; P2 conducts cross-source LLM ablation on Python
scientific computing kernels (§4.2.5) \\
Jia \& Harman (classic survey, TSE 2011) + Papadakis et al.~(survey,
Adv. Computers 2019) & Syntactic mutation & P2 is backward-compatible
(§9 degeneration theorem) and extends at the domain-semantic layer \\
\end{longtable}
}

P2's unique positioning: \textbf{meta pattern-driven + unified framework
across 4 classes + strict compatibility with classic MS as an
intensional extension}.

\textbf{CPH grounding of classical mutation testing (R2 round-2 NEW)}:
The syntactic baseline of this methodology rests on the classical
Coupling Effect Hypothesis (CPH) --- DeMillo, Lipton \& Sayward (1978)
hypothesized that tests detecting simple faults also detect complex
faults; Andrews, Briand \& Labiche (2005) and Just et al.~(2014, FSE)
empirically confirmed that mutants are valid surrogates for real faults;
Papadakis et al.~(2019) provided the most comprehensive post-Jia \&
Harman (2011) survey of mutation testing advances. §3.2.6 of this paper
argues: \textbf{CPH holds within the scope of syntactic mutation but is
non-trivial in the scope of domain-semantic mutation} --- syntactic
mutants cannot reach P2's four semantic operator classes (HP/SI/TF/OS);
even when syntactic mutation testing couples simple syntactic faults to
complex syntactic faults, that coupling does not extend to
domain-semantic faults (see §3.2.6.0 systematic-vs-incidental + §3.2.6.3
12-PUT empirics 5.14\% AST overlap rate).

\textbf{Recent work on LLM-generated mutants}: Tip et al.~(2024)
proposed LLMorpheus, using LLMs to generate mutants on JavaScript
instead of fixed operator sets, reporting fault-detection comparable to
traditional operators but with lower equivalent rates, and observing
medium-effect intervals in effect-size reports for LLM-mutants. Petrović
\& Ivanković (2018) reported approximately 20\% productive mutant ratio
on Google's internal 500,000-mutant dataset, numerically close to this
paper's §5.6.2 LRCA C1\_share measured levels (default threshold 0.16,
calibrated 0.20) --- \textbf{a numerical coincidence rather than
mechanism validation}; see §6.1 for the detailed disambiguation. This
paper's §5.7.2 measured Cliff's $\delta$ = 0.323 is in the same magnitude as
the LLM-mutant medium-effect phenomenon observed by Tip et al.~(2024) on
JavaScript. \textbf{Estimand caveat}: Tip 2024 compares ``LLM mutants vs
traditional mutants on fault detection rate'' (cross-source comparison),
while this paper's §5.7.2 compares ``aligned vs cross MP slice on the
same mutant pool'' (single-source within-pool comparison). The numerical
proximity of the two $\delta$ values does not constitute substantive support,
serving only as a reference to the medium-effect phenomenon in
LLM-mutant literature.

\textbf{Relation to existing V\&V standards (R3 round-2 NEW)}: This
paper's SMS is conceptually complementary to ASME V\&V 20-2009
\emph{Standard for Verification and Validation in Computational Fluid
Dynamics and Heat Transfer} §3 code verification --- the latter targets
code-level correctness verification of numerical solvers, while this
paper's SMS targets fault-detection adequacy of MR sets. This paper's
empirical scope is strictly bounded to single-output kernels (§3.1.1),
with a substantial scale gap to the multi-module CFD codes targeted by
V\&V 20; see §6.5.3 long-term aspiration discussion.

\subsubsection{1.4 Research Questions (Purely Empirical, 4
Items)}\label{research-questions-purely-empirical-4-items}

\begin{itemize}
\tightlist
\item
  \textbf{RQ1} What are the distributions of instantiation rate
  (inst\_rate), equivalent rate (equiv\_rate), C1 share (C1\_share), and
  survival rate (survive\_rate) for the 5 operators across 60 cells (12
  PUTs $\times$ 5 operators)?
\item
  \textbf{RQ2} What is the difference structure of SMS between
  operator-meta pattern aligned (j=k) and cross (j$\neq$k) slices? Is the
  equiv\_rate of P1 $\circ$ vacant cells systematically higher than that of \textbullet\textbullet
  populated cells?
\item
  \textbf{RQ3} What is the cross-class consistency of SMS across 4
  program classes $\times$ 5 operators (sign test + coefficient of variation)?
\item
  \textbf{RQ4} What are the empirical difference distributions between
  SMS and pattern coverage in ranking and discriminative power?
  (Descriptive, no presupposed answer)
\end{itemize}

\subsubsection{1.5 Hypothesis System (Originally 5 Items; H3 Formally
Retired, 4 Valid
Hypotheses)}\label{hypothesis-system-originally-5-items-h3-formally-retired-4-valid-hypotheses}

\begin{itemize}
\tightlist
\item
  \textbf{H1 Operator Implementation Feasibility}: $\geq$ 4 of the 5
  operators can produce $\geq$ 5 non-equivalent mutants on $\geq$ 9/12 PUTs
\item
  \textbf{H2 Meta Pattern Aligned Slice}: aligned-SMS / cross-SMS odds
  ratio $\geq$ 3.0, Cliff's $\delta$ $\geq$ 0.474
\item
  \st{\textbf{H3 Equiv Aligned Vacancy (P1 H6 Dual)}: 6 $\circ$ cells have
  equiv\_rate $\geq$ 0.85; populated cells have equiv\_rate $\leq$ 0.30}
  \textbf{(Formally Retired)}
\item
  \textbf{H4 Cross-Class Consistency}: $\Delta$SMS shows all-positive sign test
  across 4 program classes, coefficient of variation (CV) \textless{}
  0.5
\item
  \textbf{H5 LRCA Robustness}: Average suspect share (suspect\_share)
  across all 60 cells $\leq$ 0.20
\end{itemize}

\begin{quote}
\textbf{Formal Statement on Retired Hypothesis H3}: The P1-derived
version of this paper included H3 regarding the bidirectional threshold
hypothesis ``$\circ$ vacant cells equiv\_rate $\geq$ 0.85; \textbullet\textbullet populated cells
equiv\_rate $\leq$ 0.30.'' In the measured data space of LLM-generated
mutants, equivalence detection triggered in extremely few cells
(\textless{} 10 cells with non-zero equiv), causing structural collapse
of the bidirectional threshold comparison space, rendering this
hypothesis unable to make meaningful determinations on this paper's
data. This paper formally retired H3 after v3 data collection and treats
the R\_sem / R\_kill decoupling phenomenon as a descriptive observation
in §6.2 (supported jointly by §4.8.3 operator-level pilot + §5.6.2 LRCA
C1\_share), not as a substitute for formal hypothesis testing.
\textbf{To maintain citation consistency with §3 / §5 / §6 / §7,
subsequent sections continue using H1, H2, H4, H5 numbering (with H3
vacant), without renumbering throughout the paper.}
\end{quote}

\subsubsection{1.6 P2 / P4 / {[}Meng Li et al., Progress in Nuclear
Energy, under review{]}
Boundaries}\label{p2-p4-meng-li-et-al.-progress-in-nuclear-energy-under-review-boundaries}

{\def\LTcaptype{none} % do not increment counter
\begin{longtable}[]{@{}
  >{\raggedright\arraybackslash}p{(\linewidth - 4\tabcolsep) * \real{0.3333}}
  >{\raggedright\arraybackslash}p{(\linewidth - 4\tabcolsep) * \real{0.3333}}
  >{\raggedright\arraybackslash}p{(\linewidth - 4\tabcolsep) * \real{0.3333}}@{}}
\toprule\noalign{}
\begin{minipage}[b]{\linewidth}\raggedright
Dimension
\end{minipage} & \begin{minipage}[b]{\linewidth}\raggedright
P2 Main Text (This Paper)
\end{minipage} & \begin{minipage}[b]{\linewidth}\raggedright
P4 (TOSEM)
\end{minipage} \\
\midrule\noalign{}
\endhead
\bottomrule\noalign{}
\endlastfoot
Contribution Type & Tool + empirical report & Theoretical formalization
+ theorems \\
Proposition Answered & ``How does the tool behave'' (empirical) & ``How
should adequacy be defined'' (theoretical) \\
Empirical Foundation & 12 PUTs, 60 cells & No experiments, cites P1 +
P2 \\
P4-Exclusive Dimensions & Not involved & Minimal MR subset existence
theorem / reachable adequacy theorem / three-pillar coupling theorem \\
\end{longtable}
}

\paragraph{1.6.1 P2 Innovation
Attribution}\label{p2-innovation-attribution}

{\def\LTcaptype{none} % do not increment counter
\begin{longtable}[]{@{}
  >{\raggedright\arraybackslash}p{(\linewidth - 2\tabcolsep) * \real{0.5000}}
  >{\raggedright\arraybackslash}p{(\linewidth - 2\tabcolsep) * \real{0.5000}}@{}}
\toprule\noalign{}
\begin{minipage}[b]{\linewidth}\raggedright
Innovation
\end{minipage} & \begin{minipage}[b]{\linewidth}\raggedright
Content
\end{minipage} \\
\midrule\noalign{}
\endhead
\bottomrule\noalign{}
\endlastfoot
\textbf{C-I Domain-Semantic Mutation Operators} & mut\_j $\in$ MUT five
classes of semantic failure injection, dual to 5 meta patterns \\
\textbf{C-II SMS + Backward Compatibility} & Classic
\texttt{killed/(mut−equiv)} structure + three-dimensional intensional
extension + degenerates to classic MS \\
\textbf{C-III LRCA Engineering Attribution Layer} & 5 likely root causes
+ 3-layer diagnosis, purely descriptive, not in SMS formula \\
\textbf{C-IV 60-Cell Empirical Audit} & All 4 PUT classes $\times$ all 5
operators \\
\end{longtable}
}

\paragraph{1.6.2 Epistemological
Statement}\label{epistemological-statement}

SMS is an epistemological semantic detection metric, not an engineering
value proxy. Engineering value requires separate metrics (P2-CN subject
matter). \textbf{Toy-scope caveat (R2 round-2 NEW)}: The 12 PUTs in this
paper are Numerical-Recipes-style ``toy kernels'' (each PUT \textless{}
2 KB, signature limited to \texttt{float\ $\to$\ float}), intended to
provide a verifiable minimum-working-example for the methodological
backbone (§3.2.0 necessary conditions, §2.3-§4.4 equivalence judgment,
§3.2.6.3 traceability). Empirical results at this scale do \textbf{not}
constitute a methodological claim over industrial-scale, multi-module,
or multi-output scientific computing software; transfer of the empirics
to that scale is left to P5 (domain application).

\begin{center}\rule{0.5\linewidth}{0.5pt}\end{center}

\subsection{Section 2 $\cdot$ Semantic Mutation Notation
System}\label{section-2-semantic-mutation-notation-system}

\subsubsection{2.0 Dual Fundamental
Principles}\label{dual-fundamental-principles}

{\def\LTcaptype{none} % do not increment counter
\begin{longtable}[]{@{}
  >{\raggedright\arraybackslash}p{(\linewidth - 4\tabcolsep) * \real{0.3333}}
  >{\raggedright\arraybackslash}p{(\linewidth - 4\tabcolsep) * \real{0.3333}}
  >{\raggedright\arraybackslash}p{(\linewidth - 4\tabcolsep) * \real{0.3333}}@{}}
\toprule\noalign{}
\begin{minipage}[b]{\linewidth}\raggedright
Principle
\end{minipage} & \begin{minipage}[b]{\linewidth}\raggedright
Name
\end{minipage} & \begin{minipage}[b]{\linewidth}\raggedright
Operational Meaning
\end{minipage} \\
\midrule\noalign{}
\endhead
\bottomrule\noalign{}
\endlastfoot
\textbf{P-I} & Developmental & Develop upon the classical mutation
testing framework to accommodate semantic mutation requirements for
scientific computing programs and align with MR validity determination
needs \\
\textbf{P-II} & Stable/Open & Notation system stable (skeleton fixed),
content open (operators/meta patterns/root causes/classes extensible by
future researchers) \\
\end{longtable}
}

\subsubsection{2.1 Notation System Skeleton (Fixed
Layer)}\label{notation-system-skeleton-fixed-layer}

\paragraph{2.1.1 Classical Mutation Testing Vocabulary Inheritance
(Seven Core
Concepts)}\label{classical-mutation-testing-vocabulary-inheritance-seven-core-concepts}

{\def\LTcaptype{none} % do not increment counter
\begin{longtable}[]{@{}
  >{\raggedright\arraybackslash}p{(\linewidth - 4\tabcolsep) * \real{0.3333}}
  >{\raggedright\arraybackslash}p{(\linewidth - 4\tabcolsep) * \real{0.3333}}
  >{\raggedright\arraybackslash}p{(\linewidth - 4\tabcolsep) * \real{0.3333}}@{}}
\toprule\noalign{}
\begin{minipage}[b]{\linewidth}\raggedright
Classical Concept
\end{minipage} & \begin{minipage}[b]{\linewidth}\raggedright
P2 Correspondence
\end{minipage} & \begin{minipage}[b]{\linewidth}\raggedright
Extension Point
\end{minipage} \\
\midrule\noalign{}
\endhead
\bottomrule\noalign{}
\endlastfoot
Program Under Test (PUT) & \texttt{S\_i} & Only introduces index i $\in$ I,
pure naming inheritance \\
Mutation Operator & \texttt{mut\_j\ $\in$\ MUT} & Syntactic/grammatical
rules $\to$ domain-semantic failure injection \\
Mutant & \texttt{s\textquotesingle{}\ $\in$\ mut\_j(S\_i)} & Syntactic
transformation product $\to$ semantic transformation product \\
Equivalent Mutant & \texttt{equiv\_\{i,k,j\}} & Bitwise behavioral
equivalence $\to$ semantic-class equivalence (E1 $\wedge$ E2) \\
Killed Mutant & \texttt{killed\_\{i,k,j\}} & Equality oracle detects
difference $\to$ meta pattern AVP detects \\
Surviving Mutant & \texttt{survive\_\{i,k,j\}} & No structural
extension \\
Mutation Score (MS) & \texttt{SMS\ =\ killed/(mut−equiv)} & Formula
structure retained, extended only through internal semantics of
mut/equiv/killed \\
\end{longtable}
}

\paragraph{2.1.2 Complete Notation System
Table}\label{complete-notation-system-table}

\begin{verbatim}
═══════════════════ Paper Domain ═══════════════════
Paper Identity:        P1, P2, P4, P5

═══════════════════ Experimental Objects ═══════════════════
PUT Set I:  {A1,A2,A3,B1,B2,B3,C1,C2,C3,D1,D2,D3}, |I|=12
PUT:        S_i  (i $\in$ I)
Class Mapping:          cls : I $\to$ {A, B, C, D}    (open, extensible)
Valid Input Distribution:    D_S, sampling X_{K_eq} ~ D_S

═══════════════════ Metamorphic Testing Side ═══════════════════
Meta Pattern (MP):
  MP = {MP_1,...,MP_5}, k $\in$ {1,...,5}    (open, extensible)
  MP_1 Conservation  MP_2 Monotonicity  MP_3 Convergence  MP_4 Trajectory  MP_5 Partial-order
Metamorphic Relation:        MR_{i,k} (provided by P1), mr = (r,R) $\in$ MR_{i,k}
Automated Verification Pipeline (AVP):
  AVP : Programs $\times$ MR_universe $\times$ R⁺ $\to$ {pass, fail}
  AVP(s, mr, $\varepsilon$_AVP^k)  invokes corresponding verification method per MP_k

═══════════════════ Mutation Testing Side ═══════════════════
Mutation Operator Family (MUT):
  MUT = {mut_1,...,mut_5}, j $\in$ {1,...,5}    (open, extensible)
  mut_1=mut_C  mut_2=mut_M  mut_3=mut_G  mut_4=mut_T  mut_5=mut_F
  Signature: mut_j : Programs $\to$ 2^Programs
Mutant:          s' $\in$ mut_j(S_i)
Alignment:            align(j) = j

═══════════════════ Three-State Decomposition (Fixed) ═══════════════════
mut_j(S_i) = equiv_{i,k,j} $\sqcup$ killed_{i,k,j} $\sqcup$ survive_{i,k,j}    (mutually exclusive and exhaustive)

equiv Determination Dual Conditions:
  (E1) AVP-coherent:    $\forall$ mr $\in$ MR_{i,k}: AVP(S_i, mr) = AVP(s', mr)
  (E2) Output-equiv:    $\forall$ x $\in$ X_{K_eq} ~ D_S: $\|$S_i(x) − s'(x)$\|$ $\leq$ $\varepsilon$_eq

killed Determination:
  killed(s', MR_{i,k}) ⇔ $\exists$ mr $\in$ MR_{i,k}:
                          AVP(S_i, mr) = pass $\wedge$ AVP(s', mr) = fail

═══════════════════ Metric (Fixed Classical Structure) ═══════════════════
SMS_{i,k,j} := |killed_{i,k,j}| / (|mut_j(S_i)| − |equiv_{i,k,j}|)
            = |killed_{i,k,j}| / (|killed_{i,k,j}| + |survive_{i,k,j}|)
            $\in$ [0, 1]

═══════════════════ LRCA Engineering Attribution Layer (Descriptive, Not in SMS) ═══════════════════
Likely Root Cause Inventory:    C = {C1,...,C5}    (open, extensible)
  C1 True Semantic Failure / C2 Tolerance Perturbation / C3 Out-of-Distribution (OOD)
  C4 Statistical Assumption Violation / C5 Mutator Artifact
LRCA Output:       Each s' $\in$ killed annotated with root_cause(s') $\in$ C
Descriptive Quantities:          C1_share_{i,k,j}, suspect_share_{i,k,j}    (reported in §5)

═══════════════════ Slicing and Cross-Class ═══════════════════
Slicing:            aligned: j=k    cross: j$\neq$k
Cross-class:            $\Delta$SMS_c, c $\in$ {A,B,C,D}
                CV($\Delta$SMS) := std($\Delta$SMS_c) / |mean($\Delta$SMS_c)|

═══════════════════ Tolerance and Sampling ═══════════════════
Tolerance:            $\varepsilon$_AVP^k (meta pattern dependent), $\varepsilon$_eq (equivalence output tolerance)
Sampling:            K_eq = 1000 (E2 sample size), N = 20 (statistical replicates)
\end{verbatim}

\paragraph{2.1.3 Notation Stability
Declaration}\label{notation-stability-declaration}

The notation skeleton of this paper (11 core concepts + three-state
decomposition + SMS formula + AVP interface) is \textbf{fixed}. The
\textbf{content} of operator family MUT, meta patterns MP, likely root
cause inventory C, and PUT class mapping cls is open; future researchers
may extend with new elements while preserving notation semantics.

\subsubsection{2.2 Domain-Semantic Mutation Operators mut\_j $\in$
MUT}\label{domain-semantic-mutation-operators-mut_j-mut}

\paragraph{2.2.1 Operator Signature
(Fixed)}\label{operator-signature-fixed}

\begin{verbatim}
mut_j : Programs $\to$ 2^Programs
\end{verbatim}

\paragraph{2.2.2 Five Operator Instantiations in This Paper (Open,
Extensible to mut\_6,
\ldots)}\label{five-operator-instantiations-in-this-paper-open-extensible-to-mut_6}

{\def\LTcaptype{none} % do not increment counter
\begin{longtable}[]{@{}lll@{}}
\toprule\noalign{}
Operator & Failure Semantics & Alignment with MP \\
\midrule\noalign{}
\endhead
\bottomrule\noalign{}
\endlastfoot
mut\_1 = mut\_C & Conservation-breaking & MP\_1 Conservation \\
mut\_2 = mut\_M & Monotonicity-breaking & MP\_2 Monotonicity \\
mut\_3 = mut\_G & Convergence-breaking & MP\_3 Convergence \\
mut\_4 = mut\_T & Trajectory-distorting & MP\_4 Trajectory \\
mut\_5 = mut\_F & Fidelity-order-breaking & MP\_5 Partial-order \\
\end{longtable}
}

\paragraph{2.2.3 Operator Design
Principles}\label{operator-design-principles}

\begin{itemize}
\tightlist
\item
  Inject domain-semantic invariant violations
\item
  Syntactically correct and executable
\item
  Generated by humans or machines with semantic understanding capability
  (Large Language Model, LLM)
\item
  No claim of strict one-to-one operator-meta pattern correspondence
  (\texttt{align(j)=j} is a design choice, not a theorem)
\end{itemize}

\subsubsection{2.3 (Semantic) Equivalent Mutant
equiv}\label{semantic-equivalent-mutant-equiv}

\textbf{Layer 2 --- Operational (Equivalence Determination =
Instantiation of §3.2.0 Necessary Conditions)}

The equivalence determination E1 $\wedge$ E2 given in this section is an
\textbf{executable instantiation} of the necessary conditions (a)(b)(c)
in §3.2.0. Logical mapping:

\begin{itemize}
\tightlist
\item
  \textbf{E1} (AVP-coherent, ``consistent with original program
  determination via AVP for all MR\_i,k when passing'') $\leftrightarrow$ converse of
  necessary condition (c): ``whether algorithm class is consistent'' ---
  if mutant behavior is indistinguishable from original program within
  the MR framework, then algorithm class is consistent (c not
  satisfied).
\item
  \textbf{E2} (Output-equivalent on K\_eq=1000 sampling within $\varepsilon$\_eq) $\leftrightarrow$
  converse of necessary conditions (a)(b): ``whether consistent at pure
  output layer without relying on cross-function/domain knowledge
  distinction'' --- if two mutants are consistent on K\_eq numerical
  sampling, they have identical pure numerical behavior outside (a)(b).
\item
  \textbf{E1 $\wedge$ E2} = \textbf{conservative complete instantiation} of
  necessary conditions = converse of (c) $\wedge$ converse of (a)(b) = ``judged
  equivalent in all converse directions of necessary conditions.''
\end{itemize}

\textbf{Trade-offs Among Three Candidates}:

{\def\LTcaptype{none} % do not increment counter
\begin{longtable}[]{@{}
  >{\raggedright\arraybackslash}p{(\linewidth - 6\tabcolsep) * \real{0.2500}}
  >{\raggedright\arraybackslash}p{(\linewidth - 6\tabcolsep) * \real{0.2500}}
  >{\raggedright\arraybackslash}p{(\linewidth - 6\tabcolsep) * \real{0.2500}}
  >{\raggedright\arraybackslash}p{(\linewidth - 6\tabcolsep) * \real{0.2500}}@{}}
\toprule\noalign{}
\begin{minipage}[b]{\linewidth}\raggedright
Determination
\end{minipage} & \begin{minipage}[b]{\linewidth}\raggedright
False Positive (Misjudge equiv)
\end{minipage} & \begin{minipage}[b]{\linewidth}\raggedright
False Negative (Miss equiv)
\end{minipage} & \begin{minipage}[b]{\linewidth}\raggedright
Bias on SMS
\end{minipage} \\
\midrule\noalign{}
\endhead
\bottomrule\noalign{}
\endlastfoot
\textbf{E1 alone} (Semantic same) & Mutant numerically coincides on
K\_eq inputs but consistent within MR framework & E1 false $\to$ truly
non-equivalent & SMS biased low (easier to judge equiv $\to$ denominator
mut-equiv smaller) \\
\textbf{E2 alone} (Output same) & Mutant output consistent but MR
behavior different (rare, nearly impossible) & E2 false on K\_eq
sampling but consistent over full space & SMS biased high (harder to
judge equiv) \\
\textbf{E1 $\wedge$ E2} & Simultaneous false positive (harder) & E1 $\vee$ E2 false
implies non-equiv (easier to judge non-equiv) & \textbf{SMS biased high
(conservative, fewer equiv)} \\
\end{longtable}
}

\textbf{Why Choose E1 $\wedge$ E2}: On LLM-generated mutants, (i) E2 alone is
easily deceived by numerical coincidence (mutant accidentally outputs
approximations at sampling points but semantically different); (ii) E1
alone is easily deceived by insufficient AVP coverage (if
\textbar MR\_i,k\textbar{} is small, E1 pass rate is high but mutant
actually non-equivalent); (iii) E1 $\wedge$ E2 = ``AVP + numerical sanity
check'' dual-layer verification, more robust.

\textbf{Counterexample Contrast}: - E2 passes but E1 does not: extremely
rare --- mutant coincides at K\_eq sampling points but significantly
deviates at MR\_i,k trigger points (mutant should not be judged equiv,
E1 $\wedge$ E2 correctly judges non-equiv) - E1 passes but E2 does not: common
--- mutant behavior consistent within MR framework but has numerical
drift beyond $\varepsilon$\_eq (still should not judge equiv, E1 $\wedge$ E2 correct)

\textbf{Degenerate Limit Connection to §9}: Under degenerate limit
L\_equiv (L1 $\wedge$ L2), E2 alone already degenerates to classical bitwise
equivalence (Lemma 9.1); E1 degenerates to trivial condition (MP set
degenerates to equality determination under L4); three candidates are
\textbf{almost everywhere consistent} under L. Under current paper data
(non-degenerate limit), the conservative choice of E1 $\wedge$ E2 is
engineering-sound.

\paragraph{2.3.1 Determination Dual
Conditions}\label{determination-dual-conditions}

\begin{verbatim}
equiv(s', S_i, MR_{i,k}) ⇔
  (E1) AVP-coherent
  (E2) Output-equiv (K_eq=1000, $\varepsilon$_eq inherits P1 cell-specific tolerance)
\end{verbatim}

\paragraph{2.3.2 Degenerate Relation to Classical Equivalent
Mutant}\label{degenerate-relation-to-classical-equivalent-mutant}

In the limit $\varepsilon$\_eq $\to$ 0, K\_eq $\to$ full input space, (E2) degenerates to
bitwise equivalence; (E1) is implied by (E2) in that limit. equiv
reduces to classical behavioral equivalence.

\subsubsection{2.4 Semantic Mutation Score
SMS}\label{semantic-mutation-score-sms}

\paragraph{2.4.1 Formula}\label{formula}

\begin{verbatim}
SMS_{i,k,j} := |killed_{i,k,j}| / (|mut_j(S_i)| − |equiv_{i,k,j}|)
            = |killed_{i,k,j}| / (|killed_{i,k,j}| + |survive_{i,k,j}|)
            $\in$ [0, 1]
\end{verbatim}

Literally isomorphic to the Jia \& Harman MS formula.

\paragraph{2.4.2 Three Dimensions of Internal
Extension}\label{three-dimensions-of-internal-extension}

{\def\LTcaptype{none} % do not increment counter
\begin{longtable}[]{@{}
  >{\raggedright\arraybackslash}p{(\linewidth - 4\tabcolsep) * \real{0.3333}}
  >{\raggedright\arraybackslash}p{(\linewidth - 4\tabcolsep) * \real{0.3333}}
  >{\raggedright\arraybackslash}p{(\linewidth - 4\tabcolsep) * \real{0.3333}}@{}}
\toprule\noalign{}
\begin{minipage}[b]{\linewidth}\raggedright
Dimension
\end{minipage} & \begin{minipage}[b]{\linewidth}\raggedright
Classical
\end{minipage} & \begin{minipage}[b]{\linewidth}\raggedright
P2
\end{minipage} \\
\midrule\noalign{}
\endhead
\bottomrule\noalign{}
\endlastfoot
Mutation Generation & Rule-based syntactic operators & Domain-semantic
operators (LLM/human) \\
Equivalence Determination & Bitwise behavioral equivalence & Tolerance
equivalence (E1 $\wedge$ E2) \\
Test Oracle & Bitwise equality oracle & Meta pattern AVP \\
\end{longtable}
}

\paragraph{2.4.3 Backward Compatibility
Declaration}\label{backward-compatibility-declaration}

\begin{quote}
The semantic mutation notation system of this paper reuses the classical
mutation testing vocabulary of Jia \& Harman (TSE 2011) --- the seven
core concepts of program under test, mutation operator, mutant,
equivalent mutant, killed mutant, surviving mutant, and mutation score
are aligned item by item, and the SMS formula strictly preserves the
classical structure \textbf{SMS = \textbar killed\textbar{} /
(\textbar mut\textbar{} − \textbar equiv\textbar)}. The extensions of
this paper are applied only to the \textbf{internal definitions} of
three concepts: mut (syntactic operators $\to$ domain-semantic operators),
equiv (bitwise behavioral equivalence $\to$ semantic-class equivalence
E1$\wedge$E2), and killed (equality oracle $\to$ meta pattern AVP), without
introducing any new formula terms or new state classifications. Under
the degenerate limit where all extension dimensions are closed --- i.e.,
$\varepsilon$\_eq $\to$ 0, K\_eq $\to$ full input space, $\varepsilon$\_AVP\^{}k $\to$ 0, MP set reduces to
equality determination, mut\_j switches to rule-based syntactic
operators, PUT class restricted to imperative deterministic programs ---
SMS strictly regresses to the classical syntactic mutation score,
equivalent mutants degenerate to classical behavioral equivalence,
killed mutants degenerate to bitwise difference detection, and the LRCA
engineering attribution layer trivializes because C2-C5 cannot be
triggered. This backward compatibility is the intrinsic scientific
guarantee of the notation system in this paper: any empirical conclusion
based on SMS is structurally consistent with existing mutation testing
literature under classical syntactic mutation scenarios, and does not
constitute semantic fragmentation at the metric level.
\end{quote}

\subsubsection{2.5 AVP Reuse Protocol}\label{avp-reuse-protocol}

\paragraph{2.5.1 Interface Signature (Provided by
P1)}\label{interface-signature-provided-by-p1}

\begin{verbatim}
AVP : Programs $\times$ MR_universe $\times$ R⁺ $\to$ {pass, fail}
AVP(s, mr, $\varepsilon$_AVP^k)  invokes corresponding verification method per MP_k
\end{verbatim}

\paragraph{2.5.2 Meta Pattern Verification Methods (Currently 5 MP
Classes,
Extensible)}\label{meta-pattern-verification-methods-currently-5-mp-classes-extensible}

{\def\LTcaptype{none} % do not increment counter
\begin{longtable}[]{@{}
  >{\raggedright\arraybackslash}p{(\linewidth - 2\tabcolsep) * \real{0.5000}}
  >{\raggedright\arraybackslash}p{(\linewidth - 2\tabcolsep) * \real{0.5000}}@{}}
\toprule\noalign{}
\begin{minipage}[b]{\linewidth}\raggedright
Meta Pattern
\end{minipage} & \begin{minipage}[b]{\linewidth}\raggedright
Verification Method
\end{minipage} \\
\midrule\noalign{}
\endhead
\bottomrule\noalign{}
\endlastfoot
MP\_1 Conservation & Tolerance equality \textbar LHS−RHS\textbar{} $\leq$
$\varepsilon$ \\
MP\_2 / MP\_5 & Wilcoxon signed-rank test, $\alpha$=0.05 \\
MP\_3 Convergence & Convergence order estimation + asymptotic residual
ratio \\
MP\_4 Trajectory & Dynamic Time Warping (DTW) distance threshold
$\varepsilon$\_DTW \\
\end{longtable}
}

\paragraph{2.5.3 Version Pinning}\label{version-pinning}

AVP implementation version = P1 arXiv technical report commit hash
\texttt{\textless{}P1-AVP-vX.Y\textgreater{}}. P2 reproduction package
embeds complete AVP source code, decoupled from P1.

\subsubsection{2.6 LRCA Engineering Attribution
Layer}\label{lrca-engineering-attribution-layer}

\paragraph{2.6.1 Five Likely Root Causes
(Open)}\label{five-likely-root-causes-open}

{\def\LTcaptype{none} % do not increment counter
\begin{longtable}[]{@{}ll@{}}
\toprule\noalign{}
Root Cause & Meaning \\
\midrule\noalign{}
\endhead
\bottomrule\noalign{}
\endlastfoot
C1 & True Semantic Failure (what P2 seeks) \\
C2 & Numerical Tolerance Perturbation \\
C3 & Out-of-Distribution (OOD) \\
C4 & Statistical Assumption Violation \\
C5 & Mutator Artifact \\
\end{longtable}
}

\paragraph{2.6.2 Three-Layer Diagnosis + Artifact Pre-scan
(L0)}\label{three-layer-diagnosis-artifact-pre-scan-l0}

\begin{itemize}
\tightlist
\item
  \textbf{L0 Artifact Pre-scan}: Double-blind review (§4.2.4 dual LLM
  cross-source + 20\% human sampling)
\item
  \textbf{L1 Tolerance Robustness}: N=20 replicates, fail ratio $\geq$ 0.80
  considered stable
\item
  \textbf{L2 OOD Triage}: For C/D classes, distinguish valid domain
  D\_S\^{}valid from OOD region
\item
  \textbf{L3 Statistical Assumption Baseline}: For Wilcoxon/DTW,
  pre-check Independent and Identically Distributed (IID)/stationarity
  on PUT's own samples
\end{itemize}

\paragraph{2.6.3 Decision Tree}\label{decision-tree}

\begin{verbatim}
For each s' $\in$ killed_{i,k,j}:
  L1: fail ratio < 0.80 $\to$ C2
       otherwise $\to$ L2
  L2 (C/D classes): fail only in OOD $\to$ C3
       otherwise $\to$ L3
  L3 (B/D classes + Wilcoxon/DTW): assumption violation $\to$ C4
       otherwise $\to$ artifact recheck
  Artifact evidence $\to$ C5
       otherwise $\to$ C1
\end{verbatim}

Multi-label priority: \textbf{C5 \textgreater{} C4 \textgreater{} C3
\textgreater{} C2 \textgreater{} C1}.

\paragraph{2.6.4 Output Quantities}\label{output-quantities}

\begin{verbatim}
C1_share_{i,k,j}      := |{s' $\in$ killed : root_cause = C1}| / |killed|
suspect_share_{i,k,j} := 1 − C1_share_{i,k,j}
\end{verbatim}

LRCA \textbf{does not modify the SMS formula}; the killed set does not
exclude suspects.

\subsubsection{2.7 Formalization Hierarchy
Diagram}\label{formalization-hierarchy-diagram}

\begin{verbatim}
AVP (component, P1) -$\to$ equiv & killed (construction, P2 internal extension)
              -$\to$ SMS (metric, classical structure)
              -$\to$ LRCA (description, P2 engineering attribution layer, not in SMS)
\end{verbatim}

\begin{center}\rule{0.5\linewidth}{0.5pt}\end{center}

\subsection{Section 3 $\cdot$ Experimental Subjects and the 60-Cell
Operator-Program Instantiation
Matrix}\label{section-3-experimental-subjects-and-the-60-cell-operator-program-instantiation-matrix}

\subsubsection{3.1 Experimental Subject Selection (12 Programs Under
Test, Inherited from P1 {[}Meng Li et al., Progress in Nuclear Energy,
under
review{]})}\label{experimental-subject-selection-12-programs-under-test-inherited-from-p1-meng-li-et-al.-progress-in-nuclear-energy-under-review}

{\def\LTcaptype{none} % do not increment counter
\begin{longtable}[]{@{}
  >{\raggedright\arraybackslash}p{(\linewidth - 8\tabcolsep) * \real{0.2000}}
  >{\raggedright\arraybackslash}p{(\linewidth - 8\tabcolsep) * \real{0.2000}}
  >{\raggedright\arraybackslash}p{(\linewidth - 8\tabcolsep) * \real{0.2000}}
  >{\raggedright\arraybackslash}p{(\linewidth - 8\tabcolsep) * \real{0.2000}}
  >{\raggedright\arraybackslash}p{(\linewidth - 8\tabcolsep) * \real{0.2000}}@{}}
\toprule\noalign{}
\begin{minipage}[b]{\linewidth}\raggedright
Class cls
\end{minipage} & \begin{minipage}[b]{\linewidth}\raggedright
Program i
\end{minipage} & \begin{minipage}[b]{\linewidth}\raggedright
Name
\end{minipage} & \begin{minipage}[b]{\linewidth}\raggedright
Representative Mathematical Structure
\end{minipage} & \begin{minipage}[b]{\linewidth}\raggedright
Scale
\end{minipage} \\
\midrule\noalign{}
\endhead
\bottomrule\noalign{}
\endlastfoot
\textbf{A Numerical} & A1 & Lorenz system Ordinary Differential Equation
(ODE) integration & Nonlinear ODE system & \textasciitilde150 LOC \\
& A2 & LU decomposition & Linear algebra, numerical stability &
\textasciitilde80 LOC \\
& A3 & Finite Difference Method (FDM) 1D heat conduction & Parabolic
Partial Differential Equation (PDE) time-stepping & \textasciitilde200
LOC \\
\textbf{B Probabilistic} & B1 & Beta-Binomial conjugate inference &
Analytical posterior & \textasciitilde60 LOC \\
& B2 & Markov Chain Monte Carlo (MCMC) Metropolis-Hastings & Markov
chain sampling & \textasciitilde250 LOC \\
& B3 & Monte Carlo integration & Importance sampling &
\textasciitilde100 LOC \\
\textbf{C Surrogate} & C1 & Gaussian Process Regression (GPR) & Kernel
methods & \textasciitilde300 LOC \\
& C2 & Polynomial Chaos Expansion (PCE) & Orthogonal basis &
\textasciitilde250 LOC \\
& C3 & Neural Network Surrogate (NN-Surrogate, NN-Surr) & Multi-layer
perceptron substitution & \textasciitilde400 LOC \\
\textbf{D Machine Learning} & D1 & Multi-Layer Perceptron (MLP) &
Backpropagation & \textasciitilde350 LOC \\
& D2 & Support Vector Machine (SVM) & Convex optimization &
\textasciitilde200 LOC \\
& D3 & Logistic Regression (LR) & Maximum likelihood &
\textasciitilde120 LOC \\
\end{longtable}
}

\paragraph{3.1.1 Selection Rationale
(Self-Contained)}\label{selection-rationale-self-contained}

The selection of 12 PUTs not only inherits the work of P1 {[}Meng Li et
al., Progress in Nuclear Energy, under review{]}, but this paper also
independently justifies their representativeness. The justification
unfolds across three dimensions:

\textbf{(a) Library stack coverage}: The implementation stack of the 12
PUTs in this paper covers the 3 major mainstream libraries of Python
scientific computing + the numpy foundation layer:

{\def\LTcaptype{none} % do not increment counter
\begin{longtable}[]{@{}
  >{\raggedright\arraybackslash}p{(\linewidth - 4\tabcolsep) * \real{0.3333}}
  >{\raggedright\arraybackslash}p{(\linewidth - 4\tabcolsep) * \real{0.3333}}
  >{\raggedright\arraybackslash}p{(\linewidth - 4\tabcolsep) * \real{0.3333}}@{}}
\toprule\noalign{}
\begin{minipage}[b]{\linewidth}\raggedright
Library
\end{minipage} & \begin{minipage}[b]{\linewidth}\raggedright
Involved PUTs
\end{minipage} & \begin{minipage}[b]{\linewidth}\raggedright
Library Version (This Paper)
\end{minipage} \\
\midrule\noalign{}
\endhead
\bottomrule\noalign{}
\endlastfoot
numpy (linear algebra / array operations) & A1, A2, A3, B1, B2, B3 &
2.4.4 \\
scipy (optimization / integration / statistics) & A1(integrate),
A2(linalg), B2(stats) & 1.17.1 \\
scikit-learn (surrogate / ML) & C1, C2, C3, D1, D2, D3 & 1.8.0 \\
\end{longtable}
}

These 3 libraries + numpy constitute the de facto foundation stack for
Python scientific computing software, covering all scientific computing
modules in the PyPI download top-50 (2026-04 data). The 12 PUTs in this
paper do not cover GPU/distributed computing stacks (JAX / CuPy / dask)
and domain-specific libraries (BioPython / Astropy / RDKit), which is a
limitation explicitly declared in §7.2.1 R5, reserved for the P3 paper
(broader coverage).

\textbf{(b) Mathematical structure coverage}: Across 4 classes
(numerical / probabilistic / surrogate / ML), the mathematical
structures instantiated by the 12 PUTs include:

\begin{itemize}
\tightlist
\item
  ODE / PDE time-stepping (A1, A3)
\item
  Direct linear algebra (A2)
\item
  Bayesian analytical inference + MCMC + Monte Carlo integration (B1,
  B2, B3)
\item
  Kernel methods + orthogonal polynomials + neural network surrogates
  (C1, C2, C3)
\item
  Convex optimization + backpropagation + maximum likelihood (D1, D2,
  D3)
\end{itemize}

These 12 structures cover 8 of the 12 chapters in the main table of
contents of Numerical Recipes (Press et al.~2007). \textbf{The 4
uncovered chapters}: (1) advanced PDE solvers (FEM, finite-volume,
spectral), this paper's A3 only covers 1-D explicit FDM; (2) FFT /
spectral methods; (3) optimization (interior-point, trust-region, etc.),
D2 SVM only touches a subset of convex optimization; (4) symbolic
computation / computer algebra. The PUTs corresponding to these 4
chapters are important in industrial-grade scientific computing software
(CFD / quantum chemistry / signal processing), but are not covered in
this paper, reserved for the P3 paper scaling study (§7.2.1 R5 merged
declaration with this section).

\textbf{(c) Comparison with existing mutation testing benchmarks}:

{\def\LTcaptype{none} % do not increment counter
\begin{longtable}[]{@{}
  >{\raggedright\arraybackslash}p{(\linewidth - 4\tabcolsep) * \real{0.3333}}
  >{\raggedright\arraybackslash}p{(\linewidth - 4\tabcolsep) * \real{0.3333}}
  >{\raggedright\arraybackslash}p{(\linewidth - 4\tabcolsep) * \real{0.3333}}@{}}
\toprule\noalign{}
\begin{minipage}[b]{\linewidth}\raggedright
Benchmark
\end{minipage} & \begin{minipage}[b]{\linewidth}\raggedright
Covered PUTs
\end{minipage} & \begin{minipage}[b]{\linewidth}\raggedright
Relationship to This Paper's 12 PUTs
\end{minipage} \\
\midrule\noalign{}
\endhead
\bottomrule\noalign{}
\endlastfoot
\textbf{DeepCrime} (Humbatova, Jahangirova \& Tonella, ISSTA 2021) &
Deep learning systems (Keras/TensorFlow, real-fault-based) & This
paper's class D (D1/D2/D3, sklearn ML kernels) shares topical overlap
with DeepCrime on the ML category (specific frameworks and fault models
differ) \\
\textbf{Defects4J} (Just, Jalali \& Ernst, ISSTA 2014) & General Java
fault database & No direct PUT-selection overlap with this paper; cited
only as a baseline mutation-as-fault-proxy reference (§1.3.2) \\
\textbf{mutmut / cosmic-ray demo PUTs} & General Python (no domain
focus) & This paper significantly extends scientific-computing domain
focus \\
\end{longtable}
}

The 12 PUTs in this paper share topical overlap with DeepCrime on the ML
subset (D); classes A (numerical) / B (probabilistic) / C (surrogate)
are this paper's unique extensions relative to the Python LLM-mutant
literature.

\textbf{(d) Balance between scale and representativeness}: Each PUT is
50-400 LOC, with function signatures standardized to
\texttt{program(x:\ float)\ $\to$\ float} (facilitating AVP invocation and
60-cell matrix generation). The scale is smaller than industrial code
(typically 1-10 KLOC), but retains the complete semantics of the core
mathematical structure of each class.

\textbf{Limitation: The \texttt{program(x:\ float)\ $\to$\ float} signature
simplification is a substantive constraint} (not a purely engineering
tradeoff). The input of industrial-grade scientific computing software
is typically high-dimensional structures such as mesh / state-vector /
tensor (CFD grids, MD particle states, FEM element matrices), and the
output may be multi-component fields or time-series trajectories. The
scalarized PUTs in this paper impose an upper bound on the semantic
complexity of mutants, potentially systematically underestimating SMS
and cross-class differences on industrial PUTs. This signature
constraint is explicitly declared in §7.2.1 R5 and §7.5 final
limitation; the P3 paper will validate the portability of SMS on
industrial-grade PUTs.

\textbf{Statistical layer support}: The Cliff's $\delta$ in §5.7.2 and the
Friedman $\chi$$^{2}$ in §5.8.4 both provide directional evidence (the latter
significant, the former H2 rejected). The scale of 12 PUTs works for H1
(12/12 PUTs aligned \textgreater{} 0) and H4 sign test, but is clearly
underpowered for mixed-effects (§5.8.3 Singular) and Spearman $\rho$ (n=12,
p=0.74) (§7.2.2 R6). This paper does not make a strong claim that ``12
PUTs are sufficient to support all RQs.''

\subsubsection{3.2 5 Operator Semantic Injection Rules (MUT Content
Instantiation)}\label{operator-semantic-injection-rules-mut-content-instantiation}

This section defines the 5 classes of \textbf{meta-mutation operators}
in this paper --- not operators directly targeting a specific PUT, but
operator families / operator templates. For different types of programs
(a numeric / b probabilistic / c surrogate / d ML), each meta-operator
requires \textbf{specialization} to be instantiated into executable
mutant generation rules.

\textbf{Specialization rule examples}: - \textbf{HP} (hyperparameter)
specializes on class a PUTs to ``change numerical algorithm tolerance /
max\_iter''; on class c to ``change GPR kernel noise\_level /
length\_scale''; on class d to ``change MLP hidden\_dim / dropout'' -
\textbf{OS} (API replacement) specializes on class a to ``swap numerical
linear algebra APIs'' (e.g., \texttt{det} $\leftrightarrow$ \texttt{sum(diag)}); on
class b to ``swap probability distribution sampling APIs''; on class c
to ``swap surrogate classes (GPR $\leftrightarrow$ RBF $\leftrightarrow$ NN)'' - \textbf{TF} (numerical
transform) specializes on class a to ``change integration order (RK4 $\to$
Euler)''; on class b to ``change MC estimator''

§3.3 provides the full specialization grid of 5 meta-operators on 12
PUTs $\times$ 5 MPs (60 cells). \textbf{Each mut\_X subsection in §3.2
describes the abstract definition of each meta-operator; §3.3 is their
concrete specialization}.

\textbf{Hierarchical status}: The 5 meta-operators are defined at Layer
1 (necessary conditions §3.2.0); their specialization instances (60
cells) are the reference objects for Layer 3 mutant tracing (§3.2.6.3).

\paragraph{3.2.0 Necessary Conditions for Semantic Mutation (Layer 1 ---
Definitional)}\label{necessary-conditions-for-semantic-mutation-layer-1-definitional}

\begin{quote}
Added in P2 R2 revision. Provides formal criteria for ``what counts as
semantic mutation'' as the methodological foundation for the 5 classes
of semantic operators (CE/OS/HP/TF/SI) in §3.2.
\end{quote}

\textbf{Definition (Semantic Mutation Criteria)}: A mutant
\texttt{s\textquotesingle{}\ =\ mut\_j(S\_i)} is \textbf{semantic
mutation} if and only if it satisfies at least one of the following
three conditions:

\begin{enumerate}
\def\labelenumi{(\alph{enumi})}
\item
  \textbf{Cross-function-boundary replacement}: The AST node operated on
  by the mutator crosses at least one function-call or module-import
  boundary (e.g., \texttt{np.linalg.det(M)} $\to$
  \texttt{np.sum(np.diag(M))}, replacing 1 function call with a
  composition of 2 functions);
\item
  \textbf{Carries domain knowledge}: The legality of the mutation
  depends on mathematical/physical/statistical knowledge of the
  program's domain, not purely syntactic type preservation (e.g., GPR
  \texttt{noise\_level=1e-4\ $\to$\ 1e-1} knows this is a hyperparameter
  rather than a random literal constant);
\item
  \textbf{Changes algorithmic class}: The mutation changes the algorithm
  class implemented by the program (e.g., RK4 $\to$ Euler changes
  integration order, dropout-prob 0.5 $\to$ 0 changes ML model class).
\end{enumerate}

Otherwise it is \textbf{syntactic mutation} (AST-local + domain-agnostic
+ does not change algorithm class).

\textbf{Correspondence with the 5 operator classes in §3.2.1-§3.2.5}
(each operator satisfies at least one condition):

{\def\LTcaptype{none} % do not increment counter
\begin{longtable}[]{@{}
  >{\raggedright\arraybackslash}p{(\linewidth - 8\tabcolsep) * \real{0.2000}}
  >{\raggedright\arraybackslash}p{(\linewidth - 8\tabcolsep) * \real{0.2000}}
  >{\raggedright\arraybackslash}p{(\linewidth - 8\tabcolsep) * \real{0.2000}}
  >{\raggedright\arraybackslash}p{(\linewidth - 8\tabcolsep) * \real{0.2000}}
  >{\raggedright\arraybackslash}p{(\linewidth - 8\tabcolsep) * \real{0.2000}}@{}}
\toprule\noalign{}
\begin{minipage}[b]{\linewidth}\raggedright
Operator Class
\end{minipage} & \begin{minipage}[b]{\linewidth}\raggedright
(a) Cross-function boundary
\end{minipage} & \begin{minipage}[b]{\linewidth}\raggedright
(b) Domain knowledge
\end{minipage} & \begin{minipage}[b]{\linewidth}\raggedright
(c) Algorithm class
\end{minipage} & \begin{minipage}[b]{\linewidth}\raggedright
Primary satisfied condition
\end{minipage} \\
\midrule\noalign{}
\endhead
\bottomrule\noalign{}
\endlastfoot
\textbf{CE} Constant perturbation & $\times$ & $\bigtriangleup$ (domain semantics of
constants) & $\times$ & Partial (b); weakest condition \\
\textbf{OS} API replacement & \checkmark & \checkmark (mathematical equivalence between
APIs) & $\bigtriangleup$ (sometimes changes algorithm class) & (a)+(b) \\
\textbf{HP} Hyperparameter & $\times$ & \checkmark (semantic dimension of
hyperparameters) & $\bigtriangleup$ (extreme HP changes algorithm) & (b)+partial (c) \\
\textbf{TF} Numerical transform & $\bigtriangleup$ (sometimes cross-function) & \checkmark
(order/convergence of numerical methods) & \checkmark (changes
integration/interpolation order) & (b)+(c) \\
\textbf{SI/CF} Structural injection & $\bigtriangleup$ (sometimes cross-control-flow) &
\checkmark (algorithmic intent of control flow) & \checkmark (changes algorithm skeleton)
& (b)+(c) \\
\end{longtable}
}

\textbf{Only the CE class partially satisfies the necessary conditions}
(mainly relying on (b) domain semantics, neither (a) nor (c) is strong).
This is precisely why CE overlaps with the NumberReplacer of syntactic
tools in the operator-level comparison table in §3.2.6.1 --- CE is a
semantic/syntactic boundary class. OS / HP / TF / SI strongly satisfy
one or more of (a)(b)(c), and \textbf{structurally} do not belong to the
capability space of syntactic mutators.

\textbf{Hierarchical status}: The (a)(b)(c) in this section are
\textbf{Layer 1 necessary conditions} of the P2 methodology; the
equivalence criteria E1 $\wedge$ E2 in §2.3 / §4.4 are \textbf{Layer 2
instantiation} of the necessary conditions (see §4.4 revision); the
mutant tracing empirics in §3.2.6.3 are \textbf{Layer 3 application}
(see that section).

\paragraph{mut\_C
(Conservation-breaking)}\label{mut_c-conservation-breaking}

\begin{itemize}
\tightlist
\item
  A1 Lorenz: Add $\varepsilon$\_drift drift to right-hand side, Hamiltonian
  monotonically drifts slowly
\item
  A2 LU: Decomposition loop omits subtraction of k+1-th row multiplier,
  breaking determinant conservation
\item
  B1 Beta-Bin: Posterior update omits normalization (total probability $\neq$
  1)
\item
  C1 GPR: Covariance matrix positive-definiteness omits diagonal term
\item
  D1 MLP: Backpropagation gradient summation omits one term
\end{itemize}

\paragraph{mut\_M
(Monotonicity-breaking)}\label{mut_m-monotonicity-breaking}

\begin{itemize}
\tightlist
\item
  A3 FDM: $\Delta$t coefficient occasionally takes negative value
\item
  B2 MCMC: Acceptance rate min(1, ratio) changed to min(0.95, ratio)
\item
  C2 PCE: High-order coefficient sorting inserts inversion pair
\item
  D2 SVM: Decision function sign flips near boundary
\end{itemize}

\paragraph{mut\_G
(Convergence-breaking)}\label{mut_g-convergence-breaking}

\begin{itemize}
\tightlist
\item
  A1 Lorenz: Fourth-order Runge-Kutta changed to 1.5-order hybrid
\item
  A3 FDM: Second-order difference term replaced with first-order
\item
  B3 MC integration: Sample size doubling does not decay variance by 1/N
\item
  C3 NN-Surr: Training epoch truncation
\end{itemize}

\paragraph{mut\_T
(Trajectory-distorting)}\label{mut_t-trajectory-distorting}

\begin{itemize}
\tightlist
\item
  A1 Lorenz: State vector swaps y and z components
\item
  B2 MCMC: Insert independent sampling segment in chain
\item
  C3 NN-Surr: Training target adds slow phase shift
\item
  D1 MLP: Hidden layer activation by periodic mask
\end{itemize}

\paragraph{mut\_F
(Fidelity-order-breaking)}\label{mut_f-fidelity-order-breaking}

\begin{itemize}
\tightlist
\item
  A2 LU: Partial pivoting occasionally degrades to no pivoting
\item
  C1 GPR: Length scale hyperparameter occasionally switches to coarse
  prior
\item
  C2 PCE: High-order term truncation randomly retains low-order
\item
  D3 LR: Regularization coefficient occasionally takes large value
\end{itemize}

\paragraph{3.2.6 Capability Relationship Between P2 Semantic Operators
and Existing Python Mutation Testing
Tools}\label{capability-relationship-between-p2-semantic-operators-and-existing-python-mutation-testing-tools}

Classic syntactic mutation tools (Jia \& Harman 2011 survey) operate at
the AST node level, independent of the program's domain semantics. We
compare the operator capabilities of P2's 5 semantic operator classes
(CE / OS / HP / TF / SI, instantiated as the 5 mut\_k categories in
§3.2.1-5) with mainstream Python mutation testing tools:

{\def\LTcaptype{none} % do not increment counter
\begin{longtable}[]{@{}
  >{\raggedright\arraybackslash}p{(\linewidth - 4\tabcolsep) * \real{0.3333}}
  >{\raggedright\arraybackslash}p{(\linewidth - 4\tabcolsep) * \real{0.3333}}
  >{\raggedright\arraybackslash}p{(\linewidth - 4\tabcolsep) * \real{0.3333}}@{}}
\toprule\noalign{}
\begin{minipage}[b]{\linewidth}\raggedright
P2 Operator Class
\end{minipage} & \begin{minipage}[b]{\linewidth}\raggedright
Python tools (mutmut / cosmic-ray / mutpy) coverage
\end{minipage} & \begin{minipage}[b]{\linewidth}\raggedright
Reason tools do not cover
\end{minipage} \\
\midrule\noalign{}
\endhead
\bottomrule\noalign{}
\endlastfoot
\textbf{CE} Constant perturbation & \checkmark Covered (CRP operator) & Tool and
P2 operator overlap, this paper retains CE as baseline \\
\textbf{OS} API replacement (e.g., \texttt{np.linalg.det} $\to$
\texttt{np.sum(np.diag)}) & $\bigtriangleup$ Mostly not covered (§3.2.6.3 empirics
88.33\% AST-disjoint; a small number of low-complexity OS
sub-expressions occasionally hit by tools) & Tool AST replacement cannot
do semantically equivalent but algebraically different replacements
across function boundaries; \texttt{np.linalg.det(M)} $\to$
\texttt{np.sum(np.diag(M))} is scientific computing domain knowledge,
tools lack this knowledge \\
\textbf{HP} Hyperparameter semantic change (e.g., GPR
\texttt{noise\_level}, MLP \texttt{max\_iter}) & $\times$ Not covered & Tools
do not recognize semantic dimensions of sklearn / scipy model
hyperparameters, can only blindly hit literal constants (equivalent to
CRP) \\
\textbf{TF} Numerical transform (e.g., RK4 $\to$ 1.5-order hybrid
integration) & $\bigtriangleup$ Partial (AOR changes operators) & Tool AOR can only
replace \texttt{+/-/*/$\div$}, cannot change semantic order of integration
methods \\
\textbf{SI} Structural injection (control flow semantic intent) & $\times$ Not
covered & Tools do not mutate semantic intent of if-else / loops, can
only mechanically change conditions \\
\end{longtable}
}

\textbf{Of the 5 operator classes, only CE overlaps with the CRP of
traditional tools; the remaining 4 classes (OS / HP / TF / SI) are
semantic-layer supplements beyond the capability of first-order
syntactic tools}. This paper uses LLM-generated semantic operators to
cover these 4 classes, which is the engineering connotation of
``domain-semantic mutation operators'' in innovation attribution C-I in
§1.6.1.

\textbf{Higher-Order Mutation (HOM) caveat}: Jia \& Harman (2009 SBSE)
and Kintis et al.~(2018 STVR) proposed that HOM, by combining multiple
first-order syntactic mutations (e.g., AOR + SDL), could in principle
produce composite syntactic mutants partially equivalent to OS / HP / TF
/ SI. For example, AOR (\texttt{*} $\to$ \texttt{+}) + SDL (delete one line)
might simulate partial effects of OS API replacement on some PUTs. The
§5 empirics in this paper did not conduct comparative experiments on
HOM; §7 R12 (NEW) lists ``empirical testing of HOM equivalence'' as a
residual threat; the ``tool unreachability'' claim in §3.2.6 is strictly
limited to first-order syntactic tools (mutmut / cosmic-ray default
configurations belong to this class).

\textbf{Tool selection justification}: If comparative experiments are
conducted (§5.10 plan), both mutmut {[}Hovde 2018{]} + cosmic-ray
{[}Tomilin 2017{]} dual tools should be used simultaneously to exclude
single-tool operator set bias; mutpy {[}Hovstadius 2014{]} is not
selected due to Python 3.10+ incompatibility (testing sklearn/scipy
reports errors); Pitest is a Java tool, not comparable. Multi-tool
comparison can exclude reasonable reviewer concerns about ``subjective
tool selection.''

\textbf{Expected findings from comparative experiments} (based on the
argument that OS / HP / TF / SI are not covered by tools): The
aligned-cross Cliff's $\delta$ on 60-cell SMS for mutants produced by syntactic
tools should be significantly lower than the measured $\delta$ = 0.446 (v3b)
for semantic operators in this paper, because syntactic tools do not
distinguish MR-MP alignment, and their AST operations do not carry
semantic alignment information.

\subparagraph{3.2.6.0 Systematic vs Incidental: Syntactic Tool
Occasional Hits $\neq$ Semantic Mutation
Method}\label{systematic-vs-incidental-syntactic-tool-occasional-hits-semantic-mutation-method}

\begin{quote}
Added in P2 R2 revision. Responds to possible reviewer challenge: ``Some
mutants occasionally produced by syntactic tools may also cross function
boundaries (satisfying §3.2.0 (a)) or change algorithmic class ((c)),
aren't they byproducts but still semantic mutants?''
\end{quote}

\textbf{Thesis}: Satisfying the necessary conditions (a)(b)(c) in §3.2.0
is one of the \textbf{sufficient conditions} for semantic mutation, but
\textbf{only when satisfaction is design intent rather than stochastic
byproduct} does it constitute a systematic semantic mutation method.
Syntactic tools (mutmut / cosmic-ray) occasionally hitting (a)(c) with
12 default operators (§3.2.6.1) is a non-zero probability event --- but
\textbf{occasionality undermines two engineering functions of semantic
mutation}:

\textbf{(i) Deepening source code understanding} Semantic mutator design
requires understanding domain-level relationships of the program --- for
example, when designing the OS operator
\texttt{np.linalg.det(M)\ $\to$\ np.sum(np.diag(M))}, one must know: these
two APIs are equivalent on diagonal matrices, not equivalent on general
matrices, and the algebraic relationship between determinant and
diagonal sum is \texttt{det\ =\ $\prod$\ eigvals} while
\texttt{sum(diag)\ =\ trace\ =\ $\sum$\ eigvals}. Syntactic tool AST
traversal \textbf{does not require} this understanding; even if similar
replacements are occasionally produced, they do not constitute
systematic interpretation of source code semantics.

\textbf{(ii) Revealing deep faults} Domain semantic errors --- physical
constant errors, unit conversion errors, boundary condition errors,
hyperparameter semantic errors, numerical method order errors --- are
all \textbf{not} AST-local errors. Syntactic mutator design goals are to
trigger syntactic faults (operator typos / off-by-one / negation flips);
the probability of occasionally hitting domain errors $\ll$ the probability
of designed triggering, and lacks repeatability. Semantic mutator design
goals \textbf{directly} correspond to these deep fault classes.

\textbf{Conclusion}: Syntactic tools occasionally producing mutants
satisfying §3.2.0 (a)(b)(c) is a stochastic byproduct --- neither
repeatable (the same tool with the same seed may not hit) nor carrying
understanding/fault-revealing engineering value. Systematic semantic
mutation requires (a)(b)(c) to be design intent. This is the
\textbf{positive reinforcement} of the ``AST-local + domain-agnostic''
criterion in the operator-level comparison table in §3.2.6.1: not only
is the operator set structure of syntactic tools unable to reach P2
necessary conditions, even occasional hits do not constitute systematic
semantic mutation in the methodological sense.

\subparagraph{3.2.6.1 Operator-Level Comparison Table (R-15
Response)}\label{operator-level-comparison-table-r-15-response}

\begin{quote}
Reviewer R-15 requires upgrading the ``tool unreachability'' claim in
§3.2.6 from categorical argument to \textbf{operator-level} comparative
evidence. This section lists all entries in the default operator sets of
mutmut and cosmic-ray, mapping one-to-one with the 37 operators in P2's
\texttt{operator\_registry.py}, to confirm that the 4 classes beyond CE
(OS / HP / TF / SI) are completely inexpressible at the operator level.
\end{quote}

{\def\LTcaptype{none} % do not increment counter
\begin{longtable}[]{@{}
  >{\raggedright\arraybackslash}p{(\linewidth - 8\tabcolsep) * \real{0.2000}}
  >{\raggedright\arraybackslash}p{(\linewidth - 8\tabcolsep) * \real{0.2000}}
  >{\raggedright\arraybackslash}p{(\linewidth - 8\tabcolsep) * \real{0.2000}}
  >{\raggedright\arraybackslash}p{(\linewidth - 8\tabcolsep) * \real{0.2000}}
  >{\raggedright\arraybackslash}p{(\linewidth - 8\tabcolsep) * \real{0.2000}}@{}}
\toprule\noalign{}
\begin{minipage}[b]{\linewidth}\raggedright
Tool Operator (mutmut + cosmic-ray)
\end{minipage} & \begin{minipage}[b]{\linewidth}\raggedright
Operator Instance
\end{minipage} & \begin{minipage}[b]{\linewidth}\raggedright
Operated AST Node
\end{minipage} & \begin{minipage}[b]{\linewidth}\raggedright
P2 Correspondence
\end{minipage} & \begin{minipage}[b]{\linewidth}\raggedright
Covers P2 Class
\end{minipage} \\
\midrule\noalign{}
\endhead
\bottomrule\noalign{}
\endlastfoot
\texttt{NumberReplacer} / number constant replacement & \texttt{1} $\to$
\texttt{2}, \texttt{0.05} $\to$ \texttt{0.06} &
\texttt{Num}/\texttt{Constant} & CE partial & $\bigtriangleup$ Literal values only \\
\texttt{ReplaceArithmeticOperator} & \texttt{+} $\to$ \texttt{-}, \texttt{*}
$\to$ \texttt{/} & \texttt{BinOp} & --- & $\times$ No correspondence \\
\texttt{ReplaceComparisonOperator} & \texttt{\textless{}} $\to$
\texttt{\textgreater{}=}, \texttt{==} $\to$ \texttt{!=} & \texttt{Compare} &
--- & $\times$ No correspondence \\
\texttt{ReplaceLogicalOperator} (\texttt{and}$\leftrightarrow$\texttt{or}) &
\texttt{a\ and\ b} $\to$ \texttt{a\ or\ b} & \texttt{BoolOp} & --- & $\times$ No
correspondence \\
\texttt{ReplaceUnaryOperator} & \texttt{-x} $\to$ \texttt{+x} &
\texttt{UnaryOp} & --- & $\times$ No correspondence \\
\texttt{ReplaceTrueFalse} & \texttt{True} $\to$ \texttt{False} &
\texttt{NameConstant} & CE keyword & $\bigtriangleup$ Bool literal only \\
\texttt{BreakContinueReplacer} & \texttt{break} $\leftrightarrow$ \texttt{continue} &
\texttt{Break}/\texttt{Continue} & --- & $\times$ No correspondence \\
\texttt{RemoveDecorator} & \texttt{@cache} deletion &
\texttt{decorator\_list} & --- & $\times$ No correspondence \\
\texttt{RemoveExceptHandler} & \texttt{except\ E:\ ...} deletion &
\texttt{ExceptHandler} & --- & $\times$ No correspondence \\
\texttt{ZeroIterationForLoop} & \texttt{for\ x\ in\ xs:} $\to$
\texttt{for\ x\ in\ {[}{]}:} & \texttt{For} & --- & $\times$ No
correspondence \\
\texttt{ReplaceIfBlock} (\texttt{If\ True/False}) & \texttt{if\ cond:} $\to$
\texttt{if\ True:} & \texttt{If} & --- & $\times$ No correspondence \\
\texttt{MutateSubscript} & \texttt{a{[}i{]}} $\to$ \texttt{a{[}i+1{]}} &
\texttt{Subscript} & --- & $\times$ No correspondence \\
(No corresponding tool operator) & --- & --- & \textbf{OS} API
replacement & $\bigtriangleup$ 88.33\% disjoint (§3.2.6.3 empirics; small number of
low-complexity OS sub-expressions occasionally hit by BinOp) \\
(No corresponding tool operator) & --- & --- & \textbf{HP}
Hyperparameter semantic change & $\times$ Tool inexpressible \\
(No corresponding tool operator) & --- & --- & \textbf{TF} Numerical
method semantic change & $\times$ Tool inexpressible \\
(No corresponding tool operator) & --- & --- & \textbf{SI/CF} Structural
injection & $\times$ Tool inexpressible \\
\end{longtable}
}

\textbf{Operator-level conclusion}: All 12 classes in the tool default
operator set remain at AST-local operations (BinOp / Compare / BoolOp /
UnaryOp / NameConstant / Subscript / If / For / Break / Continue /
decorator / except). \textbf{Categorically}, no entry recognizes
semantic dimensions of sklearn / scipy model object hyperparameters (HP,
§3.2.6.3 empirics 0/72), semantic order of numerical methods (TF, 0/54),
or intent semantics of control flow (SI/CF, 0/33). For OS
(cross-function-boundary API replacement), the tool carries no domain
knowledge, but low-complexity OS sub-expressions can be occasionally hit
by BinOp and similar operators (§3.2.6.3 empirics: 88.33\% disjoint).
\textbf{This is structural unreachability, not a matter of operator set
size --- even if the tool operator set is expanded from 12 to 100, as
long as each entry remains AST-local and domain-agnostic, the HP/TF/SI
systematic semantic operator classes remain inexpressible, and the
systematic hit rate for OS is bounded below by the empirical 88.33\%
disjointness}.

\subparagraph{3.2.6.2 cosmic-ray on a1 Single PUT Empirical Supplement
(Optional, Future-Work
Hook)}\label{cosmic-ray-on-a1-single-put-empirical-supplement-optional-future-work-hook}

\begin{quote}
The above operator-level comparison is already sufficient evidence for
the §3.2.6 thesis. This provides an \textbf{optional lightweight
empirical}, applicable to scenarios where R-15 reviewer further requests
``actual run numbers.''
\end{quote}

\texttt{scripts/run\_cosmic\_ray\_a1.sh} (NEW) encapsulates the
end-to-end cosmic-ray process on a1 PUT:

\begin{enumerate}
\def\labelenumi{\arabic{enumi}.}
\tightlist
\item
  Install \texttt{cosmic-ray};
\item
  Configure \texttt{tests/puts/test\_a1.py} as baseline test with
  \texttt{cosmic-ray.toml};
\item
  Run \texttt{cosmic-ray\ exec} for all mutants;
\item
  Report (mutants\_generated, killed, survived, incompetent) quadruple.
\end{enumerate}

\textbf{Expected result} (empirical confirmation of operator-level
comparison): - mutants\_generated: tens of magnitude (a1 file 934 bytes,
limited AST nodes); - If fully classified, $\geq$ 90\% belong to BinOp /
Compare and other 12 classes; - killed-by-test\_a1.py ratio: depends on
P2 unit test granularity; P2 tests are designed as PUT output shape/type
sanity checks, \textbf{most mutants will not be killed by this test
suite} --- this is precisely the empirical manifestation of the §3.2.6
thesis: mutants produced by syntactic tools are not aligned with
MR-violation detection objectives.

Full 12 PUT $\times$ cosmic-ray ablation is reserved for R2 revision or P3
paper (§7.2.1 R5 scaling study).

\subparagraph{3.2.6.3 Mutant Tracing Empirics (Layer 3 ---
Applied)}\label{mutant-tracing-empirics-layer-3-applied}

\begin{quote}
Added in P2 R2 revision, NEW-MAJOR-1 generalization closed (2026-05-02
12-PUT full measurement). Layer 3 application layer: use §2.3 / §4.4
equivalence detection tools (Layer 2 instantiation of §3.2.0 necessary
conditions) to trace P2 mutant set for syntactic-mutant,
\textbf{positive empirical} argument that P2 mutants are not a subset
classification of syntactic mutants.
\end{quote}

\textbf{Experimental design}: Full 12 PUTs (a1/a2/a3 / b1/b2/b3 /
c1/c2/c3 / d1/d2/d3) empirics: 1. Take P2 v4 cross-source mutant pool
(\texttt{data/mutants/\$\{PUT\_ID\}\_pool\_v4/}, 292 mutants across 12
PUTs); 2. Run cosmic-ray default operators for each PUT
(\texttt{scripts/run\_cosmic\_ray\_put.sh\ \$\{PUT\_ID\}}, 1250
syntactic mutants across 12 PUTs); 3. For each mutant, use
\texttt{ast.dump(annotate\_fields=False,\ include\_attributes=False)}
for normalized AST string, perform set difference analysis
(\texttt{scripts/p2\_vs\_syntactic\_ast\_diff\_batch.py}); 4. Report
overall overlap rate and per-operator-class breakdown.

\textbf{Empirical results}
(\texttt{data/results/cosmic\_ray\_12put\_ast\_diff.json}):

{\def\LTcaptype{none} % do not increment counter
\begin{longtable}[]{@{}ll@{}}
\toprule\noalign{}
Metric & Value \\
\midrule\noalign{}
\endhead
\bottomrule\noalign{}
\endlastfoot
Total P2 mutants across 12 PUTs & 292 \\
Total cosmic-ray syntactic mutants across 12 PUTs & 1250 \\
AST-normalized overlapping mutants & 15 \\
\textbf{Overall overlap rate} & \textbf{0.0514} \\
\end{longtable}
}

\textbf{Per-operator-class overlap rate (12-PUT aggregate)}:

{\def\LTcaptype{none} % do not increment counter
\begin{longtable}[]{@{}
  >{\raggedright\arraybackslash}p{(\linewidth - 8\tabcolsep) * \real{0.2000}}
  >{\raggedright\arraybackslash}p{(\linewidth - 8\tabcolsep) * \real{0.2000}}
  >{\raggedright\arraybackslash}p{(\linewidth - 8\tabcolsep) * \real{0.2000}}
  >{\raggedright\arraybackslash}p{(\linewidth - 8\tabcolsep) * \real{0.2000}}
  >{\raggedright\arraybackslash}p{(\linewidth - 8\tabcolsep) * \real{0.2000}}@{}}
\toprule\noalign{}
\begin{minipage}[b]{\linewidth}\raggedright
Class
\end{minipage} & \begin{minipage}[b]{\linewidth}\raggedright
n\_p2
\end{minipage} & \begin{minipage}[b]{\linewidth}\raggedright
n\_overlap
\end{minipage} & \begin{minipage}[b]{\linewidth}\raggedright
overlap\_rate
\end{minipage} & \begin{minipage}[b]{\linewidth}\raggedright
Interpretation
\end{minipage} \\
\midrule\noalign{}
\endhead
\bottomrule\noalign{}
\endlastfoot
\textbf{HP} & 72 & 0 & \textbf{0.0000} & Structurally unreachable \checkmark \\
\textbf{SI} & 33 & 0 & \textbf{0.0000} & Structurally unreachable \checkmark \\
\textbf{TF} & 54 & 0 & \textbf{0.0000} & Structurally unreachable \checkmark \\
CE & 64 & 5 & 0.0781 & Boundary class (§3.2.0 only partially satisfies
(b)) \\
OS & 60 & 7 & 0.1167 & Partial instance hits (see interpretation) \\
CF & 9 & 3 & 0.3333 & b2 only, extremely small sample \\
\end{longtable}
}

Per-PUT detailed distribution (only listing PUTs with non-zero overlap):

{\def\LTcaptype{none} % do not increment counter
\begin{longtable}[]{@{}llll@{}}
\toprule\noalign{}
PUT & n\_p2 & n\_overlap & Hit operator class and LLM source \\
\midrule\noalign{}
\endhead
\bottomrule\noalign{}
\endlastfoot
a2 & 27 & 3 & CE: Claude $\times$ 3 \\
a3 & 18 & 3 & OS: DeepSeek $\times$ 3 \\
b2 & 27 & 3 & CF: DeepSeek $\times$ 3 \\
b3 & 20 & 6 & CE: DeepSeek $\times$ 2 + OS: Claude $\times$ 1, DeepSeek $\times$ 3 \\
Other 8 PUTs (a1/b1/c1/c2/c3/d1/d2/d3) & 200 & 0 & --- \\
\end{longtable}
}

\textbf{Interpretation}:

\begin{itemize}
\item
  \textbf{HP/SI/TF class aggregate zero overlap rate} (0/72, 0/33,
  0/54): cosmic-ray AST-local operators \textbf{structurally
  unreachable} §3.2.0 necessary conditions (a)(c) --- 12-PUT full
  empirics confirm §3.2.6.1 categorical argument. HP (Hyper-Parameter)
  requires recognizing semantic dimensions of sklearn / scipy model
  object hyperparameters; SI (Structural Injection) requires mutating
  semantic intent of control flow; TF (Numerical Transform) requires
  changing semantic order of integration methods --- all three classes
  are \textbf{systematically unable to be generated} by cosmic-ray
  default operators' BinOp / Compare / NumberReplacer and other
  AST-local operators.
\item
  \textbf{CE class aggregate 7.81\% overlap rate} (5/64): CE is the
  ``boundary class'' marked in §3.2.0 --- only partially satisfies
  necessary condition (b) (carries domain knowledge),
  \textbf{conceptually aligned} with cosmic-ray NumberReplacer (see
  §3.2.6.1 row 1). The aggregate 7.81\% instance overlap mainly comes
  from simple ±1-style numerical replacements generated by LLMs (Claude
  / DeepSeek) on a2/b3 (e.g., \texttt{$\alpha$=2\ $\to$\ $\alpha$=1} in
  \texttt{Beta($\alpha$,\ $\beta$)}), whose AST representation overlaps with
  cosmic-ray NumberReplacer output. This is completely consistent with
  the prediction in §3.2.6.1 table marking CE as ``$\bigtriangleup$ literal values
  only'': \textbf{conceptual intersection exists, instance intersection
  is non-zero}. But 92.19\% of CE instances remain AST-disjoint (reason:
  LLMs tend to select domain-aware half-step perturbations, such as
  \texttt{\_RHO=28.0\ $\to$\ 27.5} preserving Lorenz chaotic regime, rather
  than tool integer addition and subtraction).
\item
  \textbf{OS class aggregate 11.67\% overlap rate} (7/60, \textbf{new
  finding}): The pre-existing argument (§3.2.6.1 row 2) claimed that the
  OS class is completely unreachable by cosmic-ray; the 12-PUT empirics
  show that this claim is too strong in the \emph{categorical} sense,
  and is in practice \textbf{88.33\% disjoint + 11.67\% incidental
  hits}. The hits concentrate on DeepSeek outputs for a3 and b3, all of
  which are low-syntactic-complexity OS sub-expressions (e.g.,
  \texttt{dx**2} $\to$ \texttt{dx*dx} in a3 FDM --- an algebraically
  equivalent rewrite incidentally hit by cosmic-ray BinOp, but labeled
  OS1 by the P2 LLM because the rewrite is semantically equivalent on
  the PUT). This is the empirical confirmation of the §3.2.6.0
  systematic-vs-incidental argument: \textbf{incidental hits by
  syntactic tools on the cross-function-boundary condition (a) are
  non-zero-probability events}. But the OS class's overall 88.33\%
  AST-disjointness still systematically rules out an ``OS = AST-local''
  classification.
\item
  \textbf{CF class aggregate 33.33\% overlap rate} (3/9): CF is
  instantiated as a sub-class of SI only in b2 (MCMC); all 9 mutants
  come from the b2 pool, of which 3 DeepSeek outputs hit cosmic-ray's
  \texttt{BreakContinueReplacer} / \texttt{ZeroIterationForLoop} at the
  AST level. The sample size is too small (n=9) to obtain a robust
  generalization; but even a 33.33\% rate means 66.67\% AST-disjoint,
  consistent with the categorical claim ``SI/CF unreachable'' in the
  §3.2.6.1 table being downgraded to \emph{partial reachability}.
\item
  \textbf{Overall overlap 5.14\% $\ll$ 1.0}: i.e., the P2 mutant pool
  \textbf{is not} a subset of the syntactic mutant pool; the two are
  systematically distinct mutant spaces. 94.86\% of P2 mutants are
  AST-disjoint with the 1250-mutant cosmic-ray output, and the absolute
  majority remain disjoint even on sub-classes (CE / OS partial) that
  are \emph{conceptually} reachable by syntactic tools.
\end{itemize}

\textbf{Conclusion (Layer 3 --- refuting the ``new-concept
classification'' challenge)}: Across the full 12-PUT empirics,
\textbf{94.86\%} of P2 mutants cannot be reproduced by cosmic-ray
default operators; the three classes HP / SI / TF (159/292 = 54.5\% of
P2 mutants) are \textbf{categorically} unrepresentable by syntactic
tools; the per-class overlap rates of CE / OS / CF --- 7.81\% / 11.67\%
/ 33.33\% --- are all far below 100\%. \textbf{Structural proof} that P2
is a systematic semantic mutation method, not a ``post-classification
copy'' of syntactic mutants --- even on sub-classes that are
conceptually reachable by syntactic tools, P2 mutants remain dominated
at the \emph{instance-selection} level by the LLM's domain-aware choices
(cf.~the source of the 92\% AST-disjointness in the CE class: LLMs
select domain-meaningful perturbations rather than ±1). This together
with the §3.2.6.0 systematic-vs-incidental argument and the §3.2.6.1
operator-level cross-table forms a complete refutation evidence chain;
meanwhile, the OS row's ``$\times$ not covered'' mark in the §3.2.6.1 table is
too absolute, and is empirically refined by this section to
\textbf{88.33\% disjoint + 11.67\% incidental hits}, with the nuance
honestly reflected in the OS interpretation of this section.

\textbf{Scope caveat}: A multi-syntactic-tool (mutmut / mutpy)
cross-comparison remains for P4 --- this section's empirics are based on
the single tool cosmic-ray, because mutpy is incompatible with Python
3.10+ (errors when testing sklearn / scipy), and mutmut's operator set
overlaps strongly with cosmic-ray's (cf.~§3.2.6.1). Effect of LLM source
differences (Claude / GPT / DeepSeek) on overlap patterns: hit instances
are not evenly distributed across the three LLMs (DeepSeek 11/15, Claude
4/15, GPT 0/15; see
\texttt{data/results/cosmic\_ray\_12put\_ast\_diff.json} for the
overlap\_files list), suggesting that DeepSeek tends to generate
syntactically simpler mutations. This LLM-source bias is discussed in
§7.2 (R8) ``LLM source distributional shift'' and does not affect the
systematic-vs-incidental argument of this section (because the
systematic argument is based on categorical AST-locality, not hit
frequency).

\subsubsection{3.3 60-cell instantiation matrix (three-state
decomposition)}\label{cell-instantiation-matrix-three-state-decomposition}

\textbf{Section positioning}: §3.2 subsections mut\_X provide abstract
definitions of the 5 meta mutation operators (CE / OS / HP / TF / SI,
i.e., mut\_C / mut\_M / mut\_G / mut\_T / mut\_F). This section provides
the \textbf{full specialization grid of meta operators across 12 PUTs $\times$
5 MPs} (60 cells)---each cell is a concrete instantiation of a (meta
operator, PUT type, meta pattern) triple. This section constitutes the
\textbf{concrete instantiation} of §3.2 meta operators and serves as the
reference object for the mutant traceability audit in §3.2.6.3.

\begin{verbatim}
              MP_1   MP_2   MP_3    MP_4    MP_5
              cons.  mono.  conv.   traj.   p-ord.
            +-----+-----+-----+-----+-----+
   A1 Lorenz| \textbullet\textbullet  | \textbullet   | \textbullet\textbullet  | \textbullet\textbullet  |  $\circ$  |
   A2 LU    | \textbullet\textbullet  |  $\circ$  | \textbullet   | \textbullet   | \textbullet\textbullet  |
   A3 FDM   | \textbullet\textbullet  | \textbullet   | \textbullet\textbullet  | \textbullet\textbullet  |  $\circ$  |
   B1 BetaBin|\textbullet\textbullet  | \textbullet   |  $\circ$  |  $\circ$  | \textbullet   |
   B2 MCMC  | \textbullet   | \textbullet\textbullet  | \textbullet\textbullet  | \textbullet\textbullet  | \textbullet   |
   B3 MC    | \textbullet\textbullet  |  $\circ$  | \textbullet\textbullet  | \textbullet   |  $\circ$  |
   C1 GPR   | \textbullet   | \textbullet\textbullet  | \textbullet\textbullet  | \textbullet   | \textbullet\textbullet  |
   C2 PCE   | \textbullet\textbullet  | \textbullet   | \textbullet\textbullet  | \textbullet   | \textbullet\textbullet  |
   C3 NN-Surr|\textbullet   | \textbullet\textbullet  | \textbullet\textbullet  | \textbullet\textbullet  | \textbullet\textbullet  |
   D1 MLP   | \textbullet\textbullet  | \textbullet\textbullet  | \textbullet   | \textbullet   | \textbullet\textbullet  |
   D2 SVM   | \textbullet   | \textbullet\textbullet  | \textbullet   |  $\circ$  | \textbullet\textbullet  |
   D3 LR    | \textbullet\textbullet  | \textbullet\textbullet  | \textbullet   |  $\circ$  | \textbullet\textbullet  |
            +-----+-----+-----+-----+-----+
   \textbullet\textbullet substantial 30 cells / \textbullet moderate 24 cells / $\circ$ vacant 6 cells (inherited from P1 H6)
\end{verbatim}

\subsubsection{3.4 Experimental scale
(measured)}\label{experimental-scale-measured}

\begin{itemize}
\tightlist
\item
  Each PUT averages 24.3 mutants (v4 cross-source pool; range 10-30, c1
  GPR constrained to 10 by V1-V4 pass rate)
\item
  All 60 cells total \textasciitilde24.3 $\times$ 12 = \textasciitilde292
  mutant instantiations (reusing the same PUT pool across the 12 PUT $\times$ 5
  MP matrix), coupled with N=20 AVP repeated sampling
\item
  v4 cache\_cross contains 298 V1-V4 confirmed mutants (three LLM
  sources $\times$ 37 operators $\times$ K=3 = 333 trials, 89\% pass rate, see
  §4.2.5(d))
\end{itemize}

\subsubsection{3.5 Engineering
significance}\label{engineering-significance}

\begin{itemize}
\tightlist
\item
  \textbf{Diagonal j=k (aligned)}: H2 threshold test
\item
  \textbf{Off-diagonal j$\neq$k (cross)}: control
\item
  \textbf{Vacant cells $\circ$}: not formally adjudicated in this paper
  (original H3 retired, see §1.5; vacant slices repurposed in §6.2 as
  descriptive evidence for R\_sem/R\_kill decoupling)
\end{itemize}

\paragraph{3.5.1 c-class primary MP selection: pre-registered v3 vs
exploratory
v3b}\label{c-class-primary-mp-selection-pre-registered-v3-vs-exploratory-v3b}

\textbf{Pre-registered primary analysis (v3, main analysis of this
paper)}: The primary MP for c-class (c1 / c2 / c3, surrogate) follows
the MP5 specification from P1 {[}Meng Li et al., Progress in Nuclear
Energy, under review{]}. All H1-H5 hypothesis tests render primary
verdicts on v3 data.

\textbf{Exploratory sensitivity analysis (v3b, post-hoc selection)}:
After observing v3 data, §5.8.4 Friedman per-class shows $\chi$$^{2}$ = 4.00, p =
0.406 for c-class (no statistically significant difference among MPs).
Based on this non-significant result, we conducted a \textbf{post-hoc
data-driven primary MP shift}, selecting the MP with maximum mean SMS
across the three c-class PUTs as the new primary:

{\def\LTcaptype{none} % do not increment counter
\begin{longtable}[]{@{}ll@{}}
\toprule\noalign{}
MP & mean SMS (c1, c2, c3 average) \\
\midrule\noalign{}
\endhead
\bottomrule\noalign{}
\endlastfoot
MP1 & 0.233 \\
MP2 & 0.000 \\
MP3 & 0.000 \\
MP4 & 0.000 \\
MP5 & 0.000 \\
\end{longtable}
}

New primary = MP1 (data in
\texttt{data/results/c\_class\_mp\_ranking.json}).

\textbf{Critical caveats (honest disclosure)}:

\begin{enumerate}
\def\labelenumi{\arabic{enumi}.}
\tightlist
\item
  \textbf{This adjustment is exploratory, not a pre-registered decision
  rule}. ``Friedman per-class p = 0.406 indicates any MP could serve as
  primary'' is a necessary but insufficient condition; the actual
  selection of argmax(mean\_SMS) constitutes selection-on-response,
  inflating H4 sign test pass rate (v3 3/4 $\to$ v3b 4/4) and $\delta$ (v3 0.323 $\to$
  v3b 0.446) while introducing selection bias
\item
  \textbf{No multiple-comparison correction applied}: selecting the
  maximum from 5 candidate MPs without applying strict max-statistic
  null distribution or Bonferroni $\times$ 5 correction
\end{enumerate}

\textbf{Quantitative bound on selection inflation} (added in P0-4
revision): permutation test quantifies the max-over-5 selection rule
post-hoc. \textbf{Null design} (corrected): treat the 15 c-class (PUT,
MP) cell SMS values as fully exchangeable; each permutation shuffles the
(PUT, MP) $\to$ SMS association and recalculates the max-over-5 per-PUT
mean, for N\_PERM = 10,000 iterations. \textbf{Result}: observed c-class
aligned mean = 0.3136, null distribution mean = 0.3494 ± 0.0347,
one-sided p(observed $\geq$ null) = 0.9885 (positioned at right-tail
percentile = (1 − 0.9885) $\times$ 100\% of null distribution).
\textbf{Bonferroni upper bound}: $\alpha$\_effective = $\alpha$ / 5 = 0.01
(family-of-5 candidates per PUT). \textbf{Joint interpretation}: both
permutation and Bonferroni sensitivity analyses indicate that the c$\to$MP1
selection effect size is substantially influenced by max-over-5
selection; abstract and §5.8.2 have accordingly downgraded all v3b
results to exploratory status (P0-3). See
\texttt{data/results/c\_class\_permutation\_v4.json} for details.

\begin{enumerate}
\def\labelenumi{\arabic{enumi}.}
\setcounter{enumi}{2}
\tightlist
\item
  \textbf{The reported v3b sign test 4/4 and $\delta$ improvement +0.123 should
  be interpreted as exploratory findings rather than confirmatory
  results}
\item
  \textbf{Pre-registered primary conclusion} (v3): H4 sign test 3/4
  (partially met), $\delta$ = 0.323 (H2 rejected). This is the primary analysis
  verdict of this paper; v3b/v4 are reported as sensitivity analyses
  only
\end{enumerate}

\textbf{Future work (P4)}: pre-register c-class primary MP selection
rule on new datasets (candidate rules: argmax of pre-defined statistic
or leave-one-class-out cross-validation) to eliminate the confound
present in this paper's v3b.

\subsubsection{3.6 Expected risk profile of LRCA across 60
cells}\label{expected-risk-profile-of-lrca-across-60-cells}

\paragraph{3.6.1 PUT-class $\times$ LRCA layer risk
weights}\label{put-class-lrca-layer-risk-weights}

{\def\LTcaptype{none} % do not increment counter
\begin{longtable}[]{@{}
  >{\raggedright\arraybackslash}p{(\linewidth - 10\tabcolsep) * \real{0.1667}}
  >{\raggedright\arraybackslash}p{(\linewidth - 10\tabcolsep) * \real{0.1667}}
  >{\raggedright\arraybackslash}p{(\linewidth - 10\tabcolsep) * \real{0.1667}}
  >{\raggedright\arraybackslash}p{(\linewidth - 10\tabcolsep) * \real{0.1667}}
  >{\raggedright\arraybackslash}p{(\linewidth - 10\tabcolsep) * \real{0.1667}}
  >{\raggedright\arraybackslash}p{(\linewidth - 10\tabcolsep) * \real{0.1667}}@{}}
\toprule\noalign{}
\begin{minipage}[b]{\linewidth}\raggedright
PUT class
\end{minipage} & \begin{minipage}[b]{\linewidth}\raggedright
L0 artifact
\end{minipage} & \begin{minipage}[b]{\linewidth}\raggedright
L1 tolerance (C2)
\end{minipage} & \begin{minipage}[b]{\linewidth}\raggedright
L2 OOD (C3)
\end{minipage} & \begin{minipage}[b]{\linewidth}\raggedright
L3 statistical (C4)
\end{minipage} & \begin{minipage}[b]{\linewidth}\raggedright
Primary risk root cause
\end{minipage} \\
\midrule\noalign{}
\endhead
\bottomrule\noalign{}
\endlastfoot
A numerical & $\star$ & \textbf{$\star$$\star$$\star$} & --- & --- & C2 dominant \\
B probabilistic & $\star$ & $\star$ & --- & \textbf{$\star$$\star$$\star$} & C4 dominant \\
C surrogate & $\star$ & $\star$$\star$ & \textbf{$\star$$\star$$\star$} & $\star$ & C3 dominant \\
D ML & $\star$ & $\star$ & \textbf{$\star$$\star$$\star$} & $\star$$\star$ & C3 + C4 mixed \\
\end{longtable}
}

\paragraph{3.6.2 Operator-PUT pair root cause hotspots
(expected)}\label{operator-put-pair-root-cause-hotspots-expected}

\begin{verbatim}
                A1 A2 A3 | B1 B2 B3 | C1 C2 C3 | D1 D2 D3
           -------------------------------------------------
   mut_C   | C2 C2 C2 | C4 C4 C2 | C2 C2 C5 | C5 C2 C2
   mut_M   | C1 C1 C1 | C4 C4 C1 | C1 C1 C5 | C5 C1 C1
   mut_G   | C2 C1 C2 | C4 C4 C2 | C3 C3 C5 | C5 C1 C1
   mut_T   | C1 C2 C1 | C4 C4 C4 | C3 C3 C5 | C5 C3 C3
   mut_F   | C2 C1 C1 | C4 C4 C4 | C3 C3 C5 | C5 C1 C3
\end{verbatim}

\paragraph{3.6.3 LRCA expected suspect\_share
thresholds}\label{lrca-expected-suspect_share-thresholds}

{\def\LTcaptype{none} % do not increment counter
\begin{longtable}[]{@{}
  >{\raggedright\arraybackslash}p{(\linewidth - 4\tabcolsep) * \real{0.3333}}
  >{\raggedright\arraybackslash}p{(\linewidth - 4\tabcolsep) * \real{0.3333}}
  >{\raggedright\arraybackslash}p{(\linewidth - 4\tabcolsep) * \real{0.3333}}@{}}
\toprule\noalign{}
\begin{minipage}[b]{\linewidth}\raggedright
PUT class
\end{minipage} & \begin{minipage}[b]{\linewidth}\raggedright
Expected suspect\_share
\end{minipage} & \begin{minipage}[b]{\linewidth}\raggedright
Acceptance threshold
\end{minipage} \\
\midrule\noalign{}
\endhead
\bottomrule\noalign{}
\endlastfoot
A numerical & 0.10--0.20 & $\leq$ 0.25 \\
B probabilistic & 0.20--0.35 & $\leq$ 0.40 \\
C surrogate & 0.20--0.30 & $\leq$ 0.35 \\
D ML & 0.25--0.40 & $\leq$ 0.45 \\
\textbf{All 60 cells average} & \textbf{$\leq$ 0.20} (corresponding to H5) &
\textbf{$\leq$ 0.25} \\
\end{longtable}
}

\paragraph{3.6.4 Interface with §4 LRCA execution
protocol}\label{interface-with-4-lrca-execution-protocol}

\begin{itemize}
\tightlist
\item
  L0 artifact pre-scan $\to$ §4.2.4 two-LLM double-blind review
\item
  L1 tolerance robustness $\to$ §4.6 N=20 repetition subprocess
\item
  L2 OOD triage $\to$ §4.6 input distribution definition
\item
  L3 statistical hypothesis baseline $\to$ §4.6 hypothesis pre-check
\end{itemize}

\subsubsection{3.7 Interface with §4 experimental
procedure}\label{interface-with-4-experimental-procedure}

§4 operationalization: (a) semantic mutation generation protocol
(dual-LLM cross-source + 20\% manual sampling, addressing LLM
reproducibility); (b) AVP invocation and version pinning; (c) LRCA
three-layer diagnostic execution; (d) §5 main table SMS + sub-table
C1\_share reporting pipeline.

\begin{center}\rule{0.5\linewidth}{0.5pt}\end{center}

\subsection{Section 4 $\cdot$ Experimental
Procedure}\label{section-4-experimental-procedure}

\subsubsection{4.1 Experimental Procedure Overview (Data
Flow)}\label{experimental-procedure-overview-data-flow}

\begin{verbatim}
[§3.1 PUT Selection] --► [§4.2 Semantic Mutation Generation]
                          |
                          ▼
                     [§4.3 Mutant Pool Pre-screening (LRCA L0)]
                          |
                          ▼
                     [§4.4 Equivalence Detection (E1 $\wedge$ E2)]
                          |
                          +--► equiv Set (Excluded)
                          |
                          ▼
                     [§4.5 AVP Invocation $\to$ killed/survive Classification]
                          |
                          ▼
                     [§4.6 LRCA Three-layer Diagnosis $\to$ root_cause Annotation]
                          |
                          ▼
                     [§4.7 §5 Main Table SMS + Subtable C1_share Report]
\end{verbatim}

\subsubsection{4.2 Semantic Mutation Generation Protocol (Addressing LLM
Reproducibility
Risk)}\label{semantic-mutation-generation-protocol-addressing-llm-reproducibility-risk}

\paragraph{4.2.1 Dual-source Generation}\label{dual-source-generation}

{\def\LTcaptype{none} % do not increment counter
\begin{longtable}[]{@{}
  >{\raggedright\arraybackslash}p{(\linewidth - 4\tabcolsep) * \real{0.3333}}
  >{\raggedright\arraybackslash}p{(\linewidth - 4\tabcolsep) * \real{0.3333}}
  >{\raggedright\arraybackslash}p{(\linewidth - 4\tabcolsep) * \real{0.3333}}@{}}
\toprule\noalign{}
\begin{minipage}[b]{\linewidth}\raggedright
Source
\end{minipage} & \begin{minipage}[b]{\linewidth}\raggedright
Proportion
\end{minipage} & \begin{minipage}[b]{\linewidth}\raggedright
Role
\end{minipage} \\
\midrule\noalign{}
\endhead
\bottomrule\noalign{}
\endlastfoot
Multi-LLM consensus & 60\% & Claude Opus 4.6 (via bltcy.ai) + GPT-5.4
(via bltcy.ai) + DeepSeek chat (via deepseek.com) three-model unanimous
agreement required for pool admission (actual v4 implementation in this
paper described in §4.2.5) \\
Manual injection & 40\% & 1-2 scientific computing researchers manually
write according to §3.2 rules \\
\end{longtable}
}

\begin{quote}
Note: The ``GPT-4o + Claude Opus + Gemini'' in §4.2.1 is the P1 protocol
description; the P2 implementation in this paper changed to Claude Opus
4.6 + GPT-5.4 + DeepSeek chat (see §4.2.5 cross-source protocol for
details). The only difference between the two is the specific LLM
selection; the prompt template and consensus voting mechanism remain
unchanged. The v4 cross-source data in this paper (298 confirmed
mutants) is entirely produced by the §4.2.5 protocol.
\end{quote}

\paragraph{4.2.2 Prompt Template Specification (LLM
Path)}\label{prompt-template-specification-llm-path}

Each (mut\_j, S\_i) pair is assigned 1 prompt template containing: -
Program under test (PUT) source code snippet - Semantic intent of mut\_j
- \textbf{Prohibition} against informing the LLM of any specific form of
MR (to avoid prompt leakage) - Output requirements: syntactically
correct + executable + single-point modification (diff \textless{} 10
lines)

\paragraph{4.2.3 Reproducibility
Parameters}\label{reproducibility-parameters}

\begin{itemize}
\tightlist
\item
  LLM temperature = 0.3
\item
  random\_seed fixed
\item
  Each prompt generates 5 candidates, review retains 2-3
\item
  Prompt template, LLM version, seed all stored in reproducibility
  package
\end{itemize}

\paragraph{4.2.4 Double-blind Review Protocol (Scheme C: Dual-LLM
Cross-source + 20\% Manual
Sampling)}\label{double-blind-review-protocol-scheme-c-dual-llm-cross-source-20-manual-sampling}

\textbf{(a)} Generator \textbf{LLM-G (Claude Opus)} and reviewer
\textbf{LLM-R (GPT-4o)} must be cross-source models with strictly
separated roles

\textbf{(b)} LLM-R review only sees PUT original code + mutant code,
unaware of mut\_j category / MR content / generator identity; outputs
triple: - Syntactically correct? \{Yes, No\} - Executable? \{Yes, No,
Uncertain\} - Semantic failure injected? \{Yes, No, Uncertain\}

\textbf{(c)} Double-confirmed mutants enter pool; inconsistencies or
Uncertain cases enter \textbf{manual arbitration queue} (estimated $\leq$
10\% of total)

\textbf{(d)} From the double-confirmed pool, randomly sample
\textbf{20\% for manual review by scientific computing researchers}; if
manual-LLM-R inconsistency rate \textgreater{} 10\%, trigger
\textbf{full manual downgrade}

\textbf{(e)} Reproducibility package public: LLM-G/R model versions,
prompt templates, temperature, seed, manual arbitration records

\textbf{(f)} Same-source bias mitigation: stratify LLM-R by PUT class
(A/B use GPT-4o, C/D use Claude Opus for cross-review), or three-LLM
rotation (GPT-4o / Claude Opus / Gemini Pro)

\textbf{(g)} Prompt injection leakage mitigation: LLM-G uses Application
Programming Interface (API) path with Retrieval-Augmented Generation
(RAG) / web search disabled; §7 Limitations explicitly acknowledges
residual risk

\paragraph{4.2.5 Cross-source Mutant Pool Protocol (Phase A
Methodological
Contribution)}\label{cross-source-mutant-pool-protocol-phase-a-methodological-contribution}

To isolate the relative contributions of ``LLM same-source bias'' and
``MR-MP alignment design'' to SMS effect size (Cliff's $\delta$), we designed a
three-stage ablation:

{\def\LTcaptype{none} % do not increment counter
\begin{longtable}[]{@{}
  >{\raggedright\arraybackslash}p{(\linewidth - 6\tabcolsep) * \real{0.2500}}
  >{\raggedright\arraybackslash}p{(\linewidth - 6\tabcolsep) * \real{0.2500}}
  >{\raggedright\arraybackslash}p{(\linewidth - 6\tabcolsep) * \real{0.2500}}
  >{\raggedright\arraybackslash}p{(\linewidth - 6\tabcolsep) * \real{0.2500}}@{}}
\toprule\noalign{}
\begin{minipage}[b]{\linewidth}\raggedright
Data Version
\end{minipage} & \begin{minipage}[b]{\linewidth}\raggedright
Mutant Pool Source
\end{minipage} & \begin{minipage}[b]{\linewidth}\raggedright
c-class Primary MP
\end{minipage} & \begin{minipage}[b]{\linewidth}\raggedright
Primary Use
\end{minipage} \\
\midrule\noalign{}
\endhead
\bottomrule\noalign{}
\endlastfoot
\textbf{v3} & Single-source (Claude Opus 4.6 only) & MP5 (P1 legacy) &
H2 baseline \\
\textbf{v3b} & Single-source (same as v3) & \textbf{MP1 (data-driven,
§3.5.1)} & Isolate contribution of MR-MP alignment design \\
\textbf{v4} & \textbf{Cross-source} (Claude Opus 4.6 + GPT-5.4 +
DeepSeek chat) & MP1 (same as v3b) & Isolate contribution of LLM source
diversity \\
\end{longtable}
}

\textbf{Protocol} (\texttt{scripts/cross\_source\_campaign.py}):

\begin{enumerate}
\def\labelenumi{(\alph{enumi})}
\item
  \textbf{Three-party LLM cross-source generation}: For each (PUT,
  operator) pair, Claude / GPT / DeepSeek each run K=3 trials with
  identical prompt template (§4.2.2), temperature fixed at 0.7,
  \textbf{source\_tag field propagated to mutant file naming}
  (\texttt{\{op\_id\}\_\{source\}\_attempt\{NN\}.py}) to facilitate
  source-stratified analysis
\item
  \textbf{Mechanical validation V1-V4}
  (\texttt{src/p2/mutators/validation.py}): syntax, executability,
  non-triviality
  (\texttt{\textbar{}y\_mutant\ -\ y\_original\textbar{}\ \textgreater{}\ 1e-6}
  on probe set), signature consistency.
\end{enumerate}

\textbf{Protocol asymmetry declaration (P1-7 revision, R0/R3/DA
consensus)}: The v4 cross-source pool \textbf{does not invoke reviewer
LLM} (cost/speed priority, deferred to P4 for full three-LLM dual-blind
review); whereas v3 / v3b data collection used the §4.2.4 original
Phase-1 dual-blind protocol (Claude-Opus generation + GPT-5.4 review +
DeepSeek arbitration). This constitutes \textbf{protocol asymmetry}
between v3 / v3b vs v4: the v4 mutant pool passes V1-V4 mechanical gates
but lacks the LLM review stage. \textbf{Potential confound}: A fraction
of $\Delta$$\delta$\_\{v3b $\to$ v4\} = −0.007 may not be ``LLM source diversity
contribution'' but rather ``slight quality decline in v4 pool relative
to v3b pool.'' This paper does not isolate this confound; P4 will rerun
dual-blind reviewer on the full v4 grid, enabling separation of quality
difference from source diversity factors. This confound is parallel to
§7.1 R11 chained-conditioning

\begin{enumerate}
\def\labelenumi{(\alph{enumi})}
\setcounter{enumi}{2}
\item
  \textbf{DeepSeek model selection}: \texttt{deepseek-chat} (not
  \texttt{deepseek-v4-pro}). V4-Pro is a reasoning model, averaging
  \textasciitilde340 tokens / call (\textasciitilde230 in
  reasoning\_content), latency 11s; the chat variant averages 113 tokens
  / call, latency 1.8s, \textbf{quality equivalent} (dry-run test shows
  three LLMs produce semantically identical sum-of-diagonal substitution
  on a2\_OS1 operator)
\item
  \textbf{Cross-source pool capacity}: 37 operators $\times$ 3 sources $\times$ K=3 =
  333 trials, V1-V4 pass rate 89\%, confirmed mutants 298; \textbf{three
  sources contribute nearly equally}: Claude=101, GPT=98, DeepSeek=99
  (Phase A key engineering finding: three LLMs have comparable
  capability on scientific computing mutant generation tasks)
\item
  \textbf{Pool sampling} (\texttt{scripts/build\_pools.py},
  POOL\_VERSION=v4): per-PUT 30 mutants upper limit, measured average
  24.3, range 10-30. c1 (GPR) only 10 mutants, because c1\_HP1 / c1\_CE1
  operators have V1-V4 pass rates near zero across all three LLMs (GPR
  PUT's WhiteKernel noise term 1e-4 $\to$ 1e-1 perturbation has minimal
  output impact, nearly all trigger V3 non-trivial failure, itself
  evidence for §6.2 R\_sem/R\_kill decoupling)
\end{enumerate}

\textbf{Chained-conditioning declaration (P0-5, R0/R1/DA consensus
revision)}: The v4 cross-source pool inherits v3b's post-hoc selection
on c-class primary MP (c1/c2/c3 $\to$ MP1). Therefore $\Delta$$\delta$\_\{v3b $\to$ v4\} =
−0.007 is not a neutral-condition LLM source diversity test, but rather
a contrast under the dual conditions of \emph{conditional on v3b's
c-class selection} + \emph{identical prompt template}. This paper does
not run the v4-pre (cross-source $\times$ c$\to$MP5 pre-shift) grid point, because
(i) relative to the current v3 $\to$ v3b $\to$ v4 chain, v4-pre primarily
answers ``whether v3b selection consumed part of the variance originally
attributable to LLM diversity,'' a question isomorphic to the P4 paper's
differential prompt experiment (§4.2.5.1), more economical to resolve
together in P4; (ii) v4-pre rerun cost \textasciitilde\$20-30 + 2-3 days
wall time, asymmetric to the narrative benefit already exposed in this
paper. \textbf{This declaration makes the conditional nature of the v3b
$\to$ v4 contrast explicitly visible, not dependent on reader inference}.

\textbf{Argumentation logic of the three-stage ablation}: The two
contrasts are \textbf{reported separately} (avoiding synthetic ratios
that imply factor isolation, see §5.7.2 table): v3 $\to$ v3b jump +0.123 is
a single-class, post-hoc c-class primary MP shift, reflecting
``sensitivity of primary MP selection rule'' rather than general ``MR-MP
alignment design contribution''; v3b $\to$ v4 micro-change −0.007 is the
near-unanimous response of three LLMs conditional on v3b selection +
identical prompt template. \textbf{The v3b $\to$ v4 contrast is not a
neutral-condition LLM source diversity test} (§4.2.5.1 R-16 protocol
lists the strong-sense version). This observation suggests P4
priorities: ``differential prompt frame per LLM'' (§7.1.7 R10) and
``pre-registered primary MP selection rule'' as parallel tracks; does
not constitute quantitative decomposition of ``MR design vs LLM
diversity'' two-factor contributions. Detailed interpretation in §5.7.2
+ §6.1.

\textbf{Contrast with existing work}: Tip et al.~(2024) LLMorpheus uses
single LLM (Claude) to generate JavaScript mutants, without cross-source
contrast; among known LLM-mutant empirical work, none isolate the
contribution of LLM source diversity. The three-stage ablation in §4.2.5
of this paper is the first work to factor-decompose ``LLM source
diversity vs MR design'' in the scientific computing software domain.

\subparagraph{4.2.5.1 Differential Prompt Protocol (R-16 Response,
Future-work
Commitment)}\label{differential-prompt-protocol-r-16-response-future-work-commitment}

\begin{quote}
Reviewer R-16 concern: v4 cross-source results show only −0.007
micro-change from v3b $\to$ v4, potentially reflecting ``same-source
response of three LLMs under identical prompt template'' rather than
``true upper bound of LLM diversity.'' To separate these two
possibilities requires \textbf{one tailored differential prompt per
LLM}, testing prompt-fixed vs prompt-differentiated contrast. This
section provides complete protocol, exit criteria, and resource
estimates for this experiment; execution deferred to R2 revision or P4
paper (consistent with §1.6 P-series sequencing), code skeleton ready
(\texttt{scripts/run\_differential\_prompt.py}).
\end{quote}

\textbf{Experimental design} (2 $\times$ 3 factorial, within-(PUT, operator)
repeated):

\begin{itemize}
\tightlist
\item
  Factor A: \textbf{prompt template}, 3 levels:

  \begin{itemize}
  \tightlist
  \item
    \textbf{V\_canonical} (control, identical to v4): generic mutation
    operator instruction (\texttt{prompts/operator\_template.txt})
  \item
    \textbf{V\_persona}: per-LLM identity guidance (Claude=``numerical
    analysis editor'', GPT=``scientific software refactorer'',
    DeepSeek=``library API substitution specialist''), remaining
    structure identical to V\_canonical
  \item
    \textbf{V\_cot}: per-LLM reasoning style (Claude=extended thinking
    with \texttt{\textless{}thinking\textgreater{}} tags; GPT=explicit
    ``Let's reason step by step'' framing; DeepSeek=stepped-reasoning
    with deliberate intermediate output)
  \end{itemize}
\item
  Factor B: \textbf{LLM source}, 3 levels (Claude / GPT-5.4 / DeepSeek
  chat, same as v4)
\item
  Replication: K=3 trials per (operator, prompt, source) cell
\item
  Total trials: 37 operators $\times$ 3 prompts $\times$ 3 sources $\times$ 3 trials =
  \textbf{999 trials}
\end{itemize}

\textbf{Statistical analysis}:

\begin{enumerate}
\def\labelenumi{\arabic{enumi}.}
\tightlist
\item
  Recalculate Cliff's $\delta$ under v4 protocol (should $\approx$ 0.439, as
  within-experiment baseline);
\item
  Build separate pools for each prompt variant (per-PUT pool size
  matched to v4), calculate $\delta$\_canonical / $\delta$\_persona / $\delta$\_cot;
\item
  \textbf{Core exit criteria}:

  \begin{itemize}
  \tightlist
  \item
    \textbf{If max($\delta$) − min($\delta$) \textless{} 0.05}: confirm ``three LLMs'
    differential response to prompt style is limited, weak LLM diversity
    contribution does not stem from prompt-fixed artifact,'' §5.7.2 v3b
    $\to$ v4 −0.007 conclusion robust;
  \item
    \textbf{If difference $\geq$ 0.05 and differential prompt pushes $\delta$ past
    0.474 large-effect threshold}: H2 verdict needs revision from ``not
    met'' to ``prompt-conditional met'';
    \texttt{paper\_numbers\_v4.json} marked deprecated, add
    \texttt{paper\_numbers\_v5.json} (after R-16 completion)
  \item
    \textbf{If difference $\geq$ 0.05 but $\delta$ still \textless{} 0.474}: confirm
    prompt sensitivity exists but does not change H2 verdict, §7.1.7 R10
    multi-prompt listed as limitation
  \end{itemize}
\end{enumerate}

\textbf{Resource estimates}: - API cost: Claude calls (input
\textasciitilde600 tokens, output \textasciitilde200 tokens) $\times$ 333 $\approx$ \$9
USD; GPT-5.4 same scale $\approx$ \$7; DeepSeek chat $\approx$ \$2; \textbf{total $\approx$
\$18-30 USD} (variable depending on V\_cot token inflation on reasoning
model) - Wall time: \textasciitilde3-4 hours at concurrency=4 (three API
rate-limits coexist) - Manual: V1-V4 automatic mechanical validation;
\textbf{no reviewer LLM invocation} (inherits §4.2.5(b) MVP
simplification); per-PUT pool resampling automatic

\textbf{Current paper version handling of R-16}: Without execution,
§5.7.2 already explicitly declares ``v3b $\to$ v4 micro-change −0.007 under
prompt-fixed'' as methodological limitation; §7.1.7 R10 lists
``differential prompt frame per LLM'' as P4 research direction. The
formalized record of R-16 protocol (this section) makes this commitment
\textbf{executable + verifiable}, reviewers can independently run
\texttt{scripts/run\_differential\_prompt.py} for verification.

\subsubsection{4.3 Mutant Pool Pre-screening (LRCA
L0)}\label{mutant-pool-pre-screening-lrca-l0}

Before each mutant enters the pool: - Static syntax checking (linters /
type checkers) - Unit self-test (runs on simple inputs and produces
finite output) - Double-blind review sign-off

\textbf{Pool size target}: 30--50 initial candidates per cell $\to$ 10--15
retained after pre-screening (corresponding to W-2).

\subsubsection{4.4 Equivalence Detection Procedure (E1 $\wedge$
E2)}\label{equivalence-detection-procedure-e1-e2}

\begin{quote}
§2.3 already provided E1 $\wedge$ E2 as the Layer 2 instantiation of the §3.2.0
necessary condition (argumentation for selection among three candidates
in §2.3). This section provides the engineering implementation procedure
for E1 $\wedge$ E2.
\end{quote}

\begin{verbatim}
For each mutant s' $\in$ mut_j(S_i):
  Step 1 [E2 Output-equiv]:
    Sample K_eq=1000 samples from D_S
    Compare $\|$S_i(x) − s'(x)$\|$ $\leq$ $\varepsilon$_eq
    If all satisfied $\to$ candidate equiv

  Step 2 [E1 AVP-coherent]:
    For each mr in MR_{i,k}, compare AVP(S_i, mr) with AVP(s', mr)
    If all consistent $\to$ confirm equiv

  Step 3 [Enter equiv_{i,k,j} set]:
    Exclude from denominator |mut_j(S_i)|
\end{verbatim}

\subsubsection{4.5 AVP Invocation (Addressing Circular
Dependency)}\label{avp-invocation-addressing-circular-dependency}

\paragraph{4.5.1 Version Pinning}\label{version-pinning-1}

AVP implementation version = P1 arXiv technical report commit hash; P2
reproducibility package embeds complete AVP source code.

\paragraph{4.5.2 Meta Pattern Verification
Execution}\label{meta-pattern-verification-execution}

\begin{verbatim}
For each (s', mr $\in$ MR_{i,k}):
  Invoke corresponding method based on MP_k to which mr belongs:
    MP_1:    tolerance equality determination
    MP_2/MP_5: Wilcoxon signed-rank, $\alpha$=0.05
    MP_3:    convergence order estimation + residual ratio
    MP_4:    DTW distance threshold $\varepsilon$_DTW
  Return {pass, fail}
\end{verbatim}

\paragraph{4.5.3 Killed Determination}\label{killed-determination}

\begin{verbatim}
killed(s', MR_{i,k}) ⇔
  $\exists$ mr $\in$ MR_{i,k}: AVP(S_i, mr) = pass $\wedge$ AVP(s', mr) = fail
\end{verbatim}

OR aggregation (classical).

\subsubsection{4.6 LRCA Three-layer
Diagnosis}\label{lrca-three-layer-diagnosis}

\paragraph{4.6.1 Decision Tree}\label{decision-tree-1}

\begin{verbatim}
For each s' $\in$ killed_{i,k,j}:
  step 1 - L1 tolerance robustness:
    Repeat N=20 times, fail proportion
      < 0.80 $\to$ root_cause = C2 $\times$
      $\geq$ 0.80 $\to$ proceed to L2

  step 2 - L2 OOD triage (C/D class only):
    Sample from D_S^valid, fail only in OOD region
      $\to$ root_cause = C3 $\times$
      Otherwise $\to$ proceed to L3

  step 3 - L3 statistical hypothesis baseline (B/D class + Wilcoxon/DTW only):
    Pre-check IID/stationarity on PUT's own repeated samples
      If hypothesis violated $\to$ root_cause = C4 $\times$
      Otherwise $\to$ proceed to step 4

  step 4 - artifact review:
    1 external reviewer re-examines mutant code and prompt history
      If LLM/artifact evidence $\to$ root_cause = C5 $\times$
      Otherwise $\to$ root_cause = C1 \checkmark
\end{verbatim}

\paragraph{4.6.2 Multi-label Handling}\label{multi-label-handling}

Priority: \textbf{C5 \textgreater{} C4 \textgreater{} C3 \textgreater{}
C2 \textgreater{} C1} (take earliest confirmed non-semantic cause).

\paragraph{4.6.3 Output}\label{output}

\begin{verbatim}
C1_share_{i,k,j}      := |{s' $\in$ killed : root_cause = C1}| / |killed|
suspect_share_{i,k,j} := 1 − C1_share_{i,k,j}
\end{verbatim}

LRCA \textbf{does not modify SMS formula}, killed set does not exclude
suspects.

\paragraph{4.6.4 LRCA Threshold Calibration (Methodological
Contribution)}\label{lrca-threshold-calibration-methodological-contribution}

LRCA L1-L3 introduces 3 engineering thresholds:

{\def\LTcaptype{none} % do not increment counter
\begin{longtable}[]{@{}
  >{\raggedright\arraybackslash}p{(\linewidth - 4\tabcolsep) * \real{0.3333}}
  >{\raggedright\arraybackslash}p{(\linewidth - 4\tabcolsep) * \real{0.3333}}
  >{\raggedright\arraybackslash}p{(\linewidth - 4\tabcolsep) * \real{0.3333}}@{}}
\toprule\noalign{}
\begin{minipage}[b]{\linewidth}\raggedright
Parameter
\end{minipage} & \begin{minipage}[b]{\linewidth}\raggedright
Meaning
\end{minipage} & \begin{minipage}[b]{\linewidth}\raggedright
Default
\end{minipage} \\
\midrule\noalign{}
\endhead
\bottomrule\noalign{}
\endlastfoot
\texttt{ood\_band} & L2 OOD boundary half-width: input x $\in$ {[}0,
ood\_band) $\cup$ (1-ood\_band, 1{]} considered OOD & 0.05 \\
\texttt{tolerance\_multiplier} & L1 epsilon\_loose / epsilon\_strict &
10.0 \\
\texttt{statistical\_repeats} & L3 majority-vote repetitions (changing
when N=20 requires caution) & 20 \\
\end{longtable}
}

To avoid threshold arbitrariness, we scan these parameters on a 9-grid
(\texttt{ood\_band\ $\in$\ \{0.02,\ 0.05,\ 0.10\}} $\times$
\texttt{tolerance\_multiplier\ $\in$\ \{3.0,\ 10.0,\ 30.0\}},
\texttt{repeats} fixed at 20), recording mean\_C1\_share and H5 pass
rate across 60 cells for each combination (data in
\texttt{data/results/lrca\_calibration.json}).

Main findings:

{\def\LTcaptype{none} % do not increment counter
\begin{longtable}[]{@{}llll@{}}
\toprule\noalign{}
ood\_band & tolerance\_multiplier & mean\_C1\_share & H5 cells\_pass \\
\midrule\noalign{}
\endhead
\bottomrule\noalign{}
\endlastfoot
\textbf{0.02} & \textbf{3.0} (best) & \textbf{0.200} & \textbf{12 / 60
(20.0\%)} \\
0.02 & 10.0 & 0.200 & 12 / 60 \\
0.02 & 30.0 & 0.200 & 12 / 60 \\
0.05 (default) & 3.0 / 10.0 / 30.0 & 0.164 & 10 / 60 (16.7\%) \\
0.10 & 3.0 / 10.0 / 30.0 & 0.164 & 10 / 60 (16.7\%) \\
\end{longtable}
}

Observations:

\begin{enumerate}
\def\labelenumi{\arabic{enumi}.}
\tightlist
\item
  \textbf{\texttt{tolerance\_multiplier} has zero impact on H5 across
  all 9 grid points}: each \texttt{ood\_band} yields identical c1\_share
  and h5\_pass across three multipliers. This indicates L1 tolerance
  determination is rarely triggered in this dataset; dominant signal
  comes from L2 OOD.
\item
  \textbf{\texttt{ood\_band} is the only discriminating factor}:
  tightening from 0.05 to 0.02 raises H5 from 10/60 to 12/60 (+2 cells).
  This reflects that LLM-generated mutant R-fail concentrates in
  narrower boundary regions; expanding OOD boundary mislabels more
  borderline kills as C3 rather than C1.
\item
  \textbf{Calibration ceiling}: even the best combination H5 = 12/60
  (20\%) remains far below the 80\% threshold set in §5.2. This
  indicates the LRCA signal-to-noise ratio of LLM-generated mutant pools
  is inherently low (empirical manifestation of §7.1.7 R10), H5 is
  unattainable on this dataset, not a threshold calibration issue.
\end{enumerate}

The H5 numbers reported in §5.6.2 use the best combination from this
section (ood\_band=0.02, tolerance\_multiplier=3.0, repeats=20); default
threshold results retained as control (\texttt{lrca\_60cell\_v3.json}).

\subsubsection{4.7 Data Collection and Reporting
Pipeline}\label{data-collection-and-reporting-pipeline}

\paragraph{4.7.1 Single-cell Output
Table}\label{single-cell-output-table}

After N=20 repetitions for each (i, k, j) cell, output:

{\def\LTcaptype{none} % do not increment counter
\begin{longtable}[]{@{}
  >{\raggedright\arraybackslash}p{(\linewidth - 2\tabcolsep) * \real{0.5000}}
  >{\raggedright\arraybackslash}p{(\linewidth - 2\tabcolsep) * \real{0.5000}}@{}}
\toprule\noalign{}
\begin{minipage}[b]{\linewidth}\raggedright
Quantity
\end{minipage} & \begin{minipage}[b]{\linewidth}\raggedright
Meaning
\end{minipage} \\
\midrule\noalign{}
\endhead
\bottomrule\noalign{}
\endlastfoot
inst\_count & Total mut\_j(S\_i) entering pool \\
equiv\_count & E1 $\wedge$ E2 pass count \\
killed\_count & AVP killed count \\
survive\_count & live − killed \\
\textbf{SMS\_\{i,k,j\}} & killed/(inst − equiv) \\
C1\_share\_\{i,k,j\} & LRCA-annotated C1 proportion \\
suspect\_share\_\{i,k,j\} & 1 − C1\_share \\
inst\_rate / equiv\_rate / survive\_rate & Three rates (corresponding to
RQ1) \\
\end{longtable}
}

\paragraph{4.7.2 Experimental Timeline}\label{experimental-timeline}

{\def\LTcaptype{none} % do not increment counter
\begin{longtable}[]{@{}
  >{\raggedright\arraybackslash}p{(\linewidth - 4\tabcolsep) * \real{0.3333}}
  >{\raggedright\arraybackslash}p{(\linewidth - 4\tabcolsep) * \real{0.3333}}
  >{\raggedright\arraybackslash}p{(\linewidth - 4\tabcolsep) * \real{0.3333}}@{}}
\toprule\noalign{}
\begin{minipage}[b]{\linewidth}\raggedright
Phase
\end{minipage} & \begin{minipage}[b]{\linewidth}\raggedright
Time
\end{minipage} & \begin{minipage}[b]{\linewidth}\raggedright
Task
\end{minipage} \\
\midrule\noalign{}
\endhead
\bottomrule\noalign{}
\endlastfoot
Phase 1 & 2026 Q4 & mut\_j generation + L0 pre-screening (60 cells $\times$
30-50 candidates) \\
Phase 2 & 2027 Q1 & Equivalence detection + AVP invocation (N=20
repetitions) \\
Phase 3 & 2027 Q2 & LRCA three-layer diagnosis + data archiving \\
Phase 4 & 2027 Q3 & §5 statistical analysis + paper writing + IST
submission \\
\end{longtable}
}

\subsubsection{4.8 Pilot Calibration (37 Operators K=10/20 Trial
Run)}\label{pilot-calibration-37-operators-k1020-trial-run}

\paragraph{4.8.1 Pilot Design}\label{pilot-design}

To calibrate end-to-end executability of §4.2 generation protocol, §4.3
L0 pre-screening, and §4.5 AVP invocation, a round of
\textbf{operator-level pilot} was conducted in 2026 Q3 on 12 PUTs (37
named operators, 12 operators with is\_key=True at K=20, remaining 25 at
K=10, total 470 trials). Generator used Claude Opus (subscription
authentication, subprocess invocation), reviewer used GPT-5.4 (bltcy.ai
proxy), review prompt isomorphic to §4.2.4 double-blind protocol. Pilot
output focused on three \textbf{operator-level precursor quantities}:

\begin{itemize}
\tightlist
\item
  \textbf{R\_sem}: proportion of K attempts per operator passing V1-V6 $\wedge$
  operator\_match=Yes (semantic implementation success rate,
  corresponding to §4.2.4 double-blind review final judgment)
\item
  \textbf{D\_impl}: median of K(K-1)/2 pairwise AST tuple + literal +
  identifier multiset Jaccard distances among confirmed mutants per
  operator (implementation diversity, characterizing ``different
  implementation surfaces of the same operator'')
\item
  \textbf{R\_kill}: proportion of confirmed mutants per operator killed
  by AVP on that PUT's primary meta pattern (MP) MR (operator-MR
  alignment fault detection rate)
\end{itemize}

D\_impl and R\_kill are operator-level pre-projections of the two
components of §2.4 SMS: D\_impl is the ``implementation dimension''
extension quantity from §2.4.2, R\_kill equals the SMS upper bound for
that cell under operator-MR alignment slicing (equality when equiv=0).

\paragraph{4.8.2 Pilot Main Results (Operator-level Aggregation,
N=37)}\label{pilot-main-results-operator-level-aggregation-n37}

{\def\LTcaptype{none} % do not increment counter
\begin{longtable}[]{@{}lllll@{}}
\toprule\noalign{}
Metric & Median & Mean & R\_sem=0 Operators & D\_impl $\approx$ 0 Operators \\
\midrule\noalign{}
\endhead
\bottomrule\noalign{}
\endlastfoot
R\_sem & 0.50 & 0.468 & 0 / 37 & --- \\
D\_impl & 0.42 & 0.392 & --- & 1 / 37 \\
R\_kill & 0.00 & 0.189 & --- & --- \\
\end{longtable}
}

All 37 operators produced at least one V1-V6 passing mutant (R\_sem
\textgreater{} 0), verifying \textbf{operator implementability} of §4.2
generation protocol (corresponding to H1). D\_impl median 0.42 indicates
K attempts genuinely disperse at structural level---this is the stable
level after multiple rounds of deduplication (each attempt's prompt
injected with ``prior candidates $\cdot$ DO NOT REPEAT'') combined with V5
single-fault semantic relaxation (allowing naming/formal/library
equivalent substitutions), a necessary precondition for subsequent §5.1
SMS calculation.

\paragraph{4.8.3 Key Finding: R\_sem and R\_kill Decoupling on HP-class
Operators}\label{key-finding-r_sem-and-r_kill-decoupling-on-hp-class-operators}

Pilot data reveals a pattern \textbf{spanning all 4 PUT classes}:
hyperparameter (HP) and constructive parameter (CE/OS) class operators
approach R\_sem $\approx$ 1.0, but R\_kill under that PUT's \textbf{primary meta
pattern (MP)} MR may still be 0. Table 4.8 excerpts representative 12
operators (distributed by class).

\textbf{Table 4.8 R\_sem / R\_kill Decoupling Excerpt (12 Operators
Selected)}

{\def\LTcaptype{none} % do not increment counter
\begin{longtable}[]{@{}
  >{\raggedright\arraybackslash}p{(\linewidth - 10\tabcolsep) * \real{0.1667}}
  >{\raggedright\arraybackslash}p{(\linewidth - 10\tabcolsep) * \real{0.1667}}
  >{\raggedright\arraybackslash}p{(\linewidth - 10\tabcolsep) * \real{0.1667}}
  >{\raggedright\arraybackslash}p{(\linewidth - 10\tabcolsep) * \real{0.1667}}
  >{\raggedright\arraybackslash}p{(\linewidth - 10\tabcolsep) * \real{0.1667}}
  >{\raggedright\arraybackslash}p{(\linewidth - 10\tabcolsep) * \real{0.1667}}@{}}
\toprule\noalign{}
\begin{minipage}[b]{\linewidth}\raggedright
Operator
\end{minipage} & \begin{minipage}[b]{\linewidth}\raggedright
Class
\end{minipage} & \begin{minipage}[b]{\linewidth}\raggedright
PUT Primary MP
\end{minipage} & \begin{minipage}[b]{\linewidth}\raggedright
R\_sem
\end{minipage} & \begin{minipage}[b]{\linewidth}\raggedright
D\_impl
\end{minipage} & \begin{minipage}[b]{\linewidth}\raggedright
R\_kill
\end{minipage} \\
\midrule\noalign{}
\endhead
\bottomrule\noalign{}
\endlastfoot
a2\_CE1 (LU determinant parameter) & CE & MP1 conservation & 0.60 &
0.601 & \textbf{1.00} \\
a2\_OS1 (prod$\to$sum) & OS & MP1 conservation & 0.25 & 0.493 &
\textbf{1.00} \\
b1\_OS1 ($\alpha$/$\beta$ swap) & OS & MP2 monotonicity & 0.30 & 0.456 &
\textbf{1.00} \\
d1\_TF1 (label flip) & TF & MP2 monotonicity & 0.50 & 0.355 &
\textbf{1.00} \\
\textbf{c1\_CE1 (noise 1e-4$\to$1e-1)} & CE & MP5 asymptotic & 1.00 & 0.512
& \textbf{0.00} \\
\textbf{c2\_OS1 (basis function poly$\to$spline)} & OS & MP5 asymptotic &
0.90 & 0.422 & \textbf{0.00} \\
\textbf{c3\_HP1 (relu$\to$tanh)} & HP & MP5 asymptotic & 1.00 & 0.460 &
\textbf{0.00} \\
\textbf{c3\_TF1 (max\_iter 1000$\to$5)} & TF & MP5 asymptotic & 1.00 & 0.464
& \textbf{0.00} \\
\textbf{d1\_HP1 (MLP $\alpha$ 1e-4$\to$1.0)} & HP & MP2 monotonicity & 1.00 & 0.505
& \textbf{0.00} \\
\textbf{d3\_HP1 (LR C 1.0$\to$1e-4)} & HP & MP2 monotonicity & 1.00 & 0.478
& \textbf{0.00} \\
d2\_HP1 ($\gamma$ `scale'$\to$1e-3) & HP & MP2 monotonicity & 0.50 & 0.541 &
0.00 \\
b2\_HP1 (proposal step size $\times$0.1) & HP & MP2 monotonicity & 0.25 & 0.366
& 0.00 \\
\end{longtable}
}

\textbf{Observation}: R\_sem $\geq$ 0.9 $\wedge$ R\_kill = 0 combinations
\textbf{all} appear on HP/TF/OS class operators in c/d classes
(surrogate and classifier, 6 of 12 PUTs); R\_sem high $\wedge$ R\_kill = 1
combinations \textbf{all} appear on CE/OS class operators in a/b classes
(numerical and statistical computing).

\paragraph{4.8.4 Pilot Pre-assessment of
RQ2}\label{pilot-pre-assessment-of-rq2}

§1.4 RQ2 asks about SMS difference structure on operator-meta pattern
(MP) \textbf{aligned (j=k) vs non-aligned (j$\neq$k) slices}. Although pilot
operator-level data does not yet span the complete 60-cell matrix (Phase
2 task), it already observes bimodal distribution of R\_kill on
\textbf{operator-primary MP aligned slices} (j being that PUT's primary
meta pattern (MP)):

\begin{itemize}
\tightlist
\item
  \textbf{a/b classes (numerical/statistical)}: CE/OS/SI/CF operator
  perturbations directly violate algorithmic algebraic invariants; even
  when R\_sem is not high, once mutant implementation succeeds, primary
  MP (MP1 conservation / MP2 monotonicity) can detect $\to$ R\_kill $\in$ \{0.5,
  1.0\}
\item
  \textbf{c/d classes (surrogate/classifier)}: HP/TF operator
  perturbations alter model capacity or training dynamics, producing
  \textbf{functionally correct but parameter-deviated} mutants; primary
  MP (MP2 monotonicity, MP5 asymptotic) are statistical / asymptotic
  MRs, insensitive to parameter deviation $\to$ R\_kill ≡ 0
\end{itemize}

This provides \textbf{pre-suggestive evidence} for H2 ``aligned-SMS /
cross-SMS odds ratio $\geq$ 3.0'': even within the j=k aligned slice,
\textbf{operator fault dimension} (algebraic vs.~statistical) and
\textbf{MR detection dimension} exhibit non-trivial second-order
mismatch. In other words, §5.1 RQ2 main analysis quantity should, beyond
``diagonal vs off-diagonal,'' also report in appendix
\textbf{second-level R\_kill conditional distribution by operator class
(CE/OS vs HP/TF) within diagonal slice}, to interface with this pilot's
decoupling observation. This will be added as a subfigure to §5.5 Figure

\subsection{Section 5 $\cdot$ Statistical Analysis
Methods}\label{section-5-statistical-analysis-methods}

\subsubsection{5.1 Primary Table (Corresponding to
RQ1-RQ4)}\label{primary-table-corresponding-to-rq1-rq4}

{\def\LTcaptype{none} % do not increment counter
\begin{longtable}[]{@{}
  >{\raggedright\arraybackslash}p{(\linewidth - 4\tabcolsep) * \real{0.3333}}
  >{\raggedright\arraybackslash}p{(\linewidth - 4\tabcolsep) * \real{0.3333}}
  >{\raggedright\arraybackslash}p{(\linewidth - 4\tabcolsep) * \real{0.3333}}@{}}
\toprule\noalign{}
\begin{minipage}[b]{\linewidth}\raggedright
RQ
\end{minipage} & \begin{minipage}[b]{\linewidth}\raggedright
Primary Statistics
\end{minipage} & \begin{minipage}[b]{\linewidth}\raggedright
Reporting Format
\end{minipage} \\
\midrule\noalign{}
\endhead
\bottomrule\noalign{}
\endlastfoot
RQ1 & inst\_rate, equiv\_rate, C1\_share, survive\_rate & 60-cell
heatmap + 4-class marginal distributions \\
RQ2 & aligned-SMS vs cross-SMS; sparse $\circ$ vs dense \textbullet\textbullet equiv\_rate &
Cliff's $\delta$ + odds ratio + 95\% bootstrap confidence interval \\
RQ3 & $\Delta$SMS\_c (c $\in$ \{A,B,C,D\}); CV($\Delta$SMS) & sign test (df=3) +
descriptive forest plot \\
RQ4 & Spearman $\rho$ + Kendall $\tau$ for SMS vs pattern coverage & scatter plot
+ dual-metric ranking comparison table \\
\end{longtable}
}

\subsubsection{5.2 Hypothesis Testing Thresholds
(H1-H5)}\label{hypothesis-testing-thresholds-h1-h5}

\begin{itemize}
\item
  \textbf{H1}: Number of cells with \textbar non-equivalent
  mutants\textbar{} $\geq$ 5 across 5 operators $\times$ 12 PUTs $\geq$ 9 $\times$ 4 = 36
  (simplified to ``$\geq$ 4 operators meet threshold'' for full 5-operator
  condition)
\item
  \textbf{H2}: aligned-SMS / cross-SMS odds ratio $\geq$ 3.0, Cliff's $\delta$ $\geq$
  0.474

  Cliff's $\delta$ threshold follows Romano et al.~(2006) empirical table for
  software engineering:

  {\def\LTcaptype{none} % do not increment counter
  \begin{longtable}[]{@{}lll@{}}
  \toprule\noalign{}
  Level & Threshold & Interpretation \\
  \midrule\noalign{}
  \endhead
  \bottomrule\noalign{}
  \endlastfoot
  Negligible & \textbar $\delta$\textbar{} \textless{} 0.147 & Two groups
  equivalent \\
  Small & 0.147 $\leq$ \textbar $\delta$\textbar{} \textless{} 0.330 & Weak
  directionality \\
  Medium & 0.330 $\leq$ \textbar $\delta$\textbar{} \textless{} 0.474 & Clear
  directionality \\
  Large & \textbar $\delta$\textbar{} $\geq$ 0.474 & Strong dominance \\
  \end{longtable}
  }

  The H2 large-effect threshold is a pre-commitment (following P1
  {[}Meng Li et al., Progress in Nuclear Energy, under review{]}
  setting), not modified post hoc in this paper. §5.7.2 / §6.1 report
  which tier the observed $\delta$ falls into and its implications for the
  argument.
\item
  \st{\textbf{H3}: equiv\_rate $\geq$ 0.85 for 6 $\circ$ cells; equiv\_rate $\leq$ 0.30
  for 30 \textbullet\textbullet cells (bidirectional threshold)} \textbf{(retired, see §1.5
  note)}
\item
  \textbf{H4}: sign test all four classes positive (p=0.0625); CV($\Delta$SMS)
  \textless{} 0.5
\item
  \textbf{H5}: average suspect\_share $\leq$ 0.20 across all 60 cells
\end{itemize}

\subsubsection{5.3 Multiple Comparisons and
Power}\label{multiple-comparisons-and-power}

\paragraph{5.3.1 Multiple Comparisons}\label{multiple-comparisons}

Cell-level claims across 60 cells use Benjamini-Hochberg FDR correction,
$\alpha$\_FDR = 0.05.

\paragraph{5.3.2 Statistical Power
Statement}\label{statistical-power-statement}

\begin{itemize}
\tightlist
\item
  \textbf{Cross-class consistency H4}: 4-class sign test has only df=3,
  weak power; P2 frames H4 as \textbf{exploratory evidence},
  supplemented by a mixed-effects model (random intercept: PUT; fixed
  effects: class $\times$ operator) over 12 PUTs $\times$ 5 operators
\item
  \textbf{N=20 bootstrap intervals}: 1000-iteration bootstrap 95\% CI
  provided for Cliff's $\delta$ and odds ratio; no parametric p-value
\end{itemize}

\subsubsection{5.4 LRCA Descriptive
Analysis}\label{lrca-descriptive-analysis}

\paragraph{5.4.1 Primary Diagnostic
Metrics}\label{primary-diagnostic-metrics}

\begin{itemize}
\tightlist
\item
  C1\_share heatmap across 60 cells (primary result robustness
  diagnostic)
\item
  suspect\_share marginal distribution for 4 PUT classes (compared
  against §3.6.3 expected threshold)
\end{itemize}

\paragraph{5.4.2 SMS\_unfiltered Appendix (Reviewer Comparison
Metric)}\label{sms_unfiltered-appendix-reviewer-comparison-metric}

\begin{verbatim}
SMS_unfiltered_{i,k,j} := |killed_{i,k,j}| / (|mut_j(S_i)| − |equiv_{i,k,j}|)
                       (same form as primary formula, killed does not distinguish C1/C2-C5)
\end{verbatim}

Appendix provides cell-by-cell difference table between SMS\_unfiltered
and SMS; if relative difference \textless{} 5\%, confirms LRCA does not
affect robustness of primary conclusions.

\subsubsection{5.5 Visualization}\label{visualization}

\begin{itemize}
\tightlist
\item
  \textbf{Figure 1}: 60-cell SMS heatmap (rows: PUT $\times$ columns: MP $\times$
  sub-cells: mut\_j)
\item
  \textbf{Figure 2}: diagonal vs off-diagonal slice boxplots (H2
  intuition)
\item
  \textbf{Figure 3}: 4-class $\Delta$SMS forest plot (H4 intuition)
\item
  \textbf{Figure 4}: SMS vs C1\_share scatter + Spearman $\rho$ (LRCA
  diagnostic)
\item
  \textbf{Figure 5}: SMS vs pattern coverage dual-metric ranking
  comparison (RQ4)
\end{itemize}

\begin{center}\rule{0.5\linewidth}{0.5pt}\end{center}

\subsubsection{5.6 RQ1 Empirical Results (60 cells, Track-2
v2)}\label{rq1-empirical-results-60-cells-track-2-v2}

\paragraph{5.6.1 Data Scale and Cell-Level SMS
Distribution}\label{data-scale-and-cell-level-sms-distribution}

Each PUT has 12 mutants (operator-cache proportional sampling, §4.2.5
builder), each (mutant, MR) pair computes R\_kill under N=20 AVP
repeated sampling (§4.5.3). Full 60-cell SMS table shown in Figure 1,
all numbers summarized in \texttt{data/results/paper\_numbers.json}.

Primary statistics:

{\def\LTcaptype{none} % do not increment counter
\begin{longtable}[]{@{}
  >{\raggedright\arraybackslash}p{(\linewidth - 2\tabcolsep) * \real{0.5000}}
  >{\raggedright\arraybackslash}p{(\linewidth - 2\tabcolsep) * \real{0.5000}}@{}}
\toprule\noalign{}
\begin{minipage}[b]{\linewidth}\raggedright
Metric
\end{minipage} & \begin{minipage}[b]{\linewidth}\raggedright
Value
\end{minipage} \\
\midrule\noalign{}
\endhead
\bottomrule\noalign{}
\endlastfoot
Number of cells & 60 \\
Mean SMS & 0.104 \\
Median SMS & 0.000 \\
Standard deviation SMS & 0.213 \\
Number of cells with SMS = 0 & 45 / 60 \\
Mean mutant count / cell & 24.3 (v4 cross-source pool, 10-30 per PUT) \\
\end{longtable}
}

\begin{quote}
Note: Among the 45 cells with SMS = 0, the vast majority concentrate in
the cross-MP slice (j $\neq$ k), i.e., the 4 MPs outside each PUT's primary
aligned MP. This is the direct data-level manifestation of H1 (MR-MP
alignment); §5.7 quantifies the aligned vs non-aligned slice gap using
Cliff's $\delta$.
\end{quote}

\paragraph{5.6.1.1 Zero-Mass Dominance (Distribution Characteristics of
75\% Zero
Cells)}\label{zero-mass-dominance-distribution-characteristics-of-75-zero-cells}

45 / 60 cells = 75\% of cells have SMS = 0, a highly skewed,
zero-concentrated distribution. Structural attribution:

{\def\LTcaptype{none} % do not increment counter
\begin{longtable}[]{@{}lll@{}}
\toprule\noalign{}
Slice & Zero cells / total & Zero proportion \\
\midrule\noalign{}
\endhead
\bottomrule\noalign{}
\endlastfoot
aligned (j == k) & \textasciitilde3 / 12 & 25\% \\
cross (j $\neq$ k) & \textasciitilde42 / 48 & 88\% \\
All 60 cells & 45 / 60 & 75\% \\
\end{longtable}
}

\textbf{Impact on statistical inference}:

\begin{itemize}
\tightlist
\item
  Cliff's $\delta$ is a rank-based measure, mathematically well-defined for
  zero-concentrated distributions, but actual signal primarily comes
  from the non-zero portion of the 12 aligned cells (median(cross)=0,
  $\delta$'s numerator \texttt{\#\{a\ \textgreater{}\ b\}} saturates at cross $\approx$
  0). \textbf{RQ2's effect-size inference is essentially dominated by
  n\_aligned = 12 (not 60)}
\item
  Median odds ratio is formally undefined because median(cross) = 0;
  this paper uses ``aligned median \textgreater{} 0 = cross median'' as
  auxiliary qualitative evidence
\item
  Friedman $\chi$$^{2}$ remains valid under zero-concentrated distributions (does
  not assume continuity, uses only ranks), but rank ties attenuate the
  statistic's power
\end{itemize}

\textbf{Consistency with LLM-mutant literature}: Tip et al.~(2024)
LLMorpheus similarly reports high proportions of cross-MP failure on JS
sci-comp PUTs (specific numbers not listed in cited literature); this
zero-mass dominance appears to be a shared characteristic of LLM-mutant
+ existing MR-MP alignment designs, not a special issue with this
paper's PUT selection.

\textbf{Future work}: Whether expanding the mutant pool to 30+ / PUT
(P4) can make more cross-MP cells exhibit non-zero SMS is a meaningful
power test.

{[}Figure 1: 60-cell SMS heatmap (rows = PUT, cols = MP, $\star$ marks j = k
aligned cells){]}

\paragraph{5.6.2 LRCA C1\_share / suspect\_share
Distribution}\label{lrca-c1_share-suspect_share-distribution}

LRCA three-tier diagnostics label each killed mutant as C1/C2/C3/C4
(§4.6). Average statistics across 60 cells:

{\def\LTcaptype{none} % do not increment counter
\begin{longtable}[]{@{}
  >{\raggedright\arraybackslash}p{(\linewidth - 2\tabcolsep) * \real{0.5000}}
  >{\raggedright\arraybackslash}p{(\linewidth - 2\tabcolsep) * \real{0.5000}}@{}}
\toprule\noalign{}
\begin{minipage}[b]{\linewidth}\raggedright
Metric
\end{minipage} & \begin{minipage}[b]{\linewidth}\raggedright
Value
\end{minipage} \\
\midrule\noalign{}
\endhead
\bottomrule\noalign{}
\endlastfoot
Mean C1\_share (legit fault proportion, default threshold) & 0.164 \\
Mean C1\_share (calibrated best ood\_band=0.02) & 0.200 \\
Mean suspect\_share (calibrated best) & 0.800 \\
Cells meeting H5 threshold (suspect\_share $\leq$ 0.20) (default) & 10 / 60
(16.7\%) \\
Cells meeting H5 threshold (calibrated best) & \textbf{12 / 60
(20.0\%)} \\
\end{longtable}
}

H5 verdict: \textbf{not met} (calibrated best combination yields 20.0\%,
still far below the 80\% threshold set in §5.2).

\begin{quote}
Interpretation: 9-grid LRCA threshold scan (§4.6.4) shows that
tightening the OOD boundary from 0.05 to 0.02 can improve H5 from 10/60
to 12/60, but still far from crossing the 80\% strict threshold;
tolerance multiplier has no impact on this dataset. This reflects: (a)
R-fail signals from LLM-generated mutants concentrate in narrow
boundaries, and (b) most killed mutants fall into C2/C3/C4 rather than
C1 under LRCA three-tier thresholds, representing an inherent ceiling of
LLM-homogeneous mutant pools (§7.1.7 R10), not an LRCA threshold issue.
We therefore designate H5 as \textbf{not met, but as an empirical
starting point for LRCA calibration research directions} (§6.2
engineering implications).
\end{quote}

\subparagraph{5.6.2.1 H5 Cutoff Sensitivity (R-14
Response)}\label{h5-cutoff-sensitivity-r-14-response}

\begin{quote}
Reviewer R-14 questions: Is the 0.20 in H5 pass condition
\texttt{suspect\_share\ $\leq$\ 0.20} a ``lucky pick''---would the verdict
change if the cutoff were changed to 0.15 or 0.30? This section uses a
dense grid (cutoff $\in$ \{0.05, 0.10, \ldots, 0.50\}, step 0.05) to scan v4
data (\texttt{scripts/h5\_sensitivity.py}).
\end{quote}

\textbf{Results} (\texttt{data/results/h5\_sensitivity\_v4.json}):

{\def\LTcaptype{none} % do not increment counter
\begin{longtable}[]{@{}lll@{}}
\toprule\noalign{}
cutoff & h5\_cells\_pass & h5\_pass\_ratio \\
\midrule\noalign{}
\endhead
\bottomrule\noalign{}
\endlastfoot
0.05 & 12 / 60 & 20.0\% \\
0.10 & 12 / 60 & 20.0\% \\
0.15 & 12 / 60 & 20.0\% \\
\textbf{0.20} (paper) & \textbf{12 / 60} & \textbf{20.0\%} \\
0.25 & 12 / 60 & 20.0\% \\
0.30 & 12 / 60 & 20.0\% \\
0.35 & 12 / 60 & 20.0\% \\
0.40 & 12 / 60 & 20.0\% \\
0.45 & 13 / 60 & 21.7\% \\
0.50 & 13 / 60 & 21.7\% \\
\end{longtable}
}

\textbf{Key observation}: H5 pass-ratio is completely flat at 20.0\% for
cutoff $\in$ {[}0.05, 0.40{]}, with only one cell crossing over at 0.45.
\textbf{No cutoff can push h5\_pass\_ratio above
80\%}---\texttt{smallest\_cutoff\_for\_80pct\_pass\ =\ None}. This is
because v4's suspect\_share distribution is severely bimodal: median =
1.0 (48 cross cells have suspect\_share $\approx$ 1 almost universally), mean =
0.79; only the 12 aligned cells have suspect\_share in the low end (near
0). The middle region {[}0.20, 0.80{]} is nearly empty.

\textbf{Conclusion}: H5 verdict \textbf{not met is an intrinsic data
property, independent of cutoff choice}. The specific value 0.20 is
\textbf{not load-bearing} in the ``lucky pick'' sense R-14 is concerned
about---changing it to 0.10 or 0.40 leaves the verdict completely
unchanged. R-14's specific concern (whether 0.20 choice is fragile)
receives a positive answer; extended LRCA threshold 49-grid scan
(\texttt{scripts/calibrate\_lrca.py\ LRCA\_GRID=49}) is reserved for P4
paper with larger pool scale for deeper calibration.

{[}Figure 4: SMS vs C1\_share scatter (per cell, n=60){]}

\begin{center}\rule{0.5\linewidth}{0.5pt}\end{center}

\subsubsection{5.7 RQ2 Empirical Results (Aligned vs
Non-Aligned)}\label{rq2-empirical-results-aligned-vs-non-aligned}

\paragraph{5.7.1 Descriptive Statistics}\label{descriptive-statistics}

Partition 60 cells by j == primary\_MP(put):

{\def\LTcaptype{none} % do not increment counter
\begin{longtable}[]{@{}lll@{}}
\toprule\noalign{}
Slice & n & Mean SMS \\
\midrule\noalign{}
\endhead
\bottomrule\noalign{}
\endlastfoot
aligned (j == k, v4 cross-source + v3b primary) & 12 & 0.275 \\
cross (j $\neq$ k) & 48 & 0.061 \\
\end{longtable}
}

\begin{quote}
Data version comparison: - v3 homogeneous + c$\to$MP5: mean\_aligned =
0.183, mean\_cross = 0.064 - v3b homogeneous + c$\to$MP1: mean\_aligned =
0.241, mean\_cross = 0.049 - \textbf{v4 cross-source + c$\to$MP1} (primary):
mean\_aligned = 0.275, mean\_cross = 0.061 - aligned mean continuously
rises (source expansion makes aligned more likely to trigger R-fail),
cross mean slightly rises in v4 (source expansion likewise gives cross
slice more opportunities to trigger R-fail); growth proportions leave $\delta$
nearly unchanged.
\end{quote}

aligned-SMS is higher than cross-SMS in both mean and median, direction
consistent with H1. Note: cross slice median is 0, meaning over half of
cross-MP cells have SMS completely at 0---under cross design, MRs are
almost completely ineffective against mutants. This observation forms
the core empirical fact for §6.1 discussion.

{[}Figure 2: aligned vs cross SMS boxplots{]}

\paragraph{5.7.2 Effect Size and Hypothesis
Testing}\label{effect-size-and-hypothesis-testing}

Nonparametric effect size Cliff's $\delta$ (primary analysis for v3
pre-registered setting; v3b/v4 are exploratory sensitivity, see §3.5.1):

\begin{itemize}
\tightlist
\item
  \textbf{v3 (primary, pre-registered, c$\to$MP5): $\delta$ = 0.323}, 95\%
  bootstrap percentile CI {[}0.017, 0.622{]}
\item
  v3b (exploratory; data-driven c-class primary MP shift to MP1,
  §3.5.1): $\delta$ = 0.446, CI {[}0.154, 0.743{]}
\item
  v4 (exploratory; cross-source 3-LLM pool over identical prompt
  template, §4.2.5): $\delta$ = 0.439, 95\% bootstrap percentile CI
  \textbf{{[}0.127, 0.740{]}} (B = 10,000 iterations, R-12 response)
\item
  $\delta$ under all three stages fails to reach Romano (2006) large-effect
  threshold 0.474. v3 $\to$ v3b +0.123 reflects ``c-class primary MP
  reselection'' impact on $\delta$ (single class / three PUTs, selection rule
  defined after seeing data); v3b $\to$ v4 −0.007 reflects ``nearly
  identical response from three LLMs to identical prompt under fixed
  prompt template'' (see §4.2.5 dry-run), 95\% CI covers zero
\end{itemize}

Median odds ratio (median(aligned) / median(cross)): Because
median(cross) = 0, odds ratio is formally infinite; instead report
``aligned median 0.083 significantly higher than cross median 0.000'' as
auxiliary evidence.

H2 verdict (conjunction of two conditions: Cliff's $\delta$ $\geq$ 0.474 and median
odds ratio $\geq$ 3.0):

\begin{itemize}
\tightlist
\item
  $\delta$ threshold condition: \textbf{not met} (v3 primary $\delta$ = 0.323
  \textless{} 0.474; v3b/v4 exploratory $\delta$ = 0.446 / 0.439 still
  \textless{} 0.474)
\item
  Odds ratio condition: median(cross) = 0 makes ratio undefined; use
  ``aligned median \textgreater{} 0 = cross median'' as auxiliary
  directional observation
\end{itemize}

\textbf{H2 verdict: pre-registered point-estimate criterion not met}
(P0-8 revision: wording changed from ``rejected'' $\to$ ``not met under
pre-registered point-estimate criterion'', reason see effective-n note
below and P0-8 commit message). Under pre-registered primary analysis
(v3, c$\to$MP5, n\_aligned=12, n\_cross=48), Cliff's $\delta$ = 0.323 falls 0.151
short of Romano (2006) large-effect threshold 0.474, 95\% CI lower bound
0.017, effect size classified as small-to-medium. Two exploratory
sensitivity analyses (v3b data-driven primary MP shift, §3.5.1; v4
cross-source pool, §4.2.5) raise $\delta$ to 0.446 / 0.439 but neither crosses
the 0.474 strict threshold, and v3 $\to$ v3b modification includes known
post-hoc selection confound (c-class primary MP selection based on
observed data, see §3.5.1).

\textbf{Effective sample size note} (P1-5 revision, R0 W6 / R1 W7 /
DA-3.1 consensus): surface n\_aligned = 12 and n\_cross = 48, but
§5.6.1.1 already states that in v4 data 75\% cells (45 / 60) have SMS =
0 (zero-mass dominance). In the cross-slice, this means approximately
88\% (42 / 48) of 48 cells are zero; Cliff's $\delta$ inference is actually
dominated by 12 aligned cells + 6 non-zero cross cells,
\textbf{effective n $\approx$ 12 + 6 $\approx$ 18 rather than surface 60}. This actual
sample size constraint: - (a) explains ``power to detect large-effect ($\delta$
\textgreater{} 0.474) only 0.42'' in §5.7.3 power analysis---not caused
by nominal-vs-effective error in sample size; - (b) explains width of
95\% bootstrap CI {[}0.127, 0.740{]}---upper/lower bound ratio $\approx$ 5.83,
reflecting known liberal tendency of percentile bootstrap at effective n
$\approx$ 18; - (c) \textbf{does not change} H2 verdict direction (point
estimate 0.439 \textless{} 0.474 is an effect-size ceiling, not a sample
size issue, as §5.7.3 already argues), but readers should understand CI
width causation accordingly.

Future work (P4) when expanding sample to n $\geq$ 30 PUTs, whether zero-mass
dominance dilutes with PUT class diversification is a testable
hypothesis for effective-n improvement.

\textbf{Contextual observation} (not part of H2 verdict): This paper's
observed $\delta$ is in the same order of magnitude as the LLM-mutant
medium-effect range observed by Tip et al.~(2024) LLMorpheus on
JavaScript. \textbf{Estimand caveat}: Tip 2024 compares ``LLM mutants vs
traditional mutants on fault detection rate'' (cross-mutation-source
comparison), this paper's §5.7.2 compares ``aligned vs cross MP slice on
the same mutant pool'' (single-source within-pool comparison). The two $\delta$
values being similar does not constitute substantive support, only
serves as reference for medium-effect phenomena in LLM-mutant
literature, \textbf{does not constitute weakening or reframing of H2}.

\textbf{v3 $\to$ v3b and v3b $\to$ v4 contrasts reported separately} (avoiding
synthetic ratio implying factor isolation):

{\def\LTcaptype{none} % do not increment counter
\begin{longtable}[]{@{}
  >{\raggedright\arraybackslash}p{(\linewidth - 6\tabcolsep) * \real{0.2500}}
  >{\raggedright\arraybackslash}p{(\linewidth - 6\tabcolsep) * \real{0.2500}}
  >{\raggedright\arraybackslash}p{(\linewidth - 6\tabcolsep) * \real{0.2500}}
  >{\raggedright\arraybackslash}p{(\linewidth - 6\tabcolsep) * \real{0.2500}}@{}}
\toprule\noalign{}
\begin{minipage}[b]{\linewidth}\raggedright
Contrast
\end{minipage} & \begin{minipage}[b]{\linewidth}\raggedright
$\Delta$$\delta$
\end{minipage} & \begin{minipage}[b]{\linewidth}\raggedright
95\% CI (based on bootstrap iterations)
\end{minipage} & \begin{minipage}[b]{\linewidth}\raggedright
Interpretation
\end{minipage} \\
\midrule\noalign{}
\endhead
\bottomrule\noalign{}
\endlastfoot
v3 $\to$ v3b (c-class primary MP shift) & +0.123 & (one-sided data-driven
selection, CI not applicable, see §3.5.1 caveat) & Single class,
post-hoc selection; reflects primary MP sensitivity rather than generic
``MR design contribution'' \\
v3b $\to$ v4 (cross-source under fixed prompt) & −0.007 & CI covers zero &
Three LLMs nearly identical under prompt-fixed; does not constitute
strong test of ``LLM source diversity'' (§4.2.5 limitation) \\
\end{longtable}
}

Future work (P4) will test: (a) pre-registered primary MP selection rule
(avoiding §3.5.1 confound); (b) differential prompt frame per LLM
(testing strong sense of source diversity).

\begin{quote}
Interpretation: RQ2's formalized H2 boundary was not crossed. We have
expanded the mutant pool from 12/PUT to 30/PUT per §7.1.6 R9 (observed
mean 17.4 mutants/PUT, limited by cache capacity); after pool expansion
$\delta$ rose slightly from 0.321 to 0.323, CI narrowed from {[}0.021, 0.639{]}
to {[}0.017, 0.622{]}, indicating effect size is stable on this dataset
and not diluted by pool scale. We therefore characterize H2 as
\textbf{not meeting the large-effect threshold, but medium effect is
stable and numerically consistent with the medium-effect magnitude
reported in LLM-mutant literature (Tip 2024; estimand caveat: see
§5.7.2)}; the feasible path to cross the large-effect threshold (0.474)
is a cross-source mutant pool (mixing Claude / GPT / DeepSeek), as a
research direction for the P4 paper (§7.1.7 R10 already paved the way).
\end{quote}

\paragraph{5.7.3 Statistical Power Analysis (R-13
Response)}\label{statistical-power-analysis-r-13-response}

\begin{quote}
Reviewer R-13 requires: Before reporting ``large-effect threshold not
met'', need to state whether this study's sample size (n\_aligned,
n\_cross) = (12, 48) is sufficient to detect $\delta$ $\geq$ 0.474. This section
uses empirical distribution parametric bootstrap simulation (N\_sim =
5,000, seed = 42) to provide three-tier power.
\end{quote}

\textbf{Method}: \textbf{With replacement} sample from observed aligned
(n=12) and cross (n=48) v4 SMS pools to obtain simulated samples,
compute simulated Cliff's $\delta$; repeat N\_sim times, count frequency of
simulated $\delta$ exceeding given threshold as power approximation at that
threshold (\texttt{scripts/compute\_rq2\_power.py}).

\textbf{Results} (data source:
\texttt{data/results/rq2\_power\_v4.json}):

{\def\LTcaptype{none} % do not increment counter
\begin{longtable}[]{@{}lll@{}}
\toprule\noalign{}
Test threshold & Interpretation & Power at (12, 48) \\
\midrule\noalign{}
\endhead
\bottomrule\noalign{}
\endlastfoot
$\delta$ \textgreater{} 0.000 & Detect any random advantage & \textbf{0.997} \\
$\delta$ \textgreater{} 0.147 & Detect small effect & 0.966 \\
$\delta$ \textgreater{} 0.330 & Detect medium effect & 0.759 \\
\textbf{$\delta$ \textgreater{} 0.474} & \textbf{Detect large effect (H2
threshold)} & \textbf{0.423} \\
\end{longtable}
}

\textbf{Key interpretations}: 1. \textbf{Sample size sufficient for
``any-effect'' detection} (0.997 $\gg$ 0.80), H1 (RQ1) and RQ2 directional
conclusions are robust. 2. \textbf{Slightly insufficient for
medium-effect detection} (0.759 \textless{} 0.80), but 95\% CI already
covers 0.330, consistent with point estimate 0.439 indicating at least
medium effect exists. 3. \textbf{Power to detect large-effect (H2 strict
threshold) only 0.423}: This means even if true $\delta$ is around 0.474,
sample size has only about 42\% probability of detection. \textbf{But
this cannot conversely say ``insufficient power is the cause of H2 not
being met''}---observed $\delta$ = 0.439 \textless{} 0.474, effect size itself
is below threshold; increasing sample size will only narrow CI, will not
automatically elevate point estimate to large. 4. \textbf{Sample size
sweep} (\texttt{rq2\_power\_v4.json}
sample\_size\_curve\_for\_delta\_gt\_0): at n\_aligned $\in$ \{6, 12,
\ldots, 60\} (n\_cross = 4$\times$), power to detect $\delta$ \textgreater{} 0 reaches
0.974 at n\_aligned = 6, 0.996 at 12, then plateaus. Indicates RQ2
primary analysis (detect-any-effect) has very low sample size
requirement, \textbf{main bottleneck is effect size boundary itself, not
sampling noise}.

\textbf{Relationship to H2 verdict}: R-12/R-13 power supplement
\textbf{supports rather than weakens} the H2 ``large-effect threshold
not met'' conclusion --- observed effect is in the medium range and CI
cannot exclude \textless{} 0.474, continuing to claim ``large-effect
threshold not met'' is statistically reasonable. \textbf{Effect-size
ceiling breakthrough requires substantive improvement at MR design level
(P4 direction), not larger sample}.

\textbf{Stipulated-alternative power supplement (R1 W1 round-2 NEW)}:
The plug-in bootstrap above answers ``given the \emph{observed}
distribution ($\delta$ $\approx$ 0.439), how effectively can we detect $\delta$ \textgreater{}
threshold.'' R1 W1 further requires a stipulated-alternative design ---
``if the \emph{truth} is exactly the H2 threshold 0.474, how effectively
can our sample size confirm it?'' We implement the stipulated
alternative via a mixture-weight construction (script
\texttt{scripts/compute\_rq2\_power\_stipulated.py}, N\_sim = 2000): SMS
distributions have heavy ties at zero (45/60 cells = 0), raw shifting
Cliff's $\delta$ is discontinuous (any $\varepsilon$ \textgreater{} 0 jumps $\delta$ from 0.314 to
0.74), so we construct mixture aligned' = (probability w) sample from
(observed\_aligned + 0.001) + (probability 1 − w) sample from
observed\_aligned, calibrating w so that E{[}$\delta$\_stipulated{]} $\approx$ 0.474.
Calibrated w = 0.094, realized E{[}$\delta${]} = 0.4746.

{\def\LTcaptype{none} % do not increment counter
\begin{longtable}[]{@{}
  >{\raggedright\arraybackslash}p{(\linewidth - 6\tabcolsep) * \real{0.2500}}
  >{\raggedright\arraybackslash}p{(\linewidth - 6\tabcolsep) * \real{0.2500}}
  >{\raggedright\arraybackslash}p{(\linewidth - 6\tabcolsep) * \real{0.2500}}
  >{\raggedright\arraybackslash}p{(\linewidth - 6\tabcolsep) * \real{0.2500}}@{}}
\toprule\noalign{}
\begin{minipage}[b]{\linewidth}\raggedright
Stipulated $\delta$\_truth
\end{minipage} & \begin{minipage}[b]{\linewidth}\raggedright
Test criterion
\end{minipage} & \begin{minipage}[b]{\linewidth}\raggedright
Power at (12, 48)
\end{minipage} & \begin{minipage}[b]{\linewidth}\raggedright
Data source
\end{minipage} \\
\midrule\noalign{}
\endhead
\bottomrule\noalign{}
\endlastfoot
0.474 (H2 boundary) & $\delta$\_hat $\geq$ 0.474 (point-estimate criterion, used in
P0-8 verdict) & \textbf{0.491} &
\texttt{rq2\_power\_stipulated\_v4.json} \\
0.474 (H2 boundary) & 95\% CI lower \textgreater{} 0 (any-effect
criterion) & 0.868 & same \\
\end{longtable}
}

\textbf{Stipulated interpretation}: Even if the \emph{true} $\delta$ equals the
H2 threshold 0.474, this design (12, 48) still has only \textbf{49.1\%}
probability that a fresh sample's $\delta$\_hat will reach 0.474 --- this is
the actual power under the point-estimate criterion used by the P0-8
verdict. \textbf{This precisely supports the §5.7.2 P0-8 revision}
(changing ``is rejected'' to ``is not met under the pre-registered
point-estimate criterion''): the H2 verdict is a factual statement about
the point estimate not meeting the threshold, \textbf{not a claim that
the effect size is necessarily smaller than 0.474} --- the stipulated
power of 49.1\% shows that even when the truth is exactly at 0.474,
sample randomness still produces ``not met'' verdicts in roughly half of
replications. This is fully consistent with the §5.7.2 estimand caveat.

\begin{center}\rule{0.5\linewidth}{0.5pt}\end{center}

\subsubsection{5.8 RQ3 Empirical Results (Cross-Class Consistency Across
4 PUT
Classes)}\label{rq3-empirical-results-cross-class-consistency-across-4-put-classes}

\paragraph{5.8.1 Class Means}\label{class-means}

{\def\LTcaptype{none} % do not increment counter
\begin{longtable}[]{@{}
  >{\raggedright\arraybackslash}p{(\linewidth - 6\tabcolsep) * \real{0.2500}}
  >{\raggedright\arraybackslash}p{(\linewidth - 6\tabcolsep) * \real{0.2500}}
  >{\raggedright\arraybackslash}p{(\linewidth - 6\tabcolsep) * \real{0.2500}}
  >{\raggedright\arraybackslash}p{(\linewidth - 6\tabcolsep) * \real{0.2500}}@{}}
\toprule\noalign{}
\begin{minipage}[b]{\linewidth}\raggedright
Class
\end{minipage} & \begin{minipage}[b]{\linewidth}\raggedright
PUT set
\end{minipage} & \begin{minipage}[b]{\linewidth}\raggedright
Mean SMS (15 cells, \textbf{v3 baseline})
\end{minipage} & \begin{minipage}[b]{\linewidth}\raggedright
Mean SMS (\textbf{v4 cross-source})
\end{minipage} \\
\midrule\noalign{}
\endhead
\bottomrule\noalign{}
\endlastfoot
a (numeric) & a1, a2, a3 & 0.067 & 0.067 \\
b (probabilistic) & b1, b2, b3 & 0.156 & 0.148 \\
c (surrogate) & c1, c2, c3 & 0.047 & 0.089 \\
d (ML) & d1, d2, d3 & 0.081 & 0.112 \\
\end{longtable}
}

Across both versions the maximum between classes occurs in class b; the
minimum is class c under v3 (0.047) and class a under v4 (0.067).
\textbf{Cross-source main driver: class c +91.4\% (see §6.1)} --- class
c is most sensitive to cross-source pool expansion, and together with
the §3.5.1 v3b primary MP shift forms the empirical finding that ``class
c sensitivity to source diversity exceeds the other three classes''.

{[}Figure 3: Cross-class SMS forest plot (mean ± SEM){]}

\paragraph{5.8.2 Sign Test: Within-Class Aligned Higher Than
Cross}\label{sign-test-within-class-aligned-higher-than-cross}

For each class, compute ``that class's aligned slice mean − that class's
cross slice mean'', record 1 if sign is positive.

\begin{itemize}
\tightlist
\item
  Pass count: \textbf{v3 (pre-registered): 3 / 4 (partial)}; \textbf{v3b
  (exploratory, post-hoc): 4 / 4 (conditional on c-class primary MP
  shift, §3.5.1)}; v4 cross-source pool maintains 4 / 4 under v3b
  condition.
\end{itemize}

H4 primary conclusion based on v3 pre-registered: \textbf{partial, not
strict} (sign test 3/4). v3b 4/4 and v4 4/4 are sensitivity reports, do
not replace v3 verdict. Class c's flip under v3b stems from data-driven
primary MP adjustment (§3.5.1): class c Friedman per-class p = 0.406
shows any MP can serve as primary; v3b replaces P1's MP5 with the one
having maximum mean SMS (MP1), making class c aligned mean (0.233)
\textgreater{} cross mean (0); this is selection-on-response, therefore
exploratory only. Detailed directional interpretation in §6.3.

\paragraph{5.8.3 Mixed-Effects Model Limitation
Statement}\label{mixed-effects-model-limitation-statement}

§5.3.2 planned random-intercept-PUT, fixed-effects class $\times$ operator
model \textbf{primary model did not converge} on observed data:

\begin{itemize}
\tightlist
\item
  Primary model formula:
  \texttt{sms\ \textasciitilde{}\ C(class)\ +\ C(operator)\ +\ C(class):C(operator)\ +\ (1\ \textbar{}\ put)}
\item
  Fit error: Singular matrix (design matrix for class $\times$ operator
  interaction term has insufficient column rank, N=60 observations
  insufficient for 11-dimensional fixed-effects + 12 PUT random
  intercepts)
\item
  Fallback model (removing interaction term):
  \texttt{sms\ \textasciitilde{}\ C(class)\ +\ C(operator)\ +\ (1\ \textbar{}\ put)}
\item
  Fallback status: fixed-effects converge, but Group Var (PUT random
  intercept variance) hits boundary $\approx$ 0, essentially degenerates to OLS;
  model self-reports ``Group Var hit boundary; the random-intercept term
  is degenerate''
\end{itemize}

Class fixed-effects p-values (class a as baseline, fallback model, as
approximate descriptive metric):

{\def\LTcaptype{none} % do not increment counter
\begin{longtable}[]{@{}ll@{}}
\toprule\noalign{}
Comparison & p-value (approx.) \\
\midrule\noalign{}
\endhead
\bottomrule\noalign{}
\endlastfoot
class b vs a & 0.275 \\
class c vs a & 0.892 \\
class d vs a & 0.991 \\
\end{longtable}
}

\begin{quote}
\textbf{Honest statement}: Due to PUT random intercept variance
degeneration and sample size (60 cells / 12 PUTs) limitations, we do not
report fallback p-values as formal hypothesis tests, but as auxiliary
description. RQ3's primary conclusion shifts to four-piece direct
presentation: (a) class mean table + (b) sign test + (c) Friedman
nonparametric test (§5.8.4) + (d) forest plot, consistent with §5.3.2
already-stated ``small-N multiple comparison alternative''. This
limitation is extended into §7.2.2 R6.
\end{quote}

\paragraph{5.8.4 Friedman Nonparametric Alternative
Test}\label{friedman-nonparametric-alternative-test}

Because mixed-effects primary model is Singular, we use Friedman $\chi$$^{2}$ as
formal nonparametric alternative test:

\begin{itemize}
\tightlist
\item
  \textbf{Design}: PUT (block, n=12) $\times$ MP (treatment, k=5), value = SMS,
  df = 4
\item
  \textbf{Primary statistic}: \textbf{$\chi$$^{2}$ = 15.30, p = 0.0041}
  (\textless{} 0.05, significant)
\item
  \textbf{MP rank means} (MP1 $\to$ MP5): 2.92, 2.58, 2.08, 3.08, 4.33
\end{itemize}

Within-class Friedman (each class 3 PUTs $\times$ 5 MPs) --- \textbf{R1 W4
round-2: Bonferroni $\times$ 4 correction + Kendall's W effect size added}:

{\def\LTcaptype{none} % do not increment counter
\begin{longtable}[]{@{}
  >{\raggedright\arraybackslash}p{(\linewidth - 10\tabcolsep) * \real{0.1667}}
  >{\raggedright\arraybackslash}p{(\linewidth - 10\tabcolsep) * \real{0.1667}}
  >{\raggedright\arraybackslash}p{(\linewidth - 10\tabcolsep) * \real{0.1667}}
  >{\raggedright\arraybackslash}p{(\linewidth - 10\tabcolsep) * \real{0.1667}}
  >{\raggedright\arraybackslash}p{(\linewidth - 10\tabcolsep) * \real{0.1667}}
  >{\raggedright\arraybackslash}p{(\linewidth - 10\tabcolsep) * \real{0.1667}}@{}}
\toprule\noalign{}
\begin{minipage}[b]{\linewidth}\raggedright
Class
\end{minipage} & \begin{minipage}[b]{\linewidth}\raggedright
$\chi$$^{2}$
\end{minipage} & \begin{minipage}[b]{\linewidth}\raggedright
raw p
\end{minipage} & \begin{minipage}[b]{\linewidth}\raggedright
Bonferroni $\times$ 4 adjusted p
\end{minipage} & \begin{minipage}[b]{\linewidth}\raggedright
Kendall's W (n=3, k=5)
\end{minipage} & \begin{minipage}[b]{\linewidth}\raggedright
Effect strength
\end{minipage} \\
\midrule\noalign{}
\endhead
\bottomrule\noalign{}
\endlastfoot
a (numeric) & 4.00 & 0.406 & 1.000 & 0.333 & small \\
b (probabilistic) & 10.78 & \textbf{0.029} & 0.116 & 0.898 & large \\
c (surrogate) & 4.00 & 0.406 & 1.000 & 0.333 & small \\
d (ML) & 5.00 & 0.287 & 1.000 & 0.417 & small-medium \\
\end{longtable}
}

\textbf{R1 W4 round-2 correction interpretation}: Applying Bonferroni $\times$
4 correction over the family of 4 per-class Friedman tests, the b-class
original ``individually significant'' raw p = 0.029 rises to adjusted p
= 0.116, \textbf{\textgreater{}} 0.05; \textbf{no per-class result
remains significant at family-wise $\alpha$ = 0.05}. This is consistent with
the small per-class sample N = 3 PUTs $\times$ 5 MPs --- even though b-class
Kendall's W = 0.898 indicates a strong rank-concordance pattern, 3 PUTs
is insufficient to yield a conclusive within-class verdict after
multiple-comparison correction. \textbf{This is consistent with the
paper's H4 primary verdict still resting on the §5.8.2 sign test}:
per-class Friedman serves as sensitivity / descriptive reporting only
and does not drive the H4 verdict directly. Kendall's W provides the
effect-size complement (b-class 0.898 falls in the ``large concordance''
range under Cohen 1988, even though not statistically significant after
correction).

Interpretation: Friedman test on full 60-cell data is
\textbf{significant (p = 0.0041)}, indicating 5 MPs do exhibit
systematic differences across 12 PUT blocks; rank means show MP3 lowest,
MP5 highest, consistent with §5.7 aligned-cross asymmetry pattern.

\textbf{Important caveat --- Friedman main effect $\neq$ H4 cross-class
consistency}: Friedman $\chi$$^{2}$ tests ``whether rank differences exist among 5
MPs'' (MP main effect averaged over PUTs), \textbf{not what H4 tests:
``whether direction is consistent across 4 PUT classes in aligned-cross
comparison''}. These are logically independent statistical claims:

{\def\LTcaptype{none} % do not increment counter
\begin{longtable}[]{@{}
  >{\raggedright\arraybackslash}p{(\linewidth - 6\tabcolsep) * \real{0.2500}}
  >{\raggedright\arraybackslash}p{(\linewidth - 6\tabcolsep) * \real{0.2500}}
  >{\raggedright\arraybackslash}p{(\linewidth - 6\tabcolsep) * \real{0.2500}}
  >{\raggedright\arraybackslash}p{(\linewidth - 6\tabcolsep) * \real{0.2500}}@{}}
\toprule\noalign{}
\begin{minipage}[b]{\linewidth}\raggedright
Statistic
\end{minipage} & \begin{minipage}[b]{\linewidth}\raggedright
What it tests
\end{minipage} & \begin{minipage}[b]{\linewidth}\raggedright
This paper's result
\end{minipage} & \begin{minipage}[b]{\linewidth}\raggedright
Support for H4
\end{minipage} \\
\midrule\noalign{}
\endhead
\bottomrule\noalign{}
\endlastfoot
Friedman $\chi$$^{2}$ (PUT $\times$ MP) & MP rank differences (averaged over 12 PUTs) &
$\chi$$^{2}$ = 15.30, p = 0.0041 (significant) & Does not directly support \\
Sign test (4 classes) & 4/4 classes aligned mean \textgreater{} cross
mean & v3 primary: 3/4 (partial); v3b exploratory: 4/4 &
\textbf{Directly corresponds to H4} \\
\end{longtable}
}

Within-class Friedman only class b (probabilistic) individually
significant, reflecting probabilistic PUT class has larger sensitivity
differences to different MPs. Non-significance of Friedman per-class
(a/c/d classes p \textgreater{} 0.28) \textbf{does not mean ``any MP
equivalent'' within these classes}, but rather that under n=3 PUT power,
MP differences were not detected. §3.5.1 v3b exploratory uses class c
non-significance as basis for primary MP reselection, which is
selection-on-non-significance, itself having §3.5.1 already-stated
confound.

Friedman's methodological contribution is \textbf{providing formalized
nonparametric p-value for RQ3} (replacing mixed-effects Singular,
§5.8.3), not \textbf{directly verifying H4}. H4 verdict still based on
§5.8.2 sign test: v3 primary 3/4 (partial), v3b exploratory 4/4
(post-hoc, with caveats).

\begin{center}\rule{0.5\linewidth}{0.5pt}\end{center}

\subsubsection{5.9 RQ4 Empirical Results (SMS vs Pattern
Coverage)}\label{rq4-empirical-results-sms-vs-pattern-coverage}

\paragraph{5.9.1 Pattern Coverage
Operationalization}\label{pattern-coverage-operationalization}

For each PUT, compute (MP\_k, R\_outcome $\in$ \{True, False\}) binary tuple
coverage: each PUT has 5 MPs $\times$ 2 outcomes = 10 cells, coverage = number
of actually triggered cells / 10. Essentially the simplest
implementation of §1.4 RQ4 baseline (per-PUT granularity, not
distinguishing mutants).

PC range for 12 PUTs: {[}0.500, 1.000{]}, mean 0.733.

\paragraph{5.9.2 Correlation with SMS}\label{correlation-with-sms}

Pair by PUT (one PC value per PUT, paired with that PUT's mean SMS over
5 MPs):

\begin{itemize}
\tightlist
\item
  Spearman $\rho$ = 0.163 (p = 0.613) (v4 primary,
  \texttt{paper\_numbers\_v4.json})
\item
  Kendall $\tau$ = 0.136 (p = 0.568)
\end{itemize}

{[}Figure 5: per-PUT SMS vs PC scatter (n = 12){]}

\paragraph{5.9.3 Interpretation}\label{interpretation}

\textbf{Statistical-power caveat first}: n = 12 PUTs is severely
under-powered. For a Spearman $\rho$ at n = 12, the 95\% CI covers
approximately {[}−0.5, +0.6{]}; the measurement precision is
insufficient to distinguish ``zero correlation'', ``moderate positive
correlation'', or ``moderate negative correlation''. p = 0.61 / 0.57
does not constitute evidence that ``correlation does not exist'', only
that ``no correlation was detected at n = 12''.

\textbf{Revised qualitative observation} (weakened from the original
``almost independent''): Spearman $\rho$ = 0.163 and Kendall $\tau$ = 0.136 are
not significantly different from zero (p $\approx$ 0.61, 0.57), and the data
\textbf{do not support} either the ``SMS is independent of PC'' claim or
the ``SMS is strongly correlated with PC'' claim. \textbf{Orthogonal
semantic dimension is a hypothesis, not a finding from this dataset} ---
its empirical test requires n $\geq$ 30 PUTs or a more refined PC
operationalization (incorporating the mutant dimension). This paper
records ``no detectable statistical correlation between SMS and PC over
12 PUTs'' as the conservative finding for RQ4, leaving orthogonality to
the P4 paper (with expanded PUT scale or partial-correlation design).

\begin{quote}
Specifically, the b2 PUT has PC = 1.0 (full coverage) but mean SMS =
0.067; the b1 PUT has PC = 0.7 but mean SMS = 0.20 --- within the same
class, the PUT with higher PC has lower SMS, contradicting the naive
assumption that ``more complete PC kills more mutants'' and providing
positive support for SMS as an independent information dimension. The c3
PUT shows the same pattern: PC = 1.0 with mean SMS = 0.14; c1 / c2 have
PC = 0.7-0.8 with mean SMS = 0.0, with the inverse correlation more
pronounced within that class. The completion of RQ4 is bounded at this
level: further refinement of the PC definition (incorporating the mutant
dimension, cross-PUT joint coverage) is left for the P4 paper (see
§1.6).
\end{quote}

\subsection{Section 6 $\cdot$ Discussion}\label{section-6-discussion}

\textbf{Section positioning}: This section discusses the empirical
findings from the §5 60-cell empirical audit (H1/H2/H4/H5 verdicts,
Cliff's $\delta$, Friedman main effects, SMS-PC correlations). These are
\textbf{incidental empirical findings} after the three-pillar
methodological framework is established---they demonstrate the empirical
ceiling within the current scope of LLM-mutant + same-prompt +
single-output kernels, and do not constitute counterevidence to the
methodological framework (the framework's justification is completed in
§3.2.0 / §2.3-4.4 / §3.2.6.3).

\subsubsection{6.1 Systematic bias of SMS on aligned
slices}\label{systematic-bias-of-sms-on-aligned-slices}

The core empirical fact from RQ2: the median SMS of cross-MP slices (j $\neq$
k) is 0, meaning most non-aligned MRs are almost completely ineffective
against LLM-generated semantic mutants; the median SMS of aligned slices
is 0.267, far higher than 0. This asymmetry stems from the semantic
coherence between mut\_j and MP\_k---in aligned slices, mutants break
the algebraic properties directly asserted by the MR, thus the R
detection signal is strongest; in cross slices, the semantic dimension
broken by mutants is orthogonal to the semantic dimension detected by
the MR, so even if the mutation does change the output, the MR cannot
capture that change as R-fail.

\textbf{Key finding from Phase A cross-source pool (one of this paper's
core methodological contributions)}: We retested H2 using a cross-source
mutant pool (v4, see §4.2.5) generated by three LLMs (Claude Opus 4.6 +
GPT-5.4 + DeepSeek chat), obtaining $\delta$ = 0.439, nearly identical to the
v3b same-source pool $\delta$ = 0.446 (difference 0.007), with slightly wider
95\% CI but stable center. \textbf{This inversely falsifies our initial
hypothesis about the H2 upper bound---LLM same-source bias is not the
dominant factor in the H2 upper bound}: - The leap from v3 $\to$ v3b ($\delta$
0.323 $\to$ 0.446) came from data-driven adjustment of c-class primary MP
(§3.5.1), i.e., the MR-MP alignment design itself - The cross-source
expansion from v3b $\to$ v4 ($\delta$ 0.446 $\to$ 0.439) barely changed $\delta$, i.e., source
diversity of the mutant pool

This contrast \textbf{separately} reports two contrasts: $\Delta$$\delta$\_\{v3$\to$v3b\}
= +0.123 comes from single-class, post-hoc c-class primary MP selection
(§3.5.1 caveat, Bonferroni-bounded effective $\alpha$ see §3.5.1 + this paper's
P0-4 revisions); $\Delta$$\delta$\_\{v3b$\to$v4\} = −0.007, conditional on v3b selection +
identical prompt, reflects the near-consistent response of three LLMs
under fixed prompt. \textbf{The two contrasts each carry their own
selection / conditioning caveats and cannot be synthesized into a single
factor decomposition ratio}---see the existing contrast table in §5.7.2.
Tip et al.~(2024) LLMorpheus observed LLM-mutants in the medium-effect
range on JavaScript---this phenomenon is directionally consistent with
this study's observation that $\delta$ has not yet crossed the 0.474
large-effect threshold, suggesting that under LLM-mutant + current MR
design, medium-effect may be the norm for this experimental paradigm,
and crossing it requires not more sources but further refinement of MR
design (P4 paper research direction). \textbf{Estimand caveat}: Tip
2024's $\delta$ is ``LLM mutants vs traditional mutants on fault detection
rate'' (cross-mutation-source comparison), different from this paper's
``aligned vs cross MP slice'' (single-source within-pool) estimand;
numerical similarity serves only as a literature reference for the
medium-effect phenomenon.

\textbf{True contribution of the cross-source pool}: Although it does
not significantly push $\delta$, it significantly improves mutant quality
(C1\_share 0.164 $\to$ 0.209, +27\%) and inter-class balance (c-class mean
SMS 0.047 $\to$ 0.089, +91.4\%). This constitutes the empirical foundation
for the R\_sem/R\_kill decoupling engineering insight in §6.2.

\textbf{Contextual consistency with LLM-mutant literature (not H2
verdict)}: This paper's v3 primary $\delta$ = 0.323, CI {[}0.017, 0.622{]} is
in the same order of magnitude as the LLM-mutant medium-effect range
observed by Tip et al.~(2024) LLMorpheus on JavaScript. \textbf{Estimand
caveat}: Tip 2024 compares ``LLM mutants vs traditional mutants on fault
detection rate'' (cross-mutation-source comparison), while this paper's
§5.7.2 compares ``aligned vs cross MP slice on the same mutant pool''
(single-source within-pool comparison). The numerical similarity of the
two $\delta$ values does not constitute substantive support, serving only as a
reference for the medium-effect phenomenon in LLM-mutant literature.
This is contextual literature comparison, \textbf{not a reframing of the
H2 rejected verdict}; it merely shows that the medium-effect scale
observed in this study is consistent with the empirical norm of this
research paradigm in other domains (JS web testing), providing a
baseline for subsequent P4 cross-language portability research.

\textbf{Petrović \& Ivanković (2018) numerical coincidence statement
(mechanism difference confirmed)}: Petrović \& Ivanković reported
\textasciitilde20\% productive mutants on Google's industrial 500k
mutant data, numerically close to this paper's LRCA-calibrated best
C1\_share = 0.20. However, \textbf{this is numerical coincidence, not
mechanism validation}: Google's ``productive mutant'' is obtained from
\emph{developer survey} where developers subjectively judge usefulness,
a human judgment construct; LRCA's C1 is the automatic annotation output
of a \emph{3-layer classifier}, an algorithmic construct. The two
constructs differ significantly; numerical approximation does not
constitute proof of LRCA's equivalence to industrial practice. Rigorous
mechanism validation would require developer sampling review on 5-10 PUT
$\times$ N C1-flagged mutants (reserved for P4 paper). This paper positions
this comparison as ``encouraging numerical reference,'' not ``industrial
practice calibration.''

\textbf{Note on monotone transformation invariance (not a robustness
check)}: Cliff's $\delta$ is a rank-based statistic, function U / (n1$\cdot$n2),
\textbf{mathematically invariant under any strictly monotone
transformation} (Romano 2006 explicitly notes). Logit is strictly
monotone on (0, 1), so $\delta$\_logit ≡ $\delta$\_raw after logit transformation of
SMS is a construction result, not an empirical finding.

We still retain the calculation result as an execution correctness check
(\texttt{data/results/rq2\_cliffs\_delta\_logit\_v4.json} shows $\delta$\_logit
= 0.439 = $\delta$\_raw, difference 0.000---consistent with the rank-invariance
theorem), \textbf{but this does not constitute additional robustness
evidence for the H2 conclusion}. The true robustness threats to the H2
verdict come from the post-hoc selection confound in §3.5.1 v3b and
zero-mass dominance in §5.6.1.1, not from metric scale choice.

\subsubsection{6.2 Decoupling of R\_sem and
R\_kill}\label{decoupling-of-r_sem-and-r_kill}

The §4.8.3 operator-level pilot already observed: HP-class
(hyperparameter) operators have high R\_sem (semantic feasibility) but
low R\_kill (kill rate by MR). This chapter's §5.6 cell-level SMS
reproduces this pattern---45 / 60 cells have SMS = 0 (v4 cross-source
pool), concentrated in cross-MP slices; the v3b same-source pool gives
mean C1\_share of only 0.164 under default LRCA threshold, \textbf{v4
cross-source pool improves to 0.209 (+27\%) under default threshold},
i.e., among the few killed mutants, the cross-source pool significantly
reduces LRCA mislabeling rate.

Engineering insight: ``operator-MP alignment coverage'' in MR design is
a necessary condition for producing strong SMS signals; merely expanding
the ``semantically feasible mutant pool'' does not increase SMS but
dilutes the proportion. LRCA threshold calibration (OOD boundary 0.05,
tolerance multiplier 10$\times$) may also be overly sensitive, misjudging most
borderline kills as C2/C3/C4 rather than C1. These two jointly support
the research question in the P4 paper of ``using SMS to infer MR
under-coverage dimensions + calibrating LRCA.''

\subsubsection{6.3 Constrained interpretation of cross-class consistency
H4}\label{constrained-interpretation-of-cross-class-consistency-h4}

RQ3 data show that the mean SMS of all 4 classes is positive (v4
cross-source: a=0.067, b=0.148, c=0.089, d=0.112), but the differences
are limited. \textbf{H4 main verdict is based on v3 pre-registered: sign
test 3/4 (partial)}; after v3b data-driven primary MP adjustment, sign
test rises to 4/4 (exploratory, post-hoc, §3.5.1), v4 cross-source pool
maintains 4/4 under v3b conditions (also exploratory). Inter-class
balance significantly improved after cross-source: c-class mean SMS
0.047 $\to$ 0.089 (+91.4\%), d-class 0.081 $\to$ 0.112 (+38\%), a/b classes
nearly unchanged. This confirms that c/d classes (surrogate / ML) have
higher demand for mutant diversity than a/b classes (numeric /
probabilistic). The mixed-effects primary model is Singular, the
fallback model's PUT random intercept variance degenerates to 0, meaning
the sample size of 60 cells / 12 PUTs is insufficient to support a
two-layer structure that ``estimates both random intercepts and
fixed-effect interaction terms.''

However, §5.8.4 Friedman non-parametric test gives \textbf{significant
main effect ($\chi$$^{2}$ = 15.30, p = 0.0041)}, confirming systematic differences
among the 5 MPs on 12 PUT blocks. Within-class Friedman shows only
b-class (probabilistic) is individually significant (p = 0.029), the
other three classes p \textgreater{} 0.28---this is consistent with §5.7
aligned-cross asymmetry, indicating MP-differences are most evident on
probabilistic-class PUTs. \textbf{c-class internal p = 0.406
non-significant is precisely the basis for §3.5.1 data-driven primary MP
adjustment}: any MP can serve as primary, selecting the one with largest
mean SMS (MP1) is a reasonable post-hoc revision.

Our conclusion: cross-class consistency verdict \textbf{is based on v3
pre-registered = partial (sign test 3/4)}. Supporting sensitivity: (a)
all 4 class mean directions are positive (v3/v3b/v4); (b) v3b sign test
4/4 and v4 sign test 4/4 (both exploratory, conditional on c-class
primary MP shift, §3.5.1); (c) 60-cell Friedman p = 0.0041
(non-parametric fallback, \textbf{not part of H4 verdict}, see §5.8.4).
Mixed-effects unavailability is a sample constraint of N = 60 / 12 PUT
(§7.2.2 R6), not evidence absence.

\subsubsection{6.4 Positional relationship between SMS and Pattern
Coverage}\label{positional-relationship-between-sms-and-pattern-coverage}

RQ4's Spearman $\rho$ = 0.163, Kendall $\tau$ = 0.136 (n = 12 PUT, p $\approx$ 0.61, 0.57)
give the correlation direction between SMS and the simplest PC.
\textbf{n = 12 is severely underpowered}: at this sample size, Spearman
$\rho$'s 95\% CI is approximately {[}−0.5, +0.6{]}, p = 0.74 does not
constitute evidence of ``no correlation,'' only ``not detected.''

This paper's RQ4 conservative conclusion is: \textbf{at the 12 PUT
scale, SMS and the simplest PC have no detectable statistical
correlation}. Orthogonality is a hypothesis worth testing in P4,
\textbf{not a finding from this data}. ``SMS is a semantic-layer
sensitivity dimension orthogonal to coverage-class metrics'' as a P4
research direction (expand to n $\geq$ 30 PUT, or use extended PC
operationalization to incorporate mutant dimension) is not within this
paper's empirical scope.

\subsubsection{6.5 Stakeholder analysis: who benefits from SMS
(R-19)}\label{stakeholder-analysis-who-benefits-from-sms-r-19}

\begin{quote}
\textbf{Scope} (R3 round-2 tightening): All stakeholder pain points,
capability claims, and workflows discussed in this section are strictly
bounded to \textbf{single-output \texttt{float\ $\to$\ float} kernels
(\textless{} 2 KB source code, the 12 PUTs in this paper)} --- see §1.1
/ §3.1.1 / §7.5. None of the claims in this section apply to
industrial-scale, multi-module, or multi-output scientific computing
software; the latter is the research scope of P5 / P2-CN (domain
application). Per IST convention this section explicitly distinguishes
three stakeholder classes --- test engineers, MR designers, V\&V
documentation --- and for each class explains: (a) pain point before
SMS; (b) specific capability SMS provides; (c) executable workflow +
resource cost; (d) existing responsibilities not replaced.
\end{quote}

\paragraph{6.5.1 Test engineers}\label{test-engineers}

\textbf{Pain point} (scoped to single-output kernels): When testing a
single-output \texttt{float\ $\to$\ float} numerical / probabilistic /
surrogate / ML kernel, even when line / branch / MC-DC code coverage
reports 90\%+, the engineer still cannot directly judge whether the
accompanying MR set has adequate detection capability over numerical /
probabilistic / ML behaviors, because existing mutation testing tools
(mutmut, cosmic-ray) produce mostly syntactic-layer mutants (see
§3.2.6.1 operator-level cross-table, §3.2.6.3 12-PUT empirical 5.14\%
AST overlap rate), with weak correlation to domain-level semantic
errors.

\textbf{Capability SMS provides}: For each (PUT, MP) cell, provide a
scalar SMS value from 0--1, directly reading ``how many meaningful
semantic mutations this MR kills under this MP.'' When SMS is
significantly lower than the §5.7.2 same-class baseline (e.g., aligned
mean 0.275), it prompts that the MR's oracle tolerance or verification
logic may miss a class of domain errors.

\textbf{Workflow}: 1. For a newly written MR, run
\texttt{scripts/sms\_campaign.py} to obtain SMS value; 2. If SMS = 0
(75\% default case), first check LRCA labels in
\texttt{lrca\_60cell.json}---C2 (MP-semantic inconsistency) / C3 (AVP
tolerance dispute) / C4 (MR design defect) each point to different
repair paths; 3. Rerun after revising MR; iterate until SMS enters
reasonable range (§5.6 gives empirical ranges by PUT category).

\textbf{Air-gap incompatibility declaration} (NEW, R3 round-2): This
workflow depends on \textbf{external LLM API calls} (Claude / GPT /
DeepSeek) for semantic mutant generation. This makes the method
\textbf{incompatible with most regulated air-gapped V\&V workflows} ---
code review in high-risk domains such as nuclear engineering (IEC
60880), aerospace (DO-178C), medical devices (IEC 62304), and automotive
(ISO 26262) typically requires air-gapped build environments, where LLM
API calls cannot execute. Possible mitigation paths (P5 future work):
(i) self-hosted open-weight LLMs (Llama / DeepSeek local inference);
(ii) offline-cached pre-generated mutant pools with signature locking.
This paper's mutant pool is already pre-generated reproducibly via
commit-hash + raw-response store (§4.2.3), so engineers can
\textbf{reproduce a published mutant pool offline}, but
\textbf{generating new mutant pools still requires external LLM access
(off air-gap)}. This is a deployment scope limitation of the P2
methodology, explicitly acknowledged in §7.5 Limitations.

\textbf{SMS does not replace}: Domain expert manual review of MR
physical/mathematical correctness, regression testing against real
historical faults, performance and stability testing. SMS only evaluates
``semantic-layer failure detection capability,'' not ``engineering
practical value'' (§7.5 R6 already declared).

\paragraph{6.5.2 MR designers}\label{mr-designers}

\textbf{Pain point}: MR designers (often researchers or senior
developers working on single-output scientific computing kernels of the
type covered by §3.1.1) rely on intuition or literature patterns to
propose new MRs, lacking a quantifiable ``design feedback''
metric---before P1, whether a new MR is ``useful'' can only be verified
after deployment to real faults, with a very long feedback loop.

\textbf{Capability SMS provides}: SSOT-based offline SMS batch runs
(\texttt{scripts/sms\_campaign.py}) give MR designers a quantifiable
design-feedback metric --- comparing the new MR's aligned-cell SMS
against the existing-MR-set median (this paper's v4 main analysis
aligned mean 0.275) gives an initial verdict within hours on whether the
MR has acceptable failure-detection capability on aligned slices.

\textbf{Workflow} (reframed to \textbf{quarterly batch audit}, not
per-PR gating; R3 round-2):

\begin{quote}
\textbf{Removed in R3 round-2}: The original §6.5.2 GitHub Actions
per-PR YAML template has been deleted --- the original template
hardcoded \texttt{SMS\_VERSION=v4} + \texttt{P2\_PRIMARY\_VERSION=v3b}
(§3.5.1 v3b is an exploratory post-hoc selection, not pre-registered)
into a stakeholder-facing CI template, propagating the v3b
selection-on-the-response confound into downstream adopters' pipelines,
and the PR-CI threshold 0.10 contradicted the §6.5.3 audit threshold
0.20. This conflicted with the §3.5.1 caveat and is removed in R3
round-2 revision.
\end{quote}

\textbf{Replacement}: \textbf{Quarterly batch audit}, not per-PR gating
---

{\def\LTcaptype{none} % do not increment counter
\begin{longtable}[]{@{}
  >{\raggedright\arraybackslash}p{(\linewidth - 4\tabcolsep) * \real{0.3333}}
  >{\raggedright\arraybackslash}p{(\linewidth - 4\tabcolsep) * \real{0.3333}}
  >{\raggedright\arraybackslash}p{(\linewidth - 4\tabcolsep) * \real{0.3333}}@{}}
\toprule\noalign{}
\begin{minipage}[b]{\linewidth}\raggedright
Step
\end{minipage} & \begin{minipage}[b]{\linewidth}\raggedright
Description
\end{minipage} & \begin{minipage}[b]{\linewidth}\raggedright
Resource cost
\end{minipage} \\
\midrule\noalign{}
\endhead
\bottomrule\noalign{}
\endlastfoot
1. Offline mutant pool generation & Use LLM API to generate the mutant
set corresponding to newly added MRs (per PUT 24-30 mutants $\times$ 12 PUTs $\approx$
300 mutants) & LLM API: \textasciitilde\$5-15 per quarter; wall time
\textasciitilde30 min \\
2. SMS batch run &
\texttt{python\ scripts/sms\_campaign.py\ -\/-track\ 2} & 4-core laptop:
\textasciitilde10-20 min per pool \\
3. LRCA three-layer diagnostic & \texttt{python\ scripts/run\_lrca.py}
annotates C1-C5 & \textless1 min \\
4. Quarterly review & MR design team manually reviews MRs whose
aligned-cell SMS is significantly below the historical median (0.275) &
Team meeting: 1-2 hours \\
\end{longtable}
}

\textbf{Resource estimate} (quarterly): API \textasciitilde\$10 +
\textasciitilde30 min automation + 1-2 h manual review $\approx$ 0.5
person-day/quarter --- affordable for an active MR design team.
\textbf{Per-PR gating is not recommended} (LLM API latency 5-30 s + cost
+ non-determinism + air-gap incompatibility all make per-PR
impractical).

\textbf{SMS does not replace}: Domain expert judgment of MR physical
meaning reasonableness (§7.1.4 R4 already declared MR ``reasonableness''
is a priori); nor does it replace ``whether MR covers a real business
scenario'' stakeholder review (completed by product requirements side).

\paragraph{6.5.3 Research-grade evidence for V\&V documentation
(long-term
aspiration)}\label{research-grade-evidence-for-vv-documentation-long-term-aspiration}

\begin{quote}
\textbf{R3 round-2 retitled} from ``Auditors / certification bodies'':
This subsection makes \textbf{no} normative claim toward industrial
certification bodies (NRC, FDA, ISO 26262 review teams). Within this
paper's single-output kernels scope there is no traceable mapping to the
normative bodies of current IEC 60880 / ISO 26262 / DO-178C / ASME V\&V
20-2009 standards. The subsection is positioned as a \textbf{long-term
aspiration} --- SMS as research-grade test adequacy evidence may serve
as \textbf{supplementary} (not normative) material in V\&V
documentation.
\end{quote}

\textbf{Pain point} (scoped to V\&V documentation, not normative
certification): V\&V documentation for scientific computing software
(per ASME V\&V 20-2009 §3 code verification and similar guides) requires
quantifiable evidence of test adequacy. Current V\&V documentation often
relies on code coverage reports + MR lists + domain SME signatures and
lacks a mutation-based quantifiable supplement.

\textbf{SMS as research-grade supplementary evidence}: SMS values may
appear in V\&V documentation as supplementary evidence, alongside code
coverage and MR lists. Documentation may include: (a) aligned-cell SMS
for each critical PUT; (b) LRCA three-layer diagnostic conclusions (C1
direct readout / C2-C5 engineering attribution); (c) visualization
coverage map of the 60-cell matrix (§6.2 fig1). Reviewers can
independently run the reproduction package (\texttt{REPRODUCIBILITY.md})
to verify.

\begin{quote}
\textbf{R3 round-2 deletion}: The original §6.5.3 proposed acceptance
thresholds aligned-cell SMS $\geq$ 0.20 / $\geq$ 0.30 + C1\_share $\leq$ 0.20. These
thresholds were research recommendations on the empirical basis of this
paper, \textbf{but}: (i) they have no normative backing in IEC / ISO /
ASME standards; (ii) writing them under ``acceptance threshold
recommendations'' risks reader misinterpretation as enforce-ready
standards; (iii) the §3.5.1 caveat already declares that the 0.275
baseline is influenced by v3b post-hoc selection. R3 round-2 therefore
removes the threshold recommendations --- SMS values should be reported
as \textbf{descriptive supplementary evidence}, with concrete ``is it
adequate'' judgments left to actual V\&V documentation reviewers on
case-by-case engineering grounds.
\end{quote}

\textbf{Relation to existing V\&V standards}: This work is conceptually
complementary to ASME V\&V 20-2009 (Standard for Verification and
Validation in Computational Fluid Dynamics and Heat Transfer) §3 code
verification --- the latter emphasizes code-level correctness of
numerical solvers, while this paper's SMS emphasizes the fault-detection
adequacy of MR sets. However, this paper's empirical scope is strictly
bounded to single-output \texttt{float\ $\to$\ float} kernels, with a
substantial scale gap to the multi-module industrial CFD code base
targeted by ASME V\&V 20-2009. Incorporating SMS into the V\&V standards
body would require (i) scale-up empirics on large multi-module CFD code
(left to the P5 paper); (ii) multi-year dialogue with the ASME V\&V 20 /
IEEE 1012 / IEC 60880 committees. This paper does \textbf{not} advocate
that SMS enter any normative certification system in 2027.

\textbf{SMS does not replace}: Domain expert (SME) signatures, safety
analysis (FMECA, PIRT), system-level V\&V, operational history accident
retrospection, or ASME V\&V 20 §3 numerical solver verification. SMS is
\textbf{one link} in the test adequacy evidence chain, not a
single-point qualification determination.

\paragraph{6.5.4 Common interface across
stakeholders}\label{common-interface-across-stakeholders}

All three stakeholder classes obtain consistent numbers through the same
\texttt{data/results/paper\_numbers\_v4.json} (SSOT), and consistent
diagnostic labels through \texttt{lrca\_60cell\_v4.json}. This avoids
the document fragmentation problem of ``SMS seen by engineers
inconsistent with SMS in V\&V documentation''---this is the prerequisite
for the §6.5 stakeholder analysis to be implementable in engineering.
All numbers are to be interpreted under the single-output kernels scope
declared in §1.1 / §3.1.1.

\begin{center}\rule{0.5\linewidth}{0.5pt}\end{center}

\subsection{Section 7 $\cdot$ Risks and Mitigation +
Limitations}\label{section-7-risks-and-mitigation-limitations}

\subsubsection{7.1 Internal Threats}\label{internal-threats}

\paragraph{7.1.1 LLM Generation Reproducibility
(R1)}\label{llm-generation-reproducibility-r1}

\textbf{Risk}: Different LLM versions / temperature / seed produce
different mutant pools, affecting experimental reproducibility.

\textbf{Mitigation}: - §4.2.3 All reproducibility parameters stored in
replication package (LLM version, prompt, temperature, seed, manual
arbitration records) - §4.2.4 Dual LLM cross-source + 20\% manual
sampling - Provide ``reproducibility consistency check script'':
identical prompt + seed should produce $\geq$ 90\% overlapping mutant pool

\textbf{Residual risk}: LLM training data may be updated after
submission, rendering old API versions unavailable --- explicitly
acknowledged in §7 Limitations.

\paragraph{7.1.2 Probabilistic Approximation of equiv Determination
(R2)}\label{probabilistic-approximation-of-equiv-determination-r2}

\textbf{Risk}: K\_eq=1000 is an engineering approximation to the
undecidable problem of SMS (equivalent mutant detection), with
bidirectional bias of false-equiv and false-non-equiv.

\textbf{Mitigation}: - §2.3 Explicitly declares equiv as probabilistic
approximation rather than theorem-based determination - \st{§5 Appendix
provides sensitivity analysis for three configurations K\_eq $\in$ \{500,
1000, 2000\}} \textbf{(R1 W3 round-2 revision)}: The K\_eq sweep
sensitivity table was not executed in this submission; downgraded to
§7.5 Limitations as a residual threat. The Hoeffding-style upper bound
still provides a theoretical bound on the false-equiv probability; an
empirical K\_eq sweep is left to the P4 paper. - §7 Limitations cites
the Hoeffding-style upper bound on false-equiv probability; K\_eq
sensitivity sweep is deferred to the P4 paper.

\paragraph{7.1.3 Circular Dependency Between AVP and P1
(R3)}\label{circular-dependency-between-avp-and-p1-r3}

\textbf{Risk}: P1 is under SANER submission; if P1 modifies AVP
implementation, P2 data becomes invalid.

\textbf{Mitigation}: - §4.5.1 AVP version fixed to P1 commit hash - P2
replication package embeds complete AVP source code (P2 remains
self-consistent even if P1 undergoes major revision) - §1.6 Explicitly
declares {[}Meng Li et al., Progress in Nuclear Energy, under review{]}
citation as arXiv technical report

\paragraph{7.1.4 Boundary of LRCA Multi-label Determination
(R4)}\label{boundary-of-lrca-multi-label-determination-r4}

\textbf{Risk}: C2-C5 may trigger multiply on class B/D PUTs; priority
C5\textgreater C4\textgreater C3\textgreater C2\textgreater C1 is an
engineering choice that may mask mixed root causes.

\textbf{Mitigation}: - §2.6.3 Decision tree explicitly declares priority
- §5.4 LRCA diagnosis appends \textbf{multi-label co-occurrence table}
(all triggered layers for each killed mutant), not only reporting
priority-winning root cause - §7 Limitations explicitly acknowledges
root\_cause label as ``likely root cause,'' not definitive attribution

\textbf{§5.6.2 Empirical addendum}: LRCA's 3 engineering thresholds have
been calibrated via 9-grid search (§4.6.4); optimal combination
(ood\_band=0.02, tolerance\_multiplier=3.0, repeats=20) improved H5 pass
rate from default 10/60 to 12/60; §5.6.2 reported values adopt this
optimal combination. Calibration found tolerance\_multiplier has zero
impact on results within the 9-grid; dominant signal comes from
ood\_band. This paper uses this calibration as official results; default
threshold results retained as control (\texttt{lrca\_60cell\_v3.json}).

\paragraph{7.1.5 Operator Registry-PUT Source Code Drift
(R8)}\label{operator-registry-put-source-code-drift-r8}

\textbf{Risk}: During v2 $\to$ v2.1 revision, 6 / 37 operator definitions
were found to reference parameters that no longer exist after PUT
refactoring (e.g., GPR.alpha vs WhiteKernel.noise\_level; d1 registry
declares SVM but PUT is actually MLP). Such drift causes mutant
generation to have no executable match with PUT, polluting R\_sem
statistics.

\textbf{Mitigation}: Added pre-check consistency scan in §4.2 (key
identifiers in target\_locator must appear in PUT source code, otherwise
that operator is skipped for that PUT); v2.1 revision log recorded in
\texttt{data/operator\_campaign/v2\_revised6.log}.

\paragraph{7.1.6 Mutant Pool Size (R9)}\label{mutant-pool-size-r9}

\textbf{Risk}: 12 mutants per PUT is an engineering-cost balance:
smaller pools yield coarse SMS estimation jumps (each mutant contributes
1/12 $\approx$ 0.083 step size), larger pools exceed weekly Opus subscription
quota for LLM calls. §5.7 bootstrap CI reflects uncertainty from this
source, also partially explaining why H2 formally did not exceed the
0.474 threshold.

\textbf{Mitigation + empirical addendum}: Replication package includes
\texttt{data/operator\_campaign/cache/} (212 confirmed mutants) and
\texttt{scripts/build\_pools.py} (supports \texttt{POOL\_VERSION=v3} to
switch to 30 mutants/PUT). We empirically expanded pool to average 17.4
mutants/PUT (limited by cache capacity, some PUTs below 30): $\delta$ increased
marginally from 0.321 to 0.323, CI narrowed from {[}0.021, 0.639{]} to
{[}0.017, 0.622{]}, H2 still not met. This indicates effect size is an
\textbf{intrinsic upper bound of LLM-homogeneous mutant pools} on this
dataset, not an artifact of 12-mutant pool dilution (see §6.1). P4 paper
will retest H2 with \textbf{cross-source mutant pools} (mixing multiple
LLM backends).

\paragraph{7.1.7 Non-determinism of LLM Generation
(R10)}\label{non-determinism-of-llm-generation-r10}

\textbf{Risk}: Claude Opus subscription interface lacks seed control;
identical prompt may produce different mutant outputs at different
times.

\textbf{Mitigation} (three-part): - (a) Multi-turn de-dup enforces
structural differences among candidates (§4.2.1); - (b) K=10 / K=20
repetitions reduce single-point bias of individual operators (§4.8); -
(c) \texttt{data/operator\_campaign/raw/} commits complete prompt + raw
response, allowing replication experiments to directly reuse the same
mutant set used in this paper, bypassing non-determinism.

\paragraph{7.1.8 R11 Selection-on-response Chained Conditioning (NEW,
P0-5)}\label{r11-selection-on-response-chained-conditioning-new-p0-5}

Both v3b and v4 data are conditional on post-hoc selection of §3.5.1
c-class primary MP (c1/c2/c3 $\to$ MP1). Specific consequences:

\begin{enumerate}
\def\labelenumi{(\alph{enumi})}
\item
  \textbf{v3b sign test 4/4} and \textbf{v4 sign test 4/4} both inherit
  max-over-5 selection inflation; Bonferroni upper bound $\alpha$\_effective =
  0.01 (§3.5.1 P0-4 revision), permutation p-values in
  \texttt{data/results/c\_class\_permutation\_v4.json}.
\item
  \textbf{$\Delta$$\delta$\_\{v3b $\to$ v4\} = −0.007} is a contrast conditional on v3b
  selection + identical prompt, \textbf{not a neutral-condition test of
  LLM source diversity}. Neutral version (v4-pre $\times$ c$\to$MP5 pre-shift, and
  §4.2.5.1 differential prompt) reserved for P4 / R2 revision.
\end{enumerate}

\textbf{Mitigation}: Abstract / §5.8.2 / §6.3 (P0-3 revision) have
downgraded all v3b/v4 sign test results to exploratory; §4.2.5(e) (P0-5
this section revision) explicitly declares chained-conditioning; §3.5.1
(P0-4 revision) added permutation + Bonferroni quantification.
\textbf{This threat has been minimized within the scope of this paper
version but cannot be eliminated --- elimination requires v4-pre rerun
or P4 differential prompt experiment}.

\paragraph{7.1.9 R13 Protocol-implementation Gap Between v3/v3b and v4
(NEW,
P1-7)}\label{r13-protocol-implementation-gap-between-v3v3b-and-v4-new-p1-7}

v3 / v3b data collection used §4.2.4 original Phase-1 dual-blind
protocol (Claude-Opus generation + GPT-5.4 review + DeepSeek
arbitration); v4 cross-source pool only passed V1-V4 mechanical gates,
\textbf{lacking LLM reviewer review stage} (§4.2.5(b)).

\textbf{Potential confound}: Part of $\Delta$$\delta$\_\{v3b $\to$ v4\} = −0.007 variation
may not be ``LLM source diversity contribution,'' but rather ``slight
downward shift in mutant quality of v4 pool relative to v3b pool.''

\textbf{Mitigation (within this paper's scope)}: - v4 pool \textbf{did
use source diversity from three LLM providers} (Claude / GPT-5.4 /
DeepSeek each contributing \textasciitilde1/3 mutants, §4.2.5(d) three
providers contributed 101 / 98 / 99), in this sense $\Delta$$\delta$\_\{v3b $\to$ v4\} is
not a trivial single-source-vs-single-source comparison; - §6.2 LRCA
data shows v4 pool mean C1\_share 0.209 higher than v3b's 0.164 (quality
increased rather than decreased), weakly opposing ``v4 mutant quality
downward'' hypothesis; - But \textbf{complete separation of protocol
asymmetry and source diversity still requires P4 to rerun dual-blind
reviewer on v4 full grid} (estimated 60-100 mutant sampling level
achievable, \textasciitilde\$5-8 USD).

\textbf{Relationship to R11}: R11 concerns v3b $\to$ v4 \emph{selection}
asymmetry (c-class primary MP selection inheritance); R13 concerns v3 /
v3b $\to$ v4 \emph{protocol} asymmetry (dual-blind vs V1-V4 only). Both
asymmetries contribute to the explanation space of $\Delta$$\delta$\_\{v3b $\to$ v4\},
\textbf{cannot be merged as single ``LLM diversity contribution''
signal}.

\subsubsection{7.2 External Threats}\label{external-threats}

\paragraph{7.2.1 Representativeness of 12 PUTs
(R5)}\label{representativeness-of-12-puts-r5}

\textbf{Risk}: 12 PUTs do not represent the full domain of scientific
computing software (e.g., lacking molecular dynamics, quantum chemistry,
CFD large-scale code).

\textbf{Mitigation}: - §3.1 Follows P1 site selection ---
representativeness argument carried by P1 - §1.6 P-II open principle:
cls can be extended with new classes (E differential equation systems, F
symbolic computation, etc.)

\paragraph{7.2.2 Statistical Power of Cross-class Consistency
(R6)}\label{statistical-power-of-cross-class-consistency-r6}

\textbf{Risk}: 4-class sign test has weak power (df=3), CV unstable in
small samples.

\textbf{Mitigation}: - §5.3.2 H4 framed as \textbf{exploratory
evidence}, supplemented by mixed-effects model - §7 Explicitly
acknowledges ``cross-class consistency'' is a descriptive conclusion at
12 PUT scale

\textbf{§5.8 Empirical addendum}: Planned mixed-effects model
(\texttt{sms\ \textasciitilde{}\ C(class)\ *\ C(operator)\ +\ (1\ \textbar{}\ put)})
exhibited Singular matrix at N = 60 observations (primary model);
fallback model PUT random intercept variance degenerated to 0,
essentially degenerating to OLS. We therefore changed RQ3 main
conclusion to direct presentation via ``class means + sign test + forest
plot'' triad, rather than using mixed-effects p-value as formal
hypothesis test. This is consistent with ``small N alternative'' already
declared in §5.3.2, but should be honestly declared in conclusions that
mixed-effects is unavailable (§5.8.3).

\subsubsection{7.3 Construct Threats}\label{construct-threats}

\paragraph{7.3.1 SMS Semantic Detection Capability vs Engineering
Practical
Value}\label{sms-semantic-detection-capability-vs-engineering-practical-value}

\textbf{Risk}: SMS measures ``theoretical semantic detection
capability,'' not equivalent to ``production environment engineering
value.''

\textbf{Mitigation}: - §1.6.2 Epistemological declaration: SMS is not an
engineering value proxy metric - §1.3 P2-CN (Nuclear Power Engineering
2027 Q4) will specifically address engineering value dimension

\paragraph{7.3.2 LLM Homogeneity Bias
(R7)}\label{llm-homogeneity-bias-r7}

\textbf{Risk}: Even with dual LLM cross-source (Claude Opus + GPT-4o),
training data may still overlap, causing similar blind spots.

\textbf{Mitigation}: - §4.2.4(f) Three-LLM rotation + PUT class
subdivision - 20\% manual sampling as baseline - §7 Limitations
explicitly acknowledges residual risk, recommends future work introduce
third-party independent LLM family validation

\subsubsection{7.4 Conclusion Threats}\label{conclusion-threats}

\paragraph{7.4.1 Multiple Comparisons}\label{multiple-comparisons-1}

\textbf{Risk}: 60 cells + 5 hypotheses $\to$ multiple comparison false
positives.

\textbf{Mitigation}: §5.3.1 Benjamini-Hochberg FDR correction, $\alpha$\_FDR =
0.05.

\paragraph{7.4.2 Stability of N=20
Repetitions}\label{stability-of-n20-repetitions}

\textbf{Risk}: N=20 Cliff's $\delta$ and odds ratio 95\% CI are relatively
wide.

\textbf{Mitigation}: §5.3.2 Reports 1000-iteration bootstrap CI; §7
explicitly acknowledges effect size estimation uncertainty.

\subsubsection{7.5 Limitations Section (Final Draft Body
Text)}\label{limitations-section-final-draft-body-text}

\begin{quote}
This paper has the following known limitations: \textbf{(1) Equivalent
mutant determination is probabilistic approximation} (engineering
implementation of K\_eq=1000 input sampling, not complete theorem-based
determination); \textbf{(2) LLM generation homogeneity bias} (dual LLM
cross-source + 20\% manual sampling as mitigation, cannot be
eliminated); \textbf{(3) Statistical power of cross-class consistency}
(4-class sign test has only 3 degrees of freedom, H4 is exploratory
evidence); \textbf{(4) LRCA likely root cause is most probable
attribution} (decision tree priority is engineering choice, not causal
conclusion); \textbf{(5) AVP reuses P1 implementation} (P2 replication
package embeds AVP source code to maintain self-consistency, but
interface semantics may change as P1 evolves); \textbf{(6) SMS measures
semantic detection capability in the epistemological sense, does not
directly represent engineering practical value} (the latter is
specifically addressed by P2-CN and P5 in the nuclear engineering
domain).
\end{quote}

\begin{center}\rule{0.5\linewidth}{0.5pt}\end{center}

\subsection{Section 8 $\cdot$ References}\label{section-8-references}

\begin{quote}
Citation style: APA-7. Academic literature (papers/books) provides
complete venue + year + DOI/URL; software tools follow software citation
conventions (project homepage + version/year). All entries verified via
WebSearch / DOI as of 2026-05-01.
\end{quote}

\subsubsection{8.1 Mutation testing classics /
surveys}\label{mutation-testing-classics-surveys}

\begin{itemize}
\tightlist
\item
  \textbf{DeMillo, R. A., Lipton, R. J., \& Sayward, F. G.} (1978).
  Hints on test data selection: Help for the practicing programmer.
  \emph{Computer}, 11(4), 34--41.
  https://doi.org/10.1109/C-M.1978.218136 (Origin of the
  \textbf{Coupling Effect Hypothesis (CPH)}: complex faults are
  detectable by tests that detect simple faults; cited in §1.3.2 as
  theoretical grounding for syntactic-mutation-based adequacy and
  contrasted with this paper's domain-semantic adequacy.)
\item
  \textbf{Jia, Y., \& Harman, M.} (2011). An analysis and survey of the
  development of mutation testing. \emph{IEEE Transactions on Software
  Engineering}, 37(5), 649--678. https://doi.org/10.1109/TSE.2010.62
\item
  \textbf{Jia, Y., \& Harman, M.} (2009). Higher Order Mutation Testing.
  \emph{Information and Software Technology}, 51(10), 1379--1393.
  https://doi.org/10.1016/j.infsof.2009.04.016
\item
  \textbf{Andrews, J. H., Briand, L. C., \& Labiche, Y.} (2005). Is
  mutation an appropriate tool for testing experiments? In
  \emph{Proceedings of the 27th International Conference on Software
  Engineering} (ICSE 2005) (pp.~402--411). ACM.
  https://doi.org/10.1145/1062455.1062530 (Empirical foundation that
  mutants are valid surrogates for real faults; cited in §1.3.2 as the
  substrate for using MS as an adequacy proxy.)
\item
  \textbf{Just, R., Jalali, D., Inozemtseva, L., Ernst, M. D., Holmes,
  R., \& Fraser, G.} (2014). Are mutants a valid substitute for real
  faults in software testing? In \emph{Proceedings of the 22nd ACM
  SIGSOFT International Symposium on Foundations of Software
  Engineering} (FSE 2014) (pp.~654--665). ACM.
  https://doi.org/10.1145/2635868.2635929 (Reaffirms and refines Andrews
  2005's mutant-as-fault-proxy claim with stronger empirical evidence;
  cited in §1.3.2 as additional CPH validation.)
\item
  \textbf{Papadakis, M., Kintis, M., Zhang, J., Jia, Y., Le Traon, Y.,
  \& Harman, M.} (2019). Mutation testing advances: An analysis and
  survey. \emph{Advances in Computers}, 112, 275--378.
  https://doi.org/10.1016/bs.adcom.2018.03.015 (Most-recent
  comprehensive survey of mutation testing advances post-Jia \& Harman
  2011; cited in §1.3.2 as the contemporary literature anchor for
  syntactic mutation testing methodology.)
\item
  \textbf{Kintis, M., Papadakis, M., Papadopoulos, A., Valvis, E.,
  Malevris, N., \& Le Traon, Y.} (2018). How effective are mutation
  testing tools? An empirical analysis of Java mutation testing tools
  with manual analysis and real faults. \emph{Empirical Software
  Engineering}, 23(4), 2426--2463.
  https://doi.org/10.1007/s10664-017-9582-5
\item
  \textbf{Ammann, P., \& Offutt, J.} (2008). \emph{Introduction to
  software testing} (1st ed.). Cambridge University Press. (Standard
  pedagogical reference for mutation testing notation and equivalence
  detection; cited in §2.1.1 vocabulary inheritance.)
\end{itemize}

\subsubsection{8.2 Industrial-scale mutation testing
practice}\label{industrial-scale-mutation-testing-practice}

\begin{itemize}
\tightlist
\item
  \textbf{Petrović, G., \& Ivanković, M.} (2018). State of mutation
  testing at Google. In \emph{Proceedings of the 40th International
  Conference on Software Engineering: Software Engineering in Practice}
  (ICSE-SEIP 2018) (pp.~163--171). ACM.
  https://doi.org/10.1145/3183519.3183521
\item
  \textbf{Petrović, G., Ivanković, M., Fraser, G., \& Just, R.} (2021).
  Practical mutation testing at scale: A view from Google. \emph{IEEE
  Transactions on Software Engineering}, 48(10), 3900--3912.
  https://doi.org/10.1109/TSE.2021.3107634
\end{itemize}

\subsubsection{8.3 LLM-based mutation
generation}\label{llm-based-mutation-generation}

\begin{itemize}
\tightlist
\item
  \textbf{Tip, F., Bell, J., \& Schäfer, M.} (2024). LLMorpheus:
  Mutation testing using large language models. \emph{arXiv preprint}
  arXiv:2404.09952. https://arxiv.org/abs/2404.09952
\end{itemize}

\subsubsection{8.4 Deep-learning and general fault
benchmarks}\label{deep-learning-and-general-fault-benchmarks}

\begin{itemize}
\tightlist
\item
  \textbf{Humbatova, N., Jahangirova, G., \& Tonella, P.} (2021).
  DeepCrime: Mutation testing of deep learning systems based on real
  faults. In \emph{Proceedings of the 30th ACM SIGSOFT International
  Symposium on Software Testing and Analysis} (ISSTA 2021) (pp.~67--78).
  ACM. https://doi.org/10.1145/3460319.3464825
\item
  \textbf{Just, R., Jalali, D., \& Ernst, M. D.} (2014). Defects4J: A
  database of existing faults to enable controlled testing studies for
  Java programs. In \emph{Proceedings of the International Symposium on
  Software Testing and Analysis} (ISSTA 2014) (pp.~437--440). ACM.
  https://doi.org/10.1145/2610384.2628055
\end{itemize}

\subsubsection{8.5 Statistical
methodology}\label{statistical-methodology}

\begin{itemize}
\tightlist
\item
  \textbf{Romano, J., Kromrey, J. D., Coraggio, J., Skowronek, J., \&
  Devine, L.} (2006). Appropriate statistics for ordinal level data:
  Should we really be using t-test and Cohen's d for evaluating group
  differences on the NSSE and other surveys? Annual Meeting of the
  Florida Association of Institutional Research, Cocoa Beach, FL.
  (Source of Cliff's $\delta$ small/medium/large thresholds 0.147 / 0.330 /
  0.474; cited in §5.2 H2 definition and §5.7.2 verdict.)
\item
  \textbf{Vargha, A., \& Delaney, H. D.} (2000). A critique and
  improvement of the CL common language effect size statistics of McGraw
  and Wong. \emph{Journal of Educational and Behavioral Statistics},
  25(2), 101--132. https://doi.org/10.3102/10769986025002101
  (Establishes the canonical Â$_{1}$$_{2}$ measure equivalent to Cliff's $\delta$ + 0.5,
  used as a methodological reference for non-parametric effect-size
  reporting in §5.6.)
\end{itemize}

\subsubsection{8.6 Numerical / scientific computing
reference}\label{numerical-scientific-computing-reference}

\begin{itemize}
\tightlist
\item
  \textbf{Press, W. H., Teukolsky, S. A., Vetterling, W. T., \&
  Flannery, B. P.} (2007). \emph{Numerical Recipes: The Art of
  Scientific Computing} (3rd ed.). Cambridge University Press. (Used in
  §3.1.1 to ground the 12-PUT coverage against the 12 chapters of
  Numerical Recipes.)
\item
  \textbf{ASME V\&V 20 Committee} (2009). \emph{Standard for
  Verification and Validation in Computational Fluid Dynamics and Heat
  Transfer}. ASME V\&V 20-2009. American Society of Mechanical
  Engineers. (Cited in §1.3.2 and §6.5.3 as the normative V\&V framework
  whose §3 code verification this paper's SMS is conceptually
  complementary to; this paper does not claim normative compliance, see
  §6.5.3.)
\end{itemize}

\subsubsection{8.7 Software / mutation testing tools (cited in §3.2.6
and §7
R12)}\label{software-mutation-testing-tools-cited-in-3.2.6-and-7-r12}

\begin{itemize}
\tightlist
\item
  \textbf{Hovde, A.} (2018--present). \emph{mutmut}: A Python mutation
  testing tool. https://github.com/boxed/mutmut (Python-based
  first-order syntactic mutation tool; cited as a representative of
  ``first-order syntactic tools'' in §3.2.6 caveat.)
\item
  \textbf{Tomilin, A.} (2017--present). \emph{cosmic-ray}: Python
  mutation testing. https://github.com/sixty-north/cosmic-ray (Cited
  alongside mutmut for the §5.10 planned tool ablation.)
\item
  \textbf{Hovstadius, K.} (2014--present). \emph{mutpy}: Mutation
  testing for Python. https://github.com/mutpy/mutpy (Cited in §7 R12 as
  Python-3.10+-incompatible; not used as a comparator.)
\end{itemize}

\subsubsection{8.8 Companion P-series
papers}\label{companion-p-series-papers}

\begin{itemize}
\tightlist
\item
  \textbf{Li, M. et al.} (under review). Empirical audit of
  metamorphic-relation meta-patterns in scientific computing software
  (P1). \emph{International Conference on Software Analysis, Evolution
  and Reengineering (SANER) 2027.}
\item
  \textbf{Li, M. et al.} (under review). {[}P2-CN companion{]}.
  \emph{Progress in Nuclear Energy} (NED). (Reference D7 in §1.6.1.)
\end{itemize}

\begin{center}\rule{0.5\linewidth}{0.5pt}\end{center}

\subsection{Section 9 $\cdot$ SMS-MS Degeneration Theorem (R-8 Formalized
Proof)}\label{section-9-sms-ms-degeneration-theorem-r-8-formalized-proof}

\begin{quote}
This section formalizes the core claim in the §2.0 P-I developmental
principle: \textbf{SMS strictly regresses to the classical Jia \& Harman
(2011) syntactic Mutation Score (MS) in the degenerate limit where all
extension dimensions are closed}. This theorem guarantees that the P2
metric family is upward-compatible with the existing mutation testing
literature, and that all SMS-based empirical conclusions do not
constitute metric-level semantic fragmentation in classical syntactic
mutation scenarios.
\end{quote}

\subsubsection{9.1 Notation (following
§2.1.2)}\label{notation-following-2.1.2}

Let the program under test (PUT) be S\_i, mutation operator family mut
(syntactic or semantic), mutant set mut(S\_i), meta pattern (MP) set,
equivalence tolerance $\varepsilon$\_eq, equivalence sampling count K\_eq, and AVP
tolerance $\varepsilon$\_AVP. Three-state decomposition:
\texttt{mut(S)\ =\ killed\ $\cup$\ equiv\ $\cup$\ survive} (disjoint). SMS
formula:

\[
\text{SMS}_{i,k,j} = \frac{|\text{killed}_{i,k,j}|}{|\text{mut}_j(S_i)| - |\text{equiv}_{i,k,j}|}
\]

\subsubsection{9.2 Degenerate Limit Definition (R-8 + P1-3 revision:
rewritten from 6 axes to 3 joint
conditions)}\label{degenerate-limit-definition-r-8-p1-3-revision-rewritten-from-6-axes-to-3-joint-conditions}

The degenerate limit L consists of \textbf{3 joint conditions}, each
controlling one layer of the SMS formula (numerator / denominator mut /
denominator equiv); \textbf{L1--L6 are not 6 independent axes}, but
paired joint conditions (this revision responds to the dependency
queries in R0 W8 / R1 §4 / R2 W3 / DA-MAJOR-3).

\textbf{Joint condition L\_equiv} (controls equiv degeneration layer,
Lemma 9.1): - \textbf{L1}: $\varepsilon$\_eq $\to$ 0 (equivalence tolerance approaches
zero) - \textbf{L2}: K\_eq $\to$ $\infty$ (equivalence sampling covers the complete
input space D\_S) - Pairing rationale: On continuous D\_S, when L1 holds
alone but L2 does not, equiv remains a probabilistic approximation
(K\_eq samples cannot cover the entire D\_S); when L2 holds alone but L1
does not, the bitwise equality condition is diluted by $\varepsilon$\_eq tolerance.
Both must take the limit simultaneously for equiv to degenerate to
classical behavioral equivalence (which also holds strictly only outside
a D\_S-measure-zero set, see revised statement of Lemma 9.1).

\textbf{Joint condition L\_killed} (controls killed degeneration layer,
Lemma 9.2): - \textbf{L3}: $\varepsilon$\_AVP\^{}k $\to$ 0 for all k $\in$ MP (AVP tolerance
approaches zero) - \textbf{L4}: MP set = \{equality-checking MP\_eq\}
(R(y, y') ≡ y = y') - Pairing rationale: When L3 holds but L4 does not,
$\varepsilon$\_AVP $\to$ 0 still allows non-trivial MP relations to exist (R can be
monotonicity, convergence order, etc.), not degenerating to classical
difference detection; when L4 holds but L3 does not, equality checking
still carries $\varepsilon$\_AVP tolerance, not strictly enforced. Both must take
the limit simultaneously for killed determination to degenerate to
classical difference detection.

\textbf{Joint condition L\_mut} (controls mut degeneration layer, Lemma
9.3): - \textbf{L5}: mut\_j switches to rule-based syntactic operators
(Mothra-style AOR/ROR/SDL/CRP, etc.), independent of domain semantics -
\textbf{L6}: PUT class cls(I) $\subseteq$ \{imperative deterministic programs\}
(no probabilistic/surrogate/ML) - Pairing rationale: When L5 holds but
L6 does not, syntactic operators on probabilistic/ML programs may still
trigger subsets of domain-semantic mutation operators (e.g., literal
constant replacement of dropout probability); when L6 holds but L5 does
not, imperative deterministic programs can still be mutated by semantic
operators (OS/HP/TF/SI), mut(S) $\neq$ syntactic mutants. Both must take the
limit simultaneously for mut(S) to degenerate to the syntactic mutant
set in the Jia \& Harman literature.

\textbf{Total limit L = L\_equiv $\wedge$ L\_killed $\wedge$ L\_mut} (all three joint
conditions hold simultaneously).

\subsubsection{9.3 Lemmas: Three-State Decomposition Degenerates Under
L}\label{lemmas-three-state-decomposition-degenerates-under-l}

\textbf{Lemma 9.1} (equiv degeneration, P1-3 revision: added
measure-zero qualification). Under joint condition L\_equiv (L1 $\wedge$ L2),
semantic-class equivalence (E1 $\wedge$ E2) degenerates to classical behavioral
equivalence \textbf{almost everywhere} (almost everywhere w.r.t. measure
D\_S).

\textbf{Proof}: - E1 (type consistency) holds trivially in the $\varepsilon$\_eq $\to$ 0
limit (L6 imperative program output spaces are scalar/vector, types
statically guaranteed by the programming language). - E2
(numerical/semantic approximate equality) is defined as: for K\_eq
sampled inputs x \textasciitilde{} D\_S, \textbar S\_i(x) −
s'(x)\textbar{} \textless{} $\varepsilon$\_eq. Under L1 ($\varepsilon$\_eq $\to$ 0) $\wedge$ L2 (K\_eq $\to$ $\infty$
with measure-equivalent sampling to D\_S), this condition is
\textbf{almost everywhere equivalent to} $\forall$x $\in$ D\_S ~N, S\_i(x) = s'(x),
where N is a D\_S-measure-zero set (on continuous D\_S, strict bitwise
equality requires excluding measure-zero exceptions, such as numerical
NaN propagation points or floating-point cancellation pathological
points; on discrete D\_S, N = $\emptyset$, strict bitwise equality)---this is
consistent with the classical equivalent mutant definition in Jia \&
Harman (2011) §3 under measure-zero equivalence classes. $\blacksquare$

\textbf{Lemma 9.2} (killed degeneration). Under L3 $\wedge$ L4, killed
determination degenerates to classical difference detection.

\textbf{Proof}: - L4 restricts the MP set to \{MP\_eq\}, where MP\_eq's
relation R(y, y') ≡ y = y'. - Given mr = (r, R) $\in$ MR, where r is an
input transformation and R is an output relation. For mutant s', the
violation condition for MP\_eq is $\exists$x: S\_i(x) $\neq$ s'(r(x)). - Under L3
($\varepsilon$\_AVP $\to$ 0), AVP tolerates no numerical difference; violation is
equivalent to the exact inequality S\_i(x) $\neq$ s'(r(x)). - When r = id
(identity transformation), the violation condition becomes S\_i(x) $\neq$
s'(x), i.e., the mutant deviates from the original program on some
input---this is precisely the classical difference detection semantics.
- When r $\neq$ id, MP\_eq restricted by L4 still requires S\_i(x) =
s'(r(x)), treating it as a ``reference output oracle'' constructed from
the original program; this still does not introduce new state
classifications. $\blacksquare$

\textbf{Lemma 9.3} (mut degeneration). Under L5 $\wedge$ L6, mut\_j(S\_i)
degenerates to the syntactic mutant set in the Jia \& Harman (2011)
literature.

\textbf{Proof}: L5 explicitly switches mut\_j to rule-based syntactic
operators (AOR, ROR, SDL, CRP, UOI, etc., standard Mothra/Proteum sets);
L6 restricts PUTs to imperative deterministic programs, excluding
triggering conditions for semantic operators on probabilistic/ML
programs. Under this configuration, mut\_j(S\_i) is the syntactic mutant
set as defined in the literature, independent of domain semantics. $\blacksquare$

\subsubsection{9.4 Main Theorem: SMS $\to$ MS}\label{main-theorem-sms-ms}

\textbf{Theorem 9.1} (SMS-MS degeneration theorem, P1-3 revision). In
the degenerate limit L = L\_equiv $\wedge$ L\_killed $\wedge$ L\_mut, \textbf{almost
everywhere} (almost everywhere w.r.t. D\_S),

\[
\text{SMS}_{i,k,j} \xrightarrow{L} \text{MS}_{i,j} := \frac{|\text{killed}_{i,j}^{\text{classic}}|}{|\text{mut}_j^{\text{syntax}}(S_i)| - |\text{equiv}_{i,j}^{\text{classic}}|}
\]

where the right-hand side is the classical Mutation Score of Jia \&
Harman (2011), killed\^{}classic is the difference detection set,
equiv\^{}classic is the behavioral equivalent mutant set, and
mut\^{}syntax is the syntactic mutant set.

\textbf{Proof}: By Lemmas 9.1-9.3,

\begin{itemize}
\tightlist
\item
  \textbf{Numerator}: Under L3 $\wedge$ L4, killed\_\{i,k,j\} $\to$
  killed\_\{i,j\}\^{}\{classic\} (Lemma 9.2). L4 simultaneously makes
  the MR\_\{i,k\} set trivial over k (only MP\_eq remains), the index k
  degenerates, and the subscript k in SMS can be omitted.
\item
  \textbf{Denominator \textbar mut\_j(S\_i)\textbar{}}: Under L5 $\wedge$ L6 $\to$
  \textbar mut\_j\^{}\{syntax\}(S\_i)\textbar{} (Lemma 9.3).
\item
  \textbf{Denominator \textbar equiv\_\{i,k,j\}\textbar{}}: Under L1 $\wedge$
  L2 $\to$ \textbar equiv\_\{i,j\}\^{}\{classic\}\textbar{} (Lemma 9.1).
\end{itemize}

Substituting into the numerator and denominator on the right-hand side
of the SMS formula yields MS\_\{i,j\}. $\blacksquare$

\subsubsection{9.5 Corollary: LRCA
Trivialization}\label{corollary-lrca-trivialization}

\textbf{Corollary 9.1} (generic statement, post R2 round-2 attribution
audit). In the degenerate limit L = L1 $\wedge$ L2 $\wedge$ L3 (the three joint
conditions formalized in §9.2), the likely root cause inventory (LRCA) C
= \{C1, \ldots, C5\} degenerates to a single state \{C1\}.

\textbf{Sketch}: Each of C2--C5's triggering preconditions depends on at
least one of the L\_j conditions being violated. Concretely, when L1 $\wedge$
L2 $\wedge$ L3 hold simultaneously, every ``non-trivial space'' dimension of
the SMS formula (MP non-triviality, AVP tolerance non-zero, non-empty
equiv set, MR-design degree of freedom, class-mapping openness) is
closed at once, so the triggering set for C2--C5 is empty (read off from
the LRCA decision tree in §2.6.1 and §4.6.3). The detailed per-C\_k to
per-L\_j minimum-sufficient mapping depends on engineering details (§4.6
LRCA classifier thresholds), and we do not claim a one-to-one
correspondence at the §9 formal level --- readers wanting the concrete
mapping can trace it via the §2.6 decision tree + §4.6 LRCA three-layer
operator documentation.

Thus under L, suspect\_share $\to$ 0, LRCA reports only C1 (metric direct
readout) --- SMS degenerates to a single-layer metric, consistent with
the engineering attribution structure of Jia \& Harman (2011) MS. $\blacksquare$

\subsubsection{9.6 Empirical Consistency
Statement}\label{empirical-consistency-statement}

Theorem 9.1 + Corollary 9.1 jointly guarantee: \textbf{Any SMS-based
empirical conclusions (such as Cliff's $\delta$, Friedman $\chi$$^{2}$, Spearman $\rho$ in
§5.7-§5.9) are structurally consistent with the existing Jia \& Harman
(2011) literature in classical syntactic mutation scenarios}, and do not
constitute metric-level semantic fragmentation. This is the intrinsic
scientific guarantee of this paper's symbolic system (§2.1) under the
``symbolic stability statement'' in §2.1.3.

\begin{center}\rule{0.5\linewidth}{0.5pt}\end{center}

\subsection{Locked Decision Inventory (D1-D8 +
W-2)}\label{locked-decision-inventory-d1-d8-w-2}

\begin{enumerate}
\def\labelenumi{\arabic{enumi}.}
\tightlist
\item
  \textbf{D1} Paper identity = $\gamma$ dual-track parallel
\item
  \textbf{D2} Mutation validity = objective artifact triple (AVP +
  general defect case library + mutator literature prior, no FMECA/PIRT)
\item
  \textbf{D3} Paradigm coverage = S-A all 4 classes $\times$ all 5 operators =
  60-cell matrix
\item
  \textbf{D4} equiv concept = AVP-coherent + K\_eq=1000 tolerance
  equivalence
\item
  \textbf{D5} Defect case library = B-1 + open-source issue systematic
  mining
\item
  \textbf{D6} SMS = M-1 classical \texttt{killed/(mut−equiv)} structure,
  decoupled from MR construction process
\item
  \textbf{D7} P1/P4 interface = I-2 mid-interface, shared {[}Meng Li et
  al., Progress in Nuclear Energy, under review{]} arXiv
\item
  \textbf{D8} P2/P2-CN/P5 triangular division of labor = D-3 Pre-NED
  trial
\item
  \textbf{W-2} Mid-tier scale = 60 cells $\times$ 10-15 mutants $\times$ N=20,
  \textasciitilde\$1200 compute
\end{enumerate}

\subsection{Dual Fundamental Principles (§2
chapter-level)}\label{dual-fundamental-principles-2-chapter-level}

\begin{itemize}
\tightlist
\item
  \textbf{P-I Developmental}: classical mutation testing $\to$ scientific
  computing semanticization (adapted to MR validity determination)
\item
  \textbf{P-II Stable/Open}: symbolic skeleton fixed, content
  (operators/MP/root causes/class mappings) open to extension
\end{itemize}

\begin{center}\rule{0.5\linewidth}{0.5pt}\end{center}

\emph{This file is the P2 paper draft §1-§7 locked snapshot (refined
version).} \emph{Generation date: 2026-04-29}

\end{document}
